% =============================================================
% Appendix E: TKF-DP --- Dirichlet-process Potts coupling
% Source: ~/tkf-dp/main.tex Sections 2 and 3, quoted near-verbatim.
% Section numbering and labels have been adapted to fit the
% supplement's overall scheme. References to the present
% appendix's own equations are local; references to other
% appendices use their supplement-level labels.
% =============================================================

The TKF-DP model decorates TKF92 with three independent Dirichlet
processes that govern (i) the per-site key partition into Potts
coevolutionary cliques; (ii) the per-site assignment to a
substitution class; and (iii) the per-class-pair assignment to a
Potts coupling atom. The class-level substitution dynamics is a
multi-site reversible CTMC; the alignment likelihood remains plain
TKF92. Inference proceeds by a variational E-step on the
continuous-time class clique, Gibbs sweeps and Jain--Neal
split--merge proposals on the three Dirichlet partitions, recursive
traceback over a time-indexed gravestone-augmented pair SCFG for
the augmented branch history, and a stochastic M-step on the
corpus-level hyperparameters with closed-form
Dirichlet--multinomial / Gamma--Poisson / Laplace conjugacies
unlocking the new-class / new-atom base-measure integrals.

\subsection{The TKF-DP generative model}
\label{sec:tkfdp-model}

The TKF92 model~\citep{ThorneEtAl92} of insertion, deletion, and
substitution has the convenient property that the alignment path
and the substitution history factorize: the marginal distribution
over alignments is independent of the amino acid dynamics, and
inference proceeds via a pair HMM on the alignment with substitution
likelihoods plugged in pointwise. This factorization is exactly
what is lost when site--site coevolutionary couplings are introduced
naively, since the substitution process at any site then depends on
the joint state of its coupled partners, and any coupling rule that
consults the existing alignment to decide partners destroys the
marginal pair HMM.

We describe a minimal extension of TKF92 that re-introduces
site--site Potts couplings while preserving the
alignment--substitution factorization. Each newborn site is assigned
a discrete key drawn i.i.d.\ from a Dirichlet process, and sites
that share a key co-evolve under a Potts-modulated joint CTMC.
Because the keys are drawn independently of the alignment, the
alignment marginal is unchanged.

The alignment--substitution factorization survives at the parameter
level but, once couplings are present, breaks at the latent-variable
level: the substitution likelihood conditional on the alignment
depends on the full timed indel history, including transient
(alignment-invisible) lineages whose presence transiently modifies
the rates of co-class members. Inference therefore requires a
tractable representation of timed indel histories. We obtain one
via a gravestone-augmented latent state and a time-indexed pair
SCFG whose inside--outside is closed form
(\S\ref{sec:gravestone-scfg}).

\begin{definition}[TKF-DP]
\label{defn:tkfdp}
Fix TKF92 indel parameters $(\lambda, \mu, r)$, a key concentration
parameter $\alpha_z > 0$, and a symmetric $20 \times 20$ Potts
coupling matrix $H$. The process jointly generates a TKF92
alignment path, a key assignment for every site, and an amino acid
trajectory at every site, as follows.
\begin{enumerate}
    \item Sample stick-breaking weights once, globally:
    \[
        \pi_k = \beta_k \prod_{l < k} (1 - \beta_l),
        \qquad
        \beta_k \stackrel{\text{iid}}{\sim} \mathrm{Beta}(1, \alpha_z).
    \]
    \item Run TKF92 indel dynamics with parameters $(\lambda, \mu, r)$,
      giving a stochastic set of alive sites $S(t)$ over time.
    \item At the birth of each new site $s$, draw a key
    \[
        z_s \mid \pi \sim \mathrm{Cat}(\pi),
    \]
    independently of all other sites, of the alignment path, and of
    the substitution history. The key is fixed for the site's
    lifetime.
    \item At any time $t$, the alive sites $S(t)$ are partitioned
      into equivalence classes $\{C_k(t)\}$ by the relation $z_s =
      z_{s'}$. The amino acid states $\{x_s\}_{s \in S(t)}$ evolve
      as independent CTMCs across classes; the CTMC on class $C$ is
      a $|C|$-site reversible chain with stationary distribution
    \[
        \pi_C(x_C) \propto \exp\!\Big(-\!\!\sum_{\{i,j\} \subset C} H(x_i, x_j)\Big).
    \]
    Class composition is piecewise constant between birth/death
    events; on each event, the affected class is rebuilt and the
    surviving sites re-equilibrate under the new class generator.
\end{enumerate}
\end{definition}

\begin{remark}[Factorization]
Conditional on $\pi$, the key draws are i.i.d.\ across sites, so the
joint distribution factorizes as
\[
    P(\mathrm{alignment},\, \pi,\, \{z_s\},\, \mathrm{substitution})
    = P(\mathrm{alignment}) \cdot P(\pi) \cdot \prod_s P(z_s \mid \pi)
      \cdot P(\mathrm{substitution} \mid \mathrm{alignment}, \{z_s\}).
\]
In particular, the marginal alignment likelihood is exactly that of
TKF92, and the pair HMM one would use under $\alpha_z = \infty$ (no
coupling) is unchanged.
\end{remark}

\begin{remark}[Partition statistics]
Marginalizing $\pi$, the partition induced on any fixed finite set
of sites is the Ewens partition with parameter $\alpha_z$,
equivalent to a Chinese restaurant process draw~\citep{aldous1985}.
The expected number of co-evolving classes among $L$ sites is
$\mathcal{O}(\alpha_z \log L)$, and tuning $\alpha_z$ controls the
sparsity of couplings: $\alpha_z \gg L$ yields
$\mathcal{O}(L^2 / \alpha_z)$ co-evolving pairs with higher-order
classes suppressed by additional factors of $L / \alpha_z$. Smaller
$\alpha_z$ gives Ewens-distributed class sizes with heavier tails,
including occasional larger cliques.
\end{remark}

\subsection{IBP variant}
\label{sec:tkfdp-ibp}

To allow each site to participate in multiple independent couplings,
replace the Dirichlet process with an Indian Buffet
Process~\citep{ibp}.

\begin{definition}[TKF92-IBP]
Sample feature frequencies $\{\pi_k\}_{k \geq 1}$ from an IBP with
concentration $\alpha_z$. At the birth of each site $s$, draw a
binary feature vector $b_s \in \{0, 1\}^{\mathbb{N}}$ with
\[
    b_{s,k} \mid \pi_k \stackrel{\text{ind}}{\sim} \mathrm{Bern}(\pi_k),
\]
independently across sites. Two alive sites $s, s'$ are coupled iff
there exists $k$ with $b_{s,k} = b_{s',k} = 1$. The amino acid
dynamics is the reversible CTMC on the alive sites with stationary
distribution
\[
    \pi(x_S) \propto \exp\!\Big(-\!\!\sum_{\{s, s'\} \,:\, \exists k,\, b_{s,k} = b_{s',k} = 1} H(x_s, x_{s'})\Big).
\]
\end{definition}

The same factorization observation applies: feature draws are i.i.d.\
given $\{\pi_k\}$, so the marginal alignment likelihood remains
plain TKF92. The IBP variant naturally accommodates hub structure
(residues that interact with many partners) via heavy-tailed feature
frequencies.

\subsection{Site classes and a GTR-parameterized generator}
\label{sec:gtr-potts}

Real proteins exhibit substantial heterogeneity in equilibrium amino
acid usage and in coupling profile: positions in $\alpha$-helices,
$\beta$-sheets, buried cores, and disordered regions have distinct
frequencies and distinct interaction patterns with their structural
neighbours. We extend the model with a per-site latent class
variable that determines both the uncoupled (point) substitution
model and the per-pair Potts interaction tensor. The variable is
non-evolving, drawn once at site birth, and we put a second
Dirichlet process on it --- independent of the key DP --- so that
the number of site classes is itself learned from data, in line
with the original infinite-mixture CAT model of Lartillot and
Philippe~\citep{lartillotcat}.

\begin{definition}[Site-class TKF-DP]
Fix the alphabet $\mathcal{A}$ with $|\mathcal{A}| = A = 20$, a
symmetric exchangeability matrix
$S \in \mathbb{R}^{A \times A}_{\geq 0}$ ($S_{xy} = S_{yx}$),
per-site rate-multiplier hyperparameters $(a_\eta, b_\eta)$, a
site-class concentration $\alpha_c > 0$, a Potts-atom concentration
$\alpha_H > 0$, and base measures
\[
    G_0: \quad \pi^{(c)} \stackrel{\text{iid}}{\sim}
              \mathrm{Dirichlet}(\kappa_\pi\, \bar\pi),
    \qquad
    G_0^H: \quad H_h(i, j) \stackrel{\text{ind}}{\sim}
              \mathcal{N}\!\bigl(\mu_{kl},\, \tau_{kl}^{-1}\bigr)
              \;\text{for } i \leq j,\;
              H_h(j, i) := H_h(i, j),
\]
where $(k, l) = (\min(i, j), \max(i, j))$ index the $A(A+1)/2$
unordered amino-acid pairs, and each Potts atom
$H_h \in \mathbb{R}^{A \times A}$ is symmetric in $(i, j)$ by
construction. Site classes are drawn from a Dirichlet process
$D_c \sim \mathrm{DP}(\alpha_c, G_0)$ via stick-breaking; Potts
atoms are drawn from a separate Dirichlet process
$D_H \sim \mathrm{DP}(\alpha_H, G_0^H)$. At the birth of each site
$s$, draw an independent per-site rate multiplier and a site class:
\[
    \eta_s \stackrel{\text{iid}}{\sim} \mathrm{Gamma}(a_\eta, b_\eta),
    \qquad
    c_s \mid D_c \stackrel{\text{iid}}{\sim} D_c,
\]
in addition to the DP key $z_s$
(\S\ref{sec:tkfdp-model}); $\eta_s$, $c_s$, and $z_s$ are mutually
independent and independent of the alignment path, and are fixed
for the site's lifetime. The Potts coupling tensor is decomposed by
class-pair: for each unordered pair of classes $\{c, c'\}$ that ever
co-occur within a key DP class, draw a Potts-atom assignment
\[
    h_{cc'} \mid D_H \stackrel{\text{iid}}{\sim} D_H,
\]
at first co-occurrence and fixed thereafter; the materialized slice
for the pair is $H_{c c'}(i, j) := H_{h_{cc'}}(i, j)$. The
amino-acid trajectory on a DP key class $C$ is the single-site
reversible CTMC with stationary distribution
\[
    \pi_C(x_C) \;\propto\; \prod_{s \in C} \pi^{(c_s)}(x_s) \cdot
    \exp\!\Big(-\!\!\sum_{\{s, s'\} \subset C}\!\! H_{c_s c_{s'}}(x_s, x_{s'}) \Big),
\]
and (taking the F81 instance reversible with respect to this
stationary)
\[
    Q^s\!\bigl(x \to x'; x_{C \setminus s}\bigr)
    \;=\; \eta_s\, S_{xx'}\, \pi^{(c_s)}(x')\,
    \exp\!\Bigl(-\tfrac{1}{2} \Delta H_s\Bigr),
    \qquad
    \Delta H_s = \!\!\sum_{s' \in C \setminus s}\!\![H_{c_s c_{s'}}(x', x_{s'}) - H_{c_s c_{s'}}(x, x_{s'})].
\]
\end{definition}

\begin{remark}[GTR limit and singletons]
For a singleton key DP class, $\Delta H_s \equiv 0$ and the rate
reduces to $Q^s_{xy} = \eta_s S_{xy} \pi^{(c_s)}(y)$, the F81
instance of GTR with class-specific equilibrium and a per-site rate
multiplier. Marginalized over $c_s$ this is exactly the CAT
infinite-mixture profile model of Lartillot and
Philippe~\citep{lartillotcat}, augmented with Yang's gamma per-site
rate variation~\citep{Yang1994a}. Coupling enters only via
$\Delta H_s$ when the key DP class size exceeds one, so the model
strictly extends the CAT--gamma combination by adding a
coevolutionary tail.
\end{remark}

\begin{remark}[F81 vs.\ symmetric-Metropolis form]
The F81 form $Q_{xy} \propto S_{xy} \pi(y)$ chosen above is one of
two natural reversible instances of the GTR family with stationary
$\pi$ and exchangeability $S$; the other is the symmetric-Metropolis
instance $Q_{xy} \propto S_{xy} \sqrt{\pi(y)/\pi(x)}$. Both satisfy
detailed balance with respect to the joint Potts stationary above.
The symmetric-Metropolis form has the property that the
eigendecomposition of its symmetrized similarity transform is
independent of $\pi$, allowing one corpus-wide $A^3$ decomposition
to be reused across all classes. We adopt the F81 form here because,
combined with the secret-destination augmentation of
\S\ref{sec:base-measure-conj}, it yields strict
Dirichlet--multinomial conjugacy on $\pi^{(c)}$, which the
symmetric-Metropolis form does not. The eigendecomposition then
becomes per-class (one $A^3$ per class per outer step), which is
small at $A = 20$ and is dominated by the Potts cluster work in any
case.
\end{remark}

\begin{remark}[Three independent Dirichlet processes]
The construction has \emph{three} Dirichlet processes acting on the
substitution side: the key DP with concentration $\alpha_z$
partitions sites into co-evolving cliques (\S\ref{sec:tkfdp-model});
the site-class DP with concentration $\alpha_c$ assigns each site a
profile $\pi^{(c)}$; and a Potts-atom DP with concentration
$\alpha_H$ and base measure $G_0^H$ assigns each unordered class-pair
$\{c, c'\}$ to a coupling atom $H_{h_{cc'}}$. With $\alpha_H$ small
the $K(K+1)/2$ class-pair indices collapse onto a small number of
canonical Potts atoms shared across class-pairs; with
$\alpha_H \to \infty$ each class-pair has its own atom. The
Potts-atom DP is what makes the model identifiable for $K_c \gg 1$:
without it, the number of free Potts parameters scales as
$K_c^2 \cdot A^2$; with it the count scales as
$K_H \cdot A^2 + K_c^2$ class-pair indices, where $K_H$ is
data-driven and typically much smaller than $K_c^2$.
\end{remark}

\begin{remark}[Per-class-pair side potentials]
\label{rem:side-potentials}
A useful refinement of the per-class profile $\pi^{(c)}$ is to give
each unordered class-pair $\{c, c'\}$ its own pair of \emph{side
potentials}
$h^{(c, c')}_a, h^{(c, c')}_b \in \mathbb{R}^A$, with independent
Gaussian priors centered at zero,
\begin{equation}
h^{(c, c')}_a(x), h^{(c, c')}_b(y) \;\sim\; \mathcal{N}(0, \tau_h^{-1}).
\end{equation}
The per-pair joint stationary then becomes
\begin{equation}
\log \pi^{(c, c')}_{\mathrm{joint}}(x, y) \;=\; \log \pi^{(c)}(x) + \log \pi^{(c')}(y) - h^{(c, c')}_a(x) - h^{(c, c')}_b(y) - H_{cc'}(x, y) - \log Z_{cc'},
\end{equation}
i.e.\ $h$ is the per-pair deviation of the effective per-site
stationary
$\pi^{\mathrm{eff}, a}_{cc'}(x) \propto \pi^{(c)}(x) e^{-h^{(c, c')}_a(x)}$
from the per-class background $\pi^{(c)}$. For self-pairs $(c, c)$
the two sites are exchangeable, so the joint distribution must be
symmetric under swap $(x, y) \leftrightarrow (y, x)$. Combined with
$H_{cc}$ symmetric in its arguments, this forces
$h^{(c,c)}_a = h^{(c,c)}_b$; we parameterise self-pairs with a
single shared vector $h^{(c,c)} \in \mathbb{R}^A$ (and prior
$h^{(c,c)} \sim \mathcal{N}(0, \tau_h^{-1})$, contributing $A$ free
parameters per self-pair rather than $2A$). Off-diagonal pairs
$\{c, c'\}$ with $c \ne c'$ retain two independent vectors
$h^{(c,c')}_a, h^{(c,c')}_b$ ($2A$ free parameters each).
Reversibility of the F81-form joint generator is preserved by
replacing the destination factor $\pi^{(c)}(x')$ in the site-1 flip
rate with $\pi^{\mathrm{eff}, a}_{cc'}(x')$ (and analogously for
site 2).

The motivation is that $\pi^{(c)}$ captures the population-mean
amino-acid usage of class $c$ \emph{averaged over all class-pair
contexts}, while the actual usage of class $c$ when paired with
class $c'$ may deviate (e.g.\ a Cys-bearing class paired with itself
is more Cys-rich than the per-class average, because we are
conditioning on disulfide-forming sites). The $h$ vectors absorb
that deviation, leaving the Potts atom $H_{cc'}$ to capture the
coupling \emph{above the marginal compositional shift}. The Gaussian
prior pulls $h$ to zero so that any pair-specific deviation must be
data-supported.

The free-$h$ Gaussian-prior formulation of $h$ presented here is
the original ``A3'' parameterisation. \autoref{sec:rev-suppl} shows
that A3 is non-reversible across partial-presence pair histories;
\autoref{sec:sinkhorn-suppl} replaces it with the unique reversible
(``A1'') alternative in which $h$ is deterministically
Sinkhorn-determined by $(\eqm^{(c)},\eqm^{(c')},H_{cc'})$ rather than
sampled. Under A1 the lone single-site stationary $\eqm^{(c)}$ used
on indel branches agrees with the joint pair-marginal exactly,
restoring detailed balance at every indel seam. The remainder of
this appendix adopts A1.
\end{remark}

\begin{remark}[Hierarchical structure of the $H$ base measure]
The amino-acid-pair-indexed Gaussian base measure
$\mathcal{N}(\mu_{kl}, \tau_{kl}^{-1})$ embeds an empirical-Bayes
hierarchy: the corpus-level $(\mu_{kl}, \tau_{kl})$ encode the
population-mean coupling propensity for each amino-acid pair
(typically attractive for hydrophobic--hydrophobic,
charge-complementary, etc., and repulsive for charge-like-like), and
atom-specific deviations $H_h(i, j) - \mu_{kl}$ are learned around
that mean with atom-specific magnitude controlled by the data.
Symmetry of each $H_h$ in $(i, j)$ encodes the unordered-edge
convention of Potts couplings; this is standard practice in the DCA
literature~\citep{ekeberg2013} and is biologically natural in the
absence of any a priori labelling distinguishing the two ends of an
inter-class edge. The Gaussian prior on each $H_h$ is conjugate to
the Gaussian (Hessian-based) Laplace approximation of the path-DCA
likelihood, giving a closed-form Gaussian posterior on each
materialized atom under the local quadratic expansion of
\S\ref{sec:base-measure-conj}.
\end{remark}

\subsection{Reversibility under partial presence}
\label{sec:rev-suppl}

The generative model of \S\ref{sec:tkfdp-model}, with the side
potentials of \autoref{rem:side-potentials} treated as
free parameters with Gaussian priors, is the construction we
will refer to as the \emph{A3 free-$h$} formulation. Under A3 the
per-pair joint stationary
\begin{equation}
  \pj^{(c,c')}(a,b)\;\propto\;\eqm^{(c)}(a)\,\eqm^{(c')}(b)\,
    e^{-h^{(c,c')}_a(a)-h^{(c,c')}_b(b)-H_{cc'}(a,b)}
  \label{eq:sidepot-joint}
\end{equation}
need not have row-marginal $\eqm^{(c)}$ or column-marginal
$\eqm^{(c')}$; the freely-chosen $h$ shifts both. The lone single-site
stationary used on indel branches where one column is gapped is
$\eqm^{(c)}$, but the pair-marginalised stationary used on
both-present branches is the marginal of \eqref{eq:sidepot-joint}.
These two stationaries disagree at any nonzero $H$, with the
disagreement carried by $h$. We show below that this disagreement
costs the model reversibility on \emph{partial-presence} pair
histories (column pairs that are present at some leaves and gapped
at others), and we develop the unique fix.

For a pair of classes $(c,c')$ the bare joint
\begin{equation}
  \pj^{(c,c')}(a,b)\;\propto\;\eqm^{(c)}(a)\,\eqm^{(c')}(b)\,e^{-H_{cc'}(a,b)},
  \qquad H_{cc'} = H_{cc'}^\top \in \R^{k\times k},
  \label{eq:base-joint-suppl}
\end{equation}
with $k=|\mathcal A|$, is the $h\equiv 0$ specialisation of
\eqref{eq:sidepot-joint}. On a tree branch where both columns are
present the pair evolves under a reversible CTMC $\Q_{\mathrm{pair}}$
with stationary \eqref{eq:base-joint-suppl}; on a branch where one
column is absent (inserted or deleted) the surviving column evolves
under its lone CTMC with stationary $\eqm^{(c)}$.

\paragraph{The defect.} The per-pair Felsenstein likelihood is
\emph{root-dependent} whenever the pair has partial presence and
$H_{cc'}\neq 0$. On a diagnostic fixture (binary alphabet, tree
$((A,B)P,C)R$, one column present at $\{A,P,R\}$ and gapped at
$B,C$), the pair log-likelihood evaluated at every admissible root
is $-3.0452$ at the three roots present in both columns ($A,P,R$)
but $-2.4401$ and $-2.5831$ at the two gapped-leaf roots ($B,C$):
the same physical (alignment, partition) state receives two
different log-densities ($\sim 1.25$ nats apart) on different
visits, silently breaking detailed balance of any
column-changing MCMC move.

\subsubsection{Reversibility is marginal-consistency}
\label{sec:revcond-suppl}

Model the augmented state as $(C, x_C)$: the alive set $C$ and
residues $x_C\in\mathcal A^{C}$. The generator has three parts:
(S)~substitution within a fixed alive set, a CTMC reversible
w.r.t.\ $\eqm_C$; (D)~deletion of a site at a residue-independent
TKF rate $\mu$ (drop the coordinate); (I)~insertion of a site at
rate $\lambda$, drawing the new residue from an insertion kernel
$g_C(y\mid x_C)$.

\begin{theorem}[Reversibility criterion]
\label{thm:revcond}
Suppose the within-level CTMC is reversible w.r.t.\ $\eqm_C$ and the
candidate stationary factors as $P(C,x_C)=w_{|C|}\,\eqm_C(x_C)$.
Then the augmented process is reversible---equivalently the
phylogenetic likelihood that peels with per-level equilibrium
$\eqm_C$ is root-independent---if and only if the family
$\{\eqm_C\}$ is \emph{marginal-consistent}:
\begin{equation}
  \sum_{x_s}\eqm_{C}(x_C) \;=\; \eqm_{C\setminus s}(x_{C\setminus s})
  \qquad\text{for every }C\text{ and every }s\in C,
  \label{eq:marg-consistency}
\end{equation}
in which case the insertion kernel is forced to the true conditional
$g_C(y\mid x_C)=\eqm_{C\cup s}(y\mid x_C)$ (the
\emph{conditional-Boltzmann} insertion) and deletion is exact
marginalisation; the level weights satisfy
$w_k=(\lambda/\mu)^k w_0$.
\end{theorem}

\begin{proof}
Within a fixed $C$, detailed balance (DB) holds by the assumed
substitution reversibility ($w_{|C|}$ is a common factor). The
binding constraints are the cross-level (birth/death) edges. For
$s\notin C$ and every $x_C,y$,
\begin{equation}
  w_{|C|}\,\eqm_C(x_C)\,\lambda\,g_C(y\mid x_C)
  \;=\;
  w_{|C\cup s|}\,\eqm_{C\cup s}(x_C,y)\,\mu .
  \label{eq:seam-db}
\end{equation}
Summing over $y$ (with $\sum_y g_C=1$) gives
$w_{|C|}\eqm_C(x_C)\lambda
 = w_{|C\cup s|}\mu\,(\eqm_{C\cup s}{\downarrow}C)(x_C)$.
Both $\eqm_C$ and the marginal $\eqm_{C\cup s}{\downarrow}C$ are
normalised, so the proportionality constant is $1$, forcing
$w_{|C\cup s|}\mu = w_{|C|}\lambda$ and
\eqref{eq:marg-consistency}. Substituting back into
\eqref{eq:seam-db} yields $g_C(y\mid x_C) = \eqm_{C\cup s}(y\mid x_C)$,
the true conditional. Conversely, given a marginal-consistent
family these equations are satisfied for all edges, so the full
generator is reversible. Reversibility of a CTMC is equivalent to
root-independence of its Felsenstein likelihood (pulley principle).
\end{proof}

\begin{remark}[The base $h\equiv 0$ model is not reversible]
The bare joint \eqref{eq:base-joint-suppl} violates
\eqref{eq:marg-consistency}: its marginal
$\sum_b \pj^{(c,c')}(a,b)\propto
 \eqm^{(c)}(a)\sum_b\eqm^{(c')}(b)e^{-H_{cc'}(a,b)}$
equals $\eqm^{(c)}(a)$ only when the bracket is constant in $a$,
i.e.\ when $H_{cc'}\equiv 0$. For any nonzero coupling the pair
marginal differs from the lone stationary $\eqm^{(c)}$ used on
indel branches, and the process is non-reversible / root-dependent
on composition-changing partial-presence pairs.
\end{remark}

Equivalently (the time-reversal view): a death is residue-independent,
so for a birth to be its exact time-reverse the surviving residues
must already sit at the reduced-level law; that is precisely
\eqref{eq:marg-consistency}, and it forces the conditional-Boltzmann
insertion uniquely. The seam-DB cross ratio is $1.000000$ for a
marginal-consistent build and $2.927$ for a marginal-inconsistent
one; the conditional-Boltzmann insertion satisfies DB to
$1.8\times 10^{-16}$ while bare-$\eqm$ insertion breaks it
($5.08\times 10^{-3}$). A randomised sweep ($2000$ trials at
$k=2,3$, random symmetric $H$) gives
$\max|\Delta_{\text{root}}|\le 1.8\times 10^{-15}$ when the lone
stationary equals the joint marginal, versus up to $1.52$ nats with
bare $\eqm$.

\subsubsection{Imposing reversibility determines the side potentials}
\label{sec:sinkhorn-suppl}

Returning to the side-potential joint \eqref{eq:sidepot-joint},
imposing marginal-consistency \eqref{eq:marg-consistency} promotes
$h$ from a free parameter (the A3 formulation, where $h$ was given
a Gaussian prior and sampled) to a function of $(\eqm,H)$ determined
by Sinkhorn / Iterative Proportional Fitting.

\begin{proposition}[Sinkhorn determination]
\label{prop:sinkhorn}
Imposing marginal-consistency \eqref{eq:marg-consistency} on
\eqref{eq:sidepot-joint}---i.e.\ requiring row-marginals $\eqm^{(c)}$
and column-marginals $\eqm^{(c')}$---is exactly the
Sinkhorn / IPF matrix-scaling problem for the strictly positive
kernel $K(a,b)=\eqm^{(c)}(a)\eqm^{(c')}(b)e^{-H_{cc'}(a,b)}$ with
prescribed marginals. By Sinkhorn's theorem the scaled joint
exists and is unique, and the potentials are unique up to the
gauge $h^{(c,c')}_a\mapsto h^{(c,c')}_a+\kappa_1$,
$h^{(c,c')}_b\mapsto h^{(c,c')}_b+\kappa_2$ (absorbed by the
normaliser). The coupling is preserved exactly: every
log-odds-ratio
$\log\frac{\pj(a,b)\pj(a',b')}{\pj(a,b')\pj(a',b)}=-(H(a,b)+H(a',b')
-H(a,b')-H(a',b))$
is invariant under the scaling, so $H$ carries correlation only and
the marginals are corrected without touching the interaction.
\end{proposition}

At $k=4$ with non-uniform $\eqm$ and non-trivial symmetric $H$:
base ($h=0$) marginal error $6.2\times 10^{-2}$; Sinkhorn restores
marginals to $1.1\times 10^{-16}$ in $18$ iterations; six random
initialisations give an identical joint ($\le 5.6\times 10^{-17}$)
with potentials differing only by additive constants (the gauge);
the interaction log-odds-ratios match base versus Sinkhorn to
$1.6\times 10^{-15}$. The side potentials collapse from free
prior-sampled parameters (A3) to a deterministic function of
$(\eqm,H)$, removing both the sampling cost and the
non-reversibility of A3 in one move.

\subsubsection{The impossibility triangle}
\label{sec:triangle-suppl}

Write $\mu_L^{(c,c')}(x)=\sum_y\pj^{(c,c')}(x,y)$ for the
partner-marginal at the first site of a pair, and $g(c)$ for the
lone stationary used when a paired site's partner is absent. Three
properties:
\begin{description}[leftmargin=2.2em]
  \item[(R)] \emph{Reversible / root-independent} $\Leftrightarrow$
    marginal-consistency: $g(c)=\mu_L^{(c,c')}$ and
    $g(c')=\mu_R^{(c,c')}$ for every co-occurring pair
    (\Cref{thm:revcond}).
  \item[(I)] \emph{Table-independent}: $g$ depends only on a site's
    own class, $g=g(c)$, with no partner index. A site ``chooses a
    table'', not ``to be paired''; the lone law of a
    paired-but-currently-alone site equals the singleton law of its
    class. This is the defining factorisation of the TKF-DP
    construction (keys i.i.d.\ and independent of the alignment, so
    the alignment marginal stays exactly TKF92).
  \item[(F)] \emph{Free per-pair marginal shift}: $\exists\,c,c_1\neq c_2$
    with $\mu_L^{(c,c_1)}\neq\mu_L^{(c,c_2)}$ (the marginal of a
    class genuinely depends on its partner; the biological motivation
    for $h$).
\end{description}

\begin{theorem}[Exactly two of three]
\label{thm:triangle}
No cap-2 model satisfies (R), (I), (F) simultaneously; each pair of
the three is realisable.
\end{theorem}

\begin{proof}
(R) pins $\mu_L^{(c,c')}=g(c)$ for every partner $c'$; (I) strips
the partner index from $g$; equating,
$\mu_L^{(c,c_1)}=g(c)=\mu_L^{(c,c_2)}$ for all partners, so the
partner-marginal is partner-independent and (F) fails. The argument
references only the marginal of the joint, so it is invariant to
how the coupling is parameterised.
\end{proof}

The three corners are all attainable: \textbf{A1}~$=$~(R,I) is the
Sinkhorn fix of \autoref{sec:sinkhorn-suppl} with $g=\eqm^{(c)}$
(coupling alive, $|H|_{\max}=2.63$; F dead to $10^{-13}$);
\textbf{A2}~$=$~(R,F) is reversible only with a partner-indexed
$g(c,c')$ (no class-only $g$ exists, so (I) fails);
\textbf{A3}~$=$~(I,F) is the original free-sampled-$h$ formulation
of \autoref{rem:side-potentials}, with lone $=\eqm^{(c)}$ but
root-violation $0.30$ (non-reversible). The model adopted in the
remainder of this appendix is A1.

\begin{remark}[What A1 keeps]
\label{rem:a1-keeps}
A1 does not make the two members of a pair symmetric. Because the
two sites draw their site-classes independently (the site-class DP,
distinct from the key/coupling DP), generically $c_i\neq c_j$, so
the pair marginalises to different per-class backgrounds
$\eqm^{(c_i)}\neq\eqm^{(c_j)}$. Marginal asymmetry between paired
sites, to the full extent their classes differ, is retained; only
the partner-context-dependence of a class's marginal is lost. The
two Dirichlet processes divide labour cleanly: the site-class DP
carries marginal composition, the key/Potts DP carries correlation.
\end{remark}

\begin{remark}[Why site classes are required under A1]
\label{rem:classes-required}
With a single site class every site shares one background $\eqm$,
so a reversible cap-2 coupling can express only correlation in the
marginal sense---never per-site compositional bias. All positional
amino-acid structure must then come from the site-class DP;
reversibility makes site classes a prerequisite for any composition
signal beyond pure pair correlation, and $H$ alone buys none of it
(it cancels in every log-odds-ratio, \Cref{prop:sinkhorn}).
\end{remark}

\subsection{The order rule: compensatory potentials of order $m{-}1$}
\label{sec:order-rule-suppl}

\begin{proposition}[No single-site repair beyond pairs]
\label{prop:capk}
For clusters of size $\ge 3$ with fixed pairwise Potts atoms,
single-site side potentials cannot enforce marginal-consistency.
Marginalising one site of a size-$3$ cluster induces an effective
two-site factor
$T(x_1,x_2)=\sum_{x_3}\eqm^{(c_3)}(x_3)
 e^{-H_{13}(x_1,x_3)-H_{23}(x_2,x_3)}$ whose log is generically
\emph{non-separable} in $(x_1,x_2)$, i.e.\ it injects interaction
into the $(k-1)^2$-dimensional interaction block. Single-site
potentials move only the $2(k-1)$ main effects (zero interaction),
so they cannot reproduce the exact $3\!\to\!2$ marginal on the same
atom.
\end{proposition}

Across $100$ random clusters at each of $k=3,4,5,20$, the
induced-interaction Frobenius norm exceeds $10^{-9}$ in all $100$
cases (up to $1.85$); the best IPF size-$2$ fit to the exact
$3\!\to\!2$ marginal leaves residual $\KL=1.1\times 10^{-2}$ nats.
\textbf{Cap-2 is therefore the maximal cluster reversible by
single-site repair}, beyond its computational motivation.

\begin{proposition}[Order rule]
\label{prop:orderrule}
Order-$r$ compensatory potentials restore marginal-consistency
(hence reversibility, \Cref{thm:revcond}) for clusters up to size
$r+1$, and are generically necessary beyond it. Equivalently: a
size-$m$ cluster carrying a full $m$-way Potts tensor stays
reversible iff it acquires compensatory potentials of every order
$\le m-1$, fitted Sinkhorn / IPF-style so that each lower-order
marginal equals the corresponding reduced-cluster stationary.
Single-site repair ($r=1$) thus reaches size $2$ (Sinkhorn,
\Cref{sec:sinkhorn-suppl}); pair repair ($r=2$) reaches size $3$;
and so on.
\end{proposition}

\begin{proof}[Mechanism]
Eliminating one site from a size-$m$ cluster whose interaction
graph is complete leaves an $(m-1)$-clique among the survivors
(graph fill-in): the induced log-marginal carries an effective
interaction of order exactly $m-1$ (the connected cumulant, scaling
as $(\text{coupling})^{m-1}$). A cluster whose potentials are
capped at order $r$ parameterises exactly the log-densities of
interaction order $\le r$, so it can represent that $(m-1)$-site
induced marginal iff $m-1\le r$. By exponential-family duality,
fixing all $\le(m-1)$-subset marginals determines a unique
max-entropy Gibbs distribution whose sufficient statistics are
exactly the $\le(m-1)$-subset indicators; the order-$(m-1)$
IPF/Sinkhorn fit therefore reproduces any realisable target
lower-marginals exactly, with the order-$m$ term set to zero,
while any cap below $m-1$ is missing the required sufficient
statistics. Size $3$ is the special case where the
$m\!\to\!(m-1)$ marginal lands on $2$ sites, whose most general
interaction is already pairwise.
\end{proof}

The induced log-marginal's top nonzero ANOVA order is exactly
$m-1$ for $m=3,4,5,6$, with the leading term scaling cubically as
coupling $\to 0$ at $m=4$ (connected-cumulant signature); IPF
reaches machine-zero full-table residual precisely when the fitted
order reaches $m-1$ and only then. Pairwise repair succeeds
$12/12$ at $m=3$ and fails generically at $m\ge 4$; order-$(m-1)$
is exact in every cell.

\paragraph{Why dynamic latent fields buy out of the hierarchy.}
Carrying the full order-$(m-1)$ potential hierarchy is the price of
a \emph{direct} (potential-based) reversible cluster
(\Cref{prop:orderrule}). A latent field
(\Cref{sec:dynamic-suppl}) is the same projective family written
through a hidden mediator: a categorical latent of cardinality
$\le k$ that decouples the pair (\Cref{prop:mediator}) replaces
the entire compensatory-potential tower, at the cost of an extra
sum over the latent variable in the per-cluster likelihood.

\subsection{Class-level variational substitution likelihood}
\label{sec:class-elbo}

Inference under the model requires, for each branch of a phylogeny
and each co-evolving DP class $C$, the per-branch transition
probability $P(\mathbf{x}_C(t) \mid \mathbf{x}_C(0), \{c_s\}_{s \in C})$.
Direct evaluation requires exponentiating a generator on
$\mathcal{A}^{|C|}$ and is intractable beyond $|C| \approx 5$. The
bulk of the key-DP partition is small-class for $\alpha_z$ in the
practical regime, so exact computation handles most of the work;
what remains is to control the tail of large classes.

We derive a variational lower bound from an exact path-measure
factorization on the class clique. The factorization is \emph{exact}
(Proposition below)---unlike the static pseudolikelihood of DCA
inference~\citep{ekeberg2013}, each factor still conditions site $s$ on
the \emph{full} neighbour trajectory $\sigma_{C\setminus s}$---and the
intractable joint stationary is absorbed into the conditioned-on initial
state. The genuine approximation, and with it the loss of direct-versus-
indirect coupling resolution, enters only through the mean-field
variational family below; see \autoref{rem:mf-recovers}.
Two variational families on the resulting Girsanov ELBO are natural.
The first is the simpler scheme we develop here, parameterized by
per-site distributions at internal tree nodes plus a single constant
rate matrix per site per branch, with a closed-form ELBO via
eigendecomposition and a Felsenstein-style pruning recursion for
optimization. The second is the fully inhomogeneous mean-field
scheme of Cohn et al.~\citep{cohn2010}, with time-varying rates and
forward-backward ODE sweeps, which we describe as an upgrade for
branches where the constant-rate scheme is too loose.

\begin{proposition}[Path-measure factorization]
Let $\sigma_C \in D([0,t]; \mathcal{A}^{|C|})$ be the trajectory of
a single-site reversible CTMC on class $C$ with rate
$Q^s(x_s \to x_s'; x_{C \setminus s})$ and stationary distribution
$\pi_C$. The path measure satisfies
\[
    \mathbb{P}(\sigma_C) \;=\;
    \pi_C\!\bigl(\sigma_C(0)\bigr)
    \prod_{s \in C} \mathbb{P}_s\!\bigl(\sigma_s \,\big|\, \sigma_{C \setminus s}\bigr),
\]
where $\mathbb{P}_s(\cdot \mid \sigma_{C \setminus s})$ is the law
of an inhomogeneous CTMC on site $s$ with rate
$Q^s\!\bigl(\,\cdot\, ; \, x_{C \setminus s}(u)\bigr)$ at time $u$.
\end{proposition}

\noindent The factorization is exact: simultaneous jumps at distinct
sites have zero Lebesgue measure, every jump is therefore
attributable to a unique site, and the path log-density splits
site-wise. The only intractable factor is $\pi_C(\sigma_C(0))$, which
is conditioned upon for transition probabilities and so cancels.

\subsubsection*{Constant-rate variational family}

\begin{remark}[Variational family]
The variational family is parameterized at two levels. At each
internal tree node $v$ and site $s$, a discrete distribution
$\nu_v^s \in \Delta^{|\mathcal{A}|-1}$ over residue states; node
distributions factorize over sites,
$\nu_v(\mathbf{x}_v) = \prod_s \nu_v^s(x_v^s)$. At each branch $b$
between nodes $p$ and $c$ of length $t_b$ and each site $s \in C$,
a constant rate matrix $\hat Q^s_b$ on $[0, t_b]$. Conditional on
sampled endpoint states $(x_p^s, x_c^s) \sim \nu_p^s \otimes \nu_c^s$,
the variational law on the branch is the CTMC bridge of
$\hat Q^s_b$ from $x_p^s$ to $x_c^s$, with sites independent given
endpoints. The joint variational law on the tree is
\[
    \mathbb{Q} = \prod_v \nu_v \cdot \prod_b \prod_{s \in C}
    \mathbb{Q}^s_b\bigl(\sigma^s_b \,\big|\, x_p^s, x_c^s\bigr).
\]
\end{remark}

\begin{proposition}[Closed-form per-branch ELBO]
The path-measure KL contribution from branch $b$ and site $s$ has
Girsanov form
\[
    \mathrm{KL}_{b, s} = \mathbb{E}_{\mathbb{Q}}\!\!\int_0^{t_b}\!\!
    \sum_{x, x'} q^s_b(x; u)
    \!\Bigl[\, \mathbb{E}_{u}\bigl[Q^s(x \to x'; x_{C \setminus s}(u))\bigr]
    - \hat Q^s_b(x \to x')
    + \hat Q^s_b(x \to x') \log\!\frac{\hat Q^s_b(x \to x')}{\exp \mathbb{E}_{u}[\log Q^s(\cdot)]} \Bigr]\, du,
\]
where $q^s_b(x; u) = \mathbb{E}_{x_p, x_c}[q^s_b(x; u \mid x_p, x_c)]$
is the variational bridge marginal of site $s$ at time $u$ averaged
over endpoint distributions, and the inner expectations
$\mathbb{E}_{u}[\cdot]$ are over the product of \emph{instantaneous}
neighbour bridge marginals at $u$,
$\prod_{s' \neq s} q^{s'}_b(\cdot; u)$. The integrand is therefore a
product of two time-varying bridge marginals at the same instant
$u$, and the integral over $[0, t_b]$ does \emph{not} factorise into
time-averages: the cross-temporal correlation between
$q^s_b(\cdot; u)$ and $q^{s'}_b(\cdot; u)$ contributes to
$\mathrm{KL}_{b,s}$ and is what makes the formula a strict lower
bound. The bridge-expectation closed form below evaluates the integral
correctly.

In the constant-rate family, the stationary point of
$\mathrm{KL}_{b,s}$ in $\hat Q^s_b$ admits no closed form in general
because of the cross-temporal correlation. The geometric-mean rate
\[
    \log \hat Q^{s, *}_b(x \to x') \;=\;
    \mathbb{E}_{\bar q^{C \setminus s}_b}\!\bigl[\log Q^s(x \to x'; x_{C \setminus s})\bigr],
\]
where $\bar q^{s'}_b(x) = (1/t_b) \int_0^{t_b} q^{s'}_b(x; u) du$
is the time-averaged neighbour distribution for each $s' \neq s$,
is the stationary point under the additional ``mean-field-of-bridges''
assumption that $q^s_b(\cdot; u)$ and $q^{s'}_b(\cdot; u)$ are
approximately uncorrelated in $u$ on $[0, t_b]$. This assumption is
exact at $H = 0$ (sites independent under the variational law as
well as under the truth), close to exact for small coupling, and
only mildly violated at typical phylogenetic time scales; we take
$\hat Q^{s, *}_b$ above as the chosen variational rate and evaluate
the strict $\mathrm{KL}_{b, s}$ at this rate via the bridge-expectation
formula in the Remark below. Each term of the resulting bound is
$\geq 0$ by Jensen, so $\mathrm{KL}_{b, s} \geq 0$ and the ELBO is a
strict lower bound on $\log P$, as required.
\end{proposition}

\begin{remark}[Closed-form ELBO via pairwise bridge expectations]
\label{rem:bridge-expectations}
Each rate $\hat Q^s_b$ admits an eigendecomposition
$\hat Q^s_b = U^s D^s (V^s)$ with $V^s = (U^s)^{-1}$ and eigenvalues
$\{\lambda^s_k\}$ (we drop the branch subscript $b$ for legibility).
The forward and backward propagators have the same expansion:
\[
    [\exp(\hat Q^s u)]_{x_p, x} = \sum_k U^s_{x_p, k}\, V^s_{k, x}\, e^{\lambda^s_k u},
    \quad
    [\exp(\hat Q^s (t-u))]_{x, x_c} = \sum_l U^s_{x, l}\, V^s_{l, x_c}\, e^{\lambda^s_l (t-u)},
\]
so the per-site bridge marginal is
\[
    q^s_b(x; u \mid x_p, x_c) =
    \frac{1}{P_s} \sum_{k, l} U^s_{x_p, k}\, V^s_{k, x}\,
    U^s_{x, l}\, V^s_{l, x_c}\, e^{\lambda^s_k u}\, e^{\lambda^s_l (t-u)},
\]
with $P_s = [\exp(\hat Q^s t)]_{x_p, x_c}$ the bridge normalising
constant. The branch integrand of $\mathrm{KL}_{b, s}$ becomes a sum
over $(k, l)$ and over $(k', l')$ from the neighbour expansion, with
pair-time integrals
\[
    \mathrm{HR}(\alpha, \beta, t) \;=\;
    \begin{cases}
    (e^{\alpha t} - e^{\beta t}) / (\alpha - \beta) & \alpha \neq \beta, \\
    t e^{\alpha t} & \alpha = \beta,
    \end{cases}
\]
which is the standard bridge-expectation~\citep{HolmesRubin2002}
expression. The full integral over $[0, t_b]$ is a finite sum of
$\mathrm{HR}$ evaluations and closed-form coefficients in
$U, V, \lambda$.
\end{remark}

\subsubsection*{Optimization: Felsenstein-like pruning over node distributions}

The node distributions $\nu_v^s$ optimise the ELBO under a
Felsenstein-style upwards--downwards pass: each $\nu_v^s$ is the
geometric mean of (i) the upwards-propagated likelihood from
$\nu_v^s$'s subtree, (ii) the downwards-propagated prior from the
rest of the tree, and (iii) the per-branch forward and backward
propagators of $\hat Q^s$ on each adjacent branch. The pass is
$O(\text{edges} \cdot A^2)$ per outer ELBO step, dominated by the
per-branch matrix exponential. The rate $\hat Q^{s, *}_b$ is then
re-computed under the fresh node distributions and the loop iterates
to convergence; the ELBO is monotone increasing in each step.

\subsubsection*{Jensen gap and the Cohn upgrade as endpoints of an
adaptive refinement}

The constant-rate scheme of the previous subsection and the
inhomogeneous mean-field scheme of Cohn et al.~\citep{cohn2010} are
not separate algorithms but the $N = 1$ and $N \to \infty$ ends of a
single adaptive refinement. Subdivide a branch of length $t_b$ into
$N$ sub-intervals at times $0 = u_0 < u_1 < \cdots < u_N = t_b$ with
constant rate $\hat Q^s_{b, k}$ on each $[u_{k-1}, u_k]$, and a
variational node distribution $\nu^s_k$ at each intermediate
midpoint. The resulting variational law is a piecewise-constant CTMC
bridge: a piecewise-constant approximation to whatever time-varying
rate Cohn's scheme would use at the same fidelity. As $N \to \infty$
with mesh refining, piecewise-constant rates become dense in
continuous time-varying rates, so the variational family becomes
Cohn's and the ELBO converges monotonically up. At the optimum, the
$\nu^s_k$ are pinned by stationarity to the bridge marginals of the
surrounding piecewise-constant CTMC --- the explicit-midpoint
parameterisation is slack at the optimum.

The Jensen-gap diagnostic per segment is the integrated AM--GM
discrepancy of $Q^s$ against neighbour bridge marginals at the
segment's resolution; it is closed form via the same eigendecomposition
machinery and tells the algorithm where to refine. For typical
phylogenetic branches ($t_b \in [0.1, 1]$) most segments need
$N = 1$; high-coupling segments may benefit from $N = 2$ or $3$;
only pathological branches need $N \to \infty$.
Forced segmentation by composition-change events from the augmented
indel history (\S\ref{sec:gravestone-scfg}) composes cleanly with
this Jensen-driven refinement: subdivide first at every
composition-change event, then optionally $N$-refine within each
composition-constant segment.

\begin{remark}[Cohn's inhomogeneous fixed point as the $N \to \infty$ limit]
For a composition-constant segment requiring $N \to \infty$ refinement,
the limit is Cohn et al.'s full inhomogeneous fixed point: a
time-varying rate $\tilde Q^s(\cdot; u)$ obtained by alternating
forward sweeps $\dot q_s = q_s \tilde Q^s$ from $\delta_{x_p^s}$,
backward sweeps $\dot \rho_s = -\tilde Q^s \rho_s$ from
$\delta_{x_c^s}$, and rate updates
$\log \tilde Q^{s,*}(\cdot; u) = \mathbb{E}[\log Q^s(\cdot)]$
pointwise in $u$, with bridge marginals
$q^{\rm bridge}_s(x; u) \propto q_s(x; u)\, \rho_s(x; u)$.
The midpoint construction above is the discretisation of this
fixed point at mesh size $u_k - u_{k-1}$; the Cohn limit is recovered
as the mesh shrinks to zero. Cohn et al.~\citep[\S 7]{cohn2010}
also give the internal-node rule used by Felsenstein pruning under
the time-varying rate.
\end{remark}

This is the lower-bound construction referenced in the Discussion for
Potts components of size $> 2$: independent-coupling pairwise factorisation
gives the $N{=}1$ lower bound, midpoint augmentation extends it to a
formal ELBO, and $N \to \infty$ recovers the Cohn fixed point.
Classes whose Jensen gap remains large after refinement are
candidates for the $k$-clique cluster-variational upgrade of
Linzner \& Koeppl~\citep{linznerkoeppl2018}, or for fallback to
exact uniformization-based bridge sampling via Rao \&
Teh~\citep{raoteh2013}.

\begin{remark}[What the mean-field bound recovers: direct vs.\ indirect
coupling]
\label{rem:mf-recovers}
Because the variational family factorises over sites
($\mathbb{Q}=\prod_s \mathbb{Q}^s$), the coupling enters the ELBO only
through the rank-$1$ outer product of per-site bridge marginals,
\[
  \sum_{s<s'}\int_{\mathcal T}\sum_{a,b}
     H_{ss'}(a,b)\,\gamma_s(a;u)\,\gamma_{s'}(b;u)\,du,
\]
and the connected correlation
$\mathrm{Cov}_{\mathbb Q}\!\bigl(\mathbb 1_{x_s=a},\mathbb 1_{x_{s'}=b}\bigr)$
vanishes identically. Consequently
$\partial\,\mathrm{ELBO}/\partial H_{ss'}$ matches, per pair, that pair's
own endpoint-conditioned marginals and nothing else. This is
\emph{naive mean-field learning}: the bound never forms the connected
covariance $C$ nor inverts it, so it performs neither the
fluctuation-response deconvolution $J\approx -C^{-1}$ of mean-field
DCA~\citep{Morcos2011} nor the all-sites conditioning of pseudolikelihood
DCA~\citep{ekeberg2013}. The self-consistency of the E-step couples the
$\gamma_s$ through the \emph{forward} mean-field map $H\mapsto\gamma$
(cavity fields $\sum_{s'}H_{ss'}\gamma_{s'}$); linear response is the
Jacobian $\mathrm{d}\gamma/\mathrm{d}h=\chi=C$, and disentangling requires
forming \emph{and inverting} $\chi$---coordinate-ascent variational EM uses
only the fixed-point values $\gamma$, never $\mathrm{d}\gamma/\mathrm{d}h$,
so no implicit inversion is taken~\citep{kappenrodriguez1998,tanaka1998}.
Therefore, whenever a site has $\ge 2$ coupled partners (a genuine
$m\ge 3$ contact triangle), the fitted $H_{ss'}$ absorbs the indirect
$s$--$k$--$s'$ correlation, and the bound recovers indirect-inclusive
(mutual-information--like) couplings, not direct ones. The exact path
factorisation is not to blame---it conditions on the full correlated
neighbour trajectory and is plmDCA-like---the loss is incurred by the
mean-field product family. At genuine cap-$2$ (no triangles) the
distinction is moot and the bound is exact.
\end{remark}

\subsection{Composite-likelihood phylogenetic ELBO for multi-site Potts
coupling under a Metropolis--Hastings substitution form}
\label{sec:composite-mh-elbo}

The class-level bound of \autoref{sec:class-elbo} controls the coupled
CTMC on a single clique of arbitrary size $m$. Two costs remain for a
general phylogeny. First, the \emph{joint} clique likelihood on a tree
carries an $\mathcal{A}^{m}$-state conditional-likelihood vector at every
node, super-exponential in tree depth. Second, reversibility of the
joint stationary $\eqm_C(x_C)\propto\prod_s\eqm^{(c_s)}(x_s)\,
\exp(-\sum_{\{s,s'\}\subset C}H_{c_sc_{s'}}(x_s,x_{s'}))$ across the
insertion/deletion seam forces the order-$(m{-}1)$ potential hierarchy of
\autoref{prop:orderrule} and the Sinkhorn side-potential fix of
\autoref{prop:sinkhorn}. This subsection shows that a
\emph{composite} (pairwise) likelihood dissolves both costs at once: the
corpus objective becomes a \emph{product over coupled pairs of independent
two-site phylogenetic trees}, the Potts coupling enters \emph{only} through
the two-site stationary, and---under a Metropolis--Hastings substitution
form---the per-pair generator, its exact reversibility, and its
Holmes--Rubin bridge statistics all collapse to a single shared
exchangeability acting on single-residue moves.

\paragraph{Composite decomposition.} Let $E$ be the set of coupled
site-pairs: under the key DP this is the set of two-element blocks of the
partition (the cap-$2$ TKF-DP training regime, where every co-evolving
class has $|C|\le 2$); under structure supervision it is a fixed contact
map. Replace the joint clique likelihood on the phylogeny $\mathcal{T}$
by the composite (pseudo-)likelihood
\begin{equation}
  \mathcal{L}_{\mathrm{comp}}
    \;=\; \sum_{(i,j)\in E}\;
      \log P\!\bigl(\mathbf{x}_i,\mathbf{x}_j \;\big|\; \mathcal{T},\,
        \eqm^{(c_i)},\eqm^{(c_j)},\,H_{c_ic_j},\,S\bigr),
  \label{eq:comp-objective}
\end{equation}
a composite \emph{marginal} likelihood~\citep{lindsay1988} (see
\citealp{varin2011} for the general theory). This is the
product-of-\emph{marginals} form, and is to be distinguished from
Besag's product-of-\emph{conditionals} pseudolikelihood~\citep{besag1975}
(the object underlying plmDCA~\citep{ekeberg2013}): the two agree only in
the absence of coupling triangles. Each summand is the phylogenetic
likelihood of a single ordered residue pair $(x_i,x_j)$ evolving down
$\mathcal{T}$ under a two-site reversible CTMC. The Potts energy is
\emph{genuinely pairwise}, so the composite reproduces the model's energy
\emph{form}; but each $H_{ij}$ is fitted from the pair \emph{marginal}
$\eqm_{ij}$, which is exact for the direct coupling only when $i,j$ share
no common partner (\autoref{rem:mf-recovers}).

\paragraph{The two-site chain and its stationary.} Index the pair state
$x=(a,b)\in\mathcal{A}^2$ by $a$ (site $i$), $b$ (site $j$), and write
$H:=H_{c_ic_j}$. The composite factor treats both sites as present
throughout the branch (no partner insertion/deletion \emph{inside} the
factor), so the only requirement on its generator is reversibility with
respect to the pair stationary
\begin{equation}
  \eqm_{ij}(a,b)\;=\;\frac{1}{Z_{ij}}\;
     \eqm^{(c_i)}(a)\,\eqm^{(c_j)}(b)\,e^{-H(a,b)},
  \qquad
  Z_{ij}=\!\!\sum_{a,b}\eqm^{(c_i)}(a)\,\eqm^{(c_j)}(b)\,e^{-H(a,b)}.
  \label{eq:comp-pair-stat}
\end{equation}
This is exactly the pair-marginalised joint stationary
\eqref{eq:base-joint-suppl}. The marginal-consistency condition of
\autoref{thm:revcond}---$\sum_b\eqm_{ij}(a,b)=\eqm^{(c_i)}(a)$---is
\emph{not} imposed here, because it is an indel-seam requirement (needed
only when a partner is born or dies mid-branch, coupling a singleton
stationary to a pair stationary). Fixing both sites present within the
composite factor removes it, and with it the entire order-$(m{-}1)$
hierarchy and the Sinkhorn correction: the pair chain needs only detailed
balance for \eqref{eq:comp-pair-stat}, which any GTR or Metropolis
instance provides.

\paragraph{Metropolis--Hastings generator.} Take the single-transition
(one residue changes at a time) Metropolis--Hastings generator with a
single shared symmetric exchangeability $S$ and Hastings acceptance
$g(u,v)=\min(1,v/u)$,
\begin{align}
  Q_{ij}\bigl[(a,b),(a',b)\bigr]
     &= S(a,a')\,\min\!\Bigl(1,\;\tfrac{\eqm_{ij}(a',b)}{\eqm_{ij}(a,b)}\Bigr),
     & a'\neq a \quad(\text{site-}i\text{ move}),
     \label{eq:comp-mh-imove}\\
  Q_{ij}\bigl[(a,b),(a,b')\bigr]
     &= S(b,b')\,\min\!\Bigl(1,\;\tfrac{\eqm_{ij}(a,b')}{\eqm_{ij}(a,b)}\Bigr),
     & b'\neq b \quad(\text{site-}j\text{ move}),
     \label{eq:comp-mh-jmove}\\
  Q_{ij}\bigl[(a,b),(a',b')\bigr] &= 0,
     & a'\neq a,\; b'\neq b,
     \label{eq:comp-mh-double}
\end{align}
with $Q_{ij}[x,x]=-\sum_{y\neq x}Q_{ij}[x,y]$; simultaneous double jumps
have zero rate, as in any single-transition (CTBN-style) pair chain.
Three properties of the Hastings kernel do the work.

\begin{proposition}[Exact reversibility, spectator cancellation, symmetric
min-flux]
\label{prop:comp-mh}
For any symmetric $S$ and any $H$, the generator
\eqref{eq:comp-mh-imove}--\eqref{eq:comp-mh-double} satisfies detailed
balance with respect to \eqref{eq:comp-pair-stat}, with the symmetric flux
\begin{equation}
  F_{xy}\;:=\;\eqm_{ij}(x)\,Q_{ij}[x,y]
    \;=\; S_{xy}\,\min\!\bigl(\eqm_{ij}(x),\eqm_{ij}(y)\bigr)
    \;=\;F_{yx}.
  \label{eq:comp-min-flux}
\end{equation}
Moreover the acceptance ratio of a single-site move depends on the
Potts coupling only through the \emph{local} increment
$\Delta H_i(a\!\to\!a';b)=H(a',b)-H(a,b)$: the spectator stationary
cancels,
\begin{equation}
  \frac{\eqm_{ij}(a',b)}{\eqm_{ij}(a,b)}
    \;=\;\frac{\eqm^{(c_i)}(a')}{\eqm^{(c_i)}(a)}\;
        e^{-\Delta H_i(a\to a';\,b)},
  \label{eq:comp-spectator}
\end{equation}
so the site-$i$ block of $F$ carries the spectator as an overall factor,
$F[(a,b),(a',b)]=Z_{ij}^{-1}\,S(a,a')\,\eqm^{(c_j)}(b)\,
\min\!\bigl(\eqm^{(c_i)}(a)e^{-H(a,b)},\,\eqm^{(c_i)}(a')e^{-H(a',b)}\bigr)$,
and symmetrically for the site-$j$ block.
\end{proposition}

\begin{proof}
Detailed balance: $\eqm_{ij}(x)Q_{ij}[x,y]
=S_{xy}\,\eqm_{ij}(x)\min(1,\eqm_{ij}(y)/\eqm_{ij}(x))
=S_{xy}\min(\eqm_{ij}(x),\eqm_{ij}(y))$, manifestly symmetric in $x,y$
(and $=0$ off single-transition edges, where $S_{xy}=0$).
\eqref{eq:comp-spectator} is immediate from \eqref{eq:comp-pair-stat} since
$\eqm^{(c_j)}(b)$ appears in numerator and denominator; the flux form
follows by multiplying through by $\eqm_{ij}(a,b)=\eqm^{(c_j)}(b)\,
\eqm^{(c_i)}(a)e^{-H(a,b)}/Z_{ij}$.
\end{proof}

\emph{What M--H buys.} Compared with the F81 pair
$Q_{xy}\propto S_{xy}\eqm_{ij}(y)$---which is \emph{also} reversible with
respect to \eqref{eq:comp-pair-stat}, so reversibility alone does not
distinguish them---the Hastings form gives (i) the symmetric min-flux
\eqref{eq:comp-min-flux}, whose exchangeability is a \emph{single} shared
$S$ (the $190$-parameter symmetric object of the pair-model family, versus
$\mathcal{O}(K_c^2A^2)$ free Potts entries), with a closed-form
usage/exposure M-step below; (ii) the spectator cancellation
\eqref{eq:comp-spectator}, which localises $H$ to the increment
$\Delta H_i$ so the pair generator is assembled from a per-site
exchangeability and a Potts-modulated acceptance rather than a dense
$\mathcal{A}^2\times\mathcal{A}^2$ rate table; and (iii) support of $F$ (hence
of the Holmes--Rubin usage $N$) \emph{only} on the $2\mathcal{A}(\mathcal{A}-1)$
single-transition edges, so the bridge statistics never populate the
$\mathcal{A}^2$ dense block. (The symmetric-Metropolis $\sqrt{\cdot}$ variant of
the remark in \autoref{sec:gtr-potts} instead buys a $\eqm$-independent
symmetrised eigenbasis; the $\min$ variant here buys the exact acceptance
and the sparsest flux.)

\paragraph{The composite phylo-ELBO objective.} Each summand of
\eqref{eq:comp-objective} is the reversible-CTMC Felsenstein forward on the
$\mathcal{A}^2$-state pair chain $Q_{ij}$ over $\mathcal{T}$: a post-order
recursion of pair conditional-likelihood vectors with per-branch
propagators $P_{ij}(\tau)=\exp(Q_{ij}\tau)$ from the reversible
eigendecomposition of \autoref{rem:bridge-expectations}. The corpus
objective is thus a \emph{product of pair-trees},
\begin{equation}
  \mathcal{L}_{\mathrm{comp}}
    \;=\;\sum_{(i,j)\in E}\log\!\Bigl(
      \textstyle\sum_{x_r}\eqm_{ij}(x_r)\,F^{\,r}_{ij}(x_r)\Bigr),
  \qquad
  F^{\,v}_{ij}(x_v)=\!\!\prod_{u\in\mathrm{ch}(v)}\!\!
     \bigl[P_{ij}(\tau_u)\,F^{\,u}_{ij}\bigr](x_v),
  \label{eq:comp-forward}
\end{equation}
with leaf CLVs $F^{\ell}_{ij}(a,b)=\mathbb{I}[a{=}x^{\mathrm{obs}}_{i,\ell}]\,
\mathbb{I}[b{=}x^{\mathrm{obs}}_{j,\ell}]$. Each tree runs on $\mathcal{A}^2$
states independently of every other pair; no $\mathcal{A}^m$ state and no
$|\Theta|$ product ever appears. Where a co-evolving clique genuinely has
$m>2$ sites, each pair term may be further mean-fielded by the
product-of-single-site-trees family of \autoref{sec:class-elbo} (constant-rate
Cohn/Girsanov variational generators per site, coupled only through the
pairwise bridge expectations $\expect[\Delta H]$), yielding a strict lower
bound $\mathcal{L}_{ij}\le\log P(\text{pair}_{ij})$; at $m=2$ that gap is zero
because \eqref{eq:comp-forward} keeps the exact two-site chain. Hence at
cap-$2$, $\mathcal{L}_{\mathrm{comp}}=\log P_{\mathrm{comp}}(\text{data})$
exactly---the ELBO is tight and equals the composite log-likelihood.

\paragraph{E-step: Holmes--Rubin sufficient statistics.} For each pair and
each branch $u\!\to\!v$ of length $\tau$, one reversible forward--backward
pass over $Q_{ij}$ gives the endpoint-conditioned edge posterior
$w^{uv}_{ij}(x,y)$ (the two-slice variational marginal of the pair state at
the branch endpoints). The Holmes--Rubin bridge~\citep{HolmesRubin2002},
evaluated by the reversible eigendecomposition and the divided-difference
kernel $\mathrm{HR}(\alpha,\beta,\tau)$ of \autoref{rem:bridge-expectations},
returns the expected transition usage and dwell for that branch, aggregated
over pairs and branches into
\begin{equation}
  N_{xy}=\!\!\sum_{(i,j)\in E}\sum_{u\to v}
     \expect\!\bigl[\#\{x\!\to\!y\text{ jumps}\}\;\big|\;w^{uv}_{ij}\bigr],
  \quad
  T_x=\!\!\sum_{(i,j)\in E}\sum_{u\to v}
     \expect\!\bigl[\text{time in }x\;\big|\;w^{uv}_{ij}\bigr].
  \label{eq:comp-hr}
\end{equation}
By \autoref{prop:comp-mh}, $N_{xy}$ is supported only on single-transition
edges. These are the complete-data sufficient statistics of the pair chain.

\paragraph{M-step for $S$, $\eqm$/$H$.} The expected complete-data
log-likelihood is $\sum_{x\neq y}\bigl[N_{xy}\log Q_{ij}[x,y]-T_x\,
Q_{ij}[x,y]\bigr]$ up to $\eqm$-only terms. Two blocks alternate.
\emph{(i)~Exchangeability.} All single-transition edges sharing an
unordered residue pair $\{a,a'\}$ (across both coordinates, pooled by the
component-exchange symmetry $\eqm_{ij}(a,b)=\eqm_{ji}(b,a)$) carry the same
$S(a,a')$; stationarity in $S(a,a')$ gives the closed form
\begin{equation}
  S(a,a')\;\leftarrow\;\frac{\;\sum_{\text{edges}\,\{a,a'\}}N_{xy}\;}
        {\;\sum_{\text{edges}\,\{a,a'\}}T_x\,\min\!\bigl(1,\eqm_{ij}(y)/\eqm_{ij}(x)\bigr)\;},
  \label{eq:comp-mstep-S}
\end{equation}
usage over exposure---the exact maximiser, hence non-decreasing in the
expected complete-data objective. \emph{(ii)~Stationary/coupling.} With the
per-class profiles $\eqm^{(c)}$ frozen (LG08 or the archetype vocabulary
\eqref{eq:arch-mapping}), the free coupling is $H$; equivalently one may fit
the symmetric pair stationary $\eqm_{ij}$ directly and read
$H(a,b)=-\log\eqm_{ij}(a,b)+\log\eqm^{(c_i)}(a)+\log\eqm^{(c_j)}(b)+\text{const}$,
identified up to the usual Potts gauge (additive functions of $a$ alone or
$b$ alone, absorbable into $\eqm^{(c_i)},\eqm^{(c_j)}$). Because the $\min$
acceptance makes $\eqm_{ij}$ enter $Q_{ij}$ nonlinearly, there is no closed
form; a damped gradient ascent on $\log\eqm_{ij}$ of the same expected
complete-data objective (with residual $N_{xy}-T_xQ_{ij}[x,y]$) increases it.
\emph{Monotonicity.} Since \eqref{eq:comp-mh-imove}--\eqref{eq:comp-mh-double}
is reversible with respect to $\eqm_{ij}$ for \emph{every} iterate
$(S,\eqm_{ij})$ (\autoref{prop:comp-mh}), each pair term of
\eqref{eq:comp-forward} remains a valid observed-data likelihood of a
reversible CTMC, and the standard EM inequality applies per pair: the E-step
computes exact expected sufficient statistics \eqref{eq:comp-hr} under the
current $Q_{ij}$ and the M-step raises the expected complete-data
log-likelihood (exactly in $S$, ascent in $\eqm_{ij}$), so
$\mathcal{L}_{\mathrm{comp}}$ is non-decreasing under coordinate ascent.

\paragraph{What the composite approximates, and where it is exact.} At
cap-$2$ (every co-evolving class has $|C|\le2$, the TKF-DP training
geometry) the set $E$ is the class partition itself, the composite objective
\emph{is} the joint likelihood, and each pair-tree forward
\eqref{eq:comp-forward} is exact, so $\mathcal{L}_{\mathrm{comp}}=\log
P(\text{data})$ with no gap. For genuine $m>2$ cliques the product over the
$\binom{m}{2}$ pairs is a composite \emph{marginal}
likelihood~\citep{lindsay1988,varin2011}: it retains the pairwise energy
\emph{form} (the model has no higher-body term to lose), but each $H_{ij}$
is fitted from the pair \emph{marginal} $\eqm_{ij}$. When $i$ and $j$ share
a common partner $k$, that marginal absorbs the indirect $i$--$k$--$j$
correlation, so the estimator converges to the indirect-inclusive
(mutual-information--like) coupling, \emph{not} the direct one---it is
\emph{inconsistent} for the direct $H_{ij}$, and the failure is genuine, not
merely a loss of Godambe efficiency. This is the marginal-composite
counterpart of the mean-field limitation of \autoref{rem:mf-recovers};
contrast Besag's conditional pseudolikelihood (the plmDCA
object~\citep{ekeberg2013,besag1975}), which conditions each site on all
others and \emph{is} consistent for the direct couplings, and mean-field
DCA~\citep{Morcos2011}, which recovers direct couplings by inverting the
connected covariance. The two-site tree forwards themselves remain exact;
the bias is entirely the per-pair marginalisation. Thus the composite
phylo-ELBO is exact where TKF-DP trains (cap-$2$, no triangles), and above
it a controlled but \emph{direct-indirect-conflating} surrogate that reduces
the multi-site Potts objective to a product of independent two-site trees
with the coupling entering only through the pair stationary
\eqref{eq:comp-pair-stat}.

\subsection{Dynamic latent-field variant}
\label{sec:dynamic-suppl}

The Potts construction of \autoref{sec:gtr-potts} routes coupling
through an explicit pairwise residue-space tensor $H_{cc'}$ and pays
for arbitrary-size reversibility with the order-$(m-1)$ potential
hierarchy of \autoref{prop:orderrule}. An alternative construction
routes the coupling instead through a shared categorical
\emph{latent field} per coupled cluster, evolving on the tree as a
small reversible CTMC. The two sites of a coupled pair are then
conditionally independent given the field at every time; the cost
of arbitrary cluster size is one extra sum (per cluster) over the
small latent alphabet rather than a hierarchy of compensatory
potentials. The construction is reversible by induction, exactly
tractable at the cap-2 cluster geometry used in TKF-DP training,
and extends to general phylogenies via a structured mean-field
variational bound (\autoref{sec:variational-suppl}). Its closest
phylogenetic antecedents are the multi-level covarion / nested-iHMM
families~\citep{tuffleysteel1998, Beal2002, TehJordanBealBlei2006}
and the F81/parent-independent-mutation foundations of population
genetics~\citep{ewens1972, Kingman1975, Felsenstein81}.

\subsubsection{Categorical mediator}
\label{sec:dynamic-mediator}

\begin{proposition}[Categorical mediator]
\label{prop:mediator}
Any pair joint $\pj(x,y)$ on $\mathcal A\times\mathcal A$ can be
written $\pj(x,y)=\sum_{h} w_h\,a_h(x)\,b_h(y)$ with $\{a_h\}$,
$\{b_h\}$ probability vectors and $\{w_h\}$ a probability---i.e.\
there is a latent $h$ rendering the two sites conditionally
independent, $x\perp y\mid h$. The minimal cardinality $|h|$ is the
\emph{nonnegative rank} $\nnrank(\pj)\le k$: it is $1$ iff the
coupling is separable (independent sites), and up to $k$ for
full-rank coupling. Computing $\nnrank$ is NP-hard in
general~\citep{Vavasis2009}; for $k\le 3$ Cohen--Rothblum gives
$\nnrank=\operatorname{rank}$~\citep{CohenRothblum1993}, and the
first gap appears at $k=4$ (a strictly-positive symmetric size-$4$
joint of ordinary rank $3$ can have nonnegative rank $4$).
\end{proposition}

A Potts edge ``is'' a hidden common cause: a categorical field of
size $\nnrank$ that decouples the pair. In cap-2 every edge is its
own small ($\le k$-state) mediator.

\paragraph{Gaussian copula is a strict subfamily.} A scalar latent
Gaussian per site (one correlation $\rho$) induces only
\emph{monotone} association. For binary alphabets that is the whole
story; for $k\ge 3$ a Potts has $\binom{k}{2}$ coupling degrees of
freedom and can be non-monotone or cyclic, which no single-$\rho$
copula reproduces. Matching arbitrary $H$ needs a \emph{vector}
latent Gaussian (rank up to $k-1$) plus general links, i.e.\ a
latent-factor model.

\subsubsection{Instantly-re-equilibrating dynamic latent-field selector}
\label{sec:reequil-suppl}

A categorical latent field $\theta(t)$ evolves on the tree as a
small reversible CTMC with state space the number of fields
$L \ll q^m$; a coupled site evolves under the field-conditioned
single-site generator $\Q(\theta)$ between field jumps and
\emph{instantly re-equilibrates}---resamples from $f(\cdot\mid\theta')$---at
each jump. This is not a continuous-time Bayesian network (the
residue process has instantaneous jumps slaved to $\theta$), but it
is exactly solvable at the cap-2 cluster geometry used in TKF-DP
training; on general phylogenies the joint multi-component
Felsenstein admits the variational rescue of
\autoref{sec:variational-suppl}.

\begin{proposition}[Hierarchical Felsenstein]
\label{prop:hier-fels}
Resampling at every field jump makes residues conditionally
independent across sites \emph{and} memoryless across jumps given
the $\theta$-trajectory. Hence conditional on a $\theta$-trajectory
$\bm{\theta}$ on the tree, the joint likelihood factorises over
sites,
\begin{equation}
  L(\text{obs}\mid\bm{\theta}) = \prod_{s} L_s(\text{obs at site }s\mid\bm{\theta}).
\end{equation}
At the cap-2 cluster geometry (one internal parent, two leaves;
the cherry geometry of TKF-DP training and IPHMM pairwise
postprocessing), integrating $\bm{\theta}$ via a single Felsenstein
step over the field alphabet gives the exact joint cluster emission
at cost
$O(\mathrm{edges}\cdot L^2 + \mathrm{edges}\cdot L\cdot m\cdot k^2)$,
\emph{linear in cluster size}~$m$. For phylogenies of depth $>1$,
the per-node conditional-likelihood representation under exact
marginalisation expands non-linearly with tree depth; a structured
mean-field ELBO (\autoref{sec:variational-suppl}) restores
tractable linear-in-$m$ cost on general trees and is tight at
cap-2 and at the $\rho_{\text{chain}}\to\{0,\infty\}$ rate limits.
\end{proposition}

The $\theta$-CTMC is reversible by construction; residues are
conditionally independent given $\theta$, so masking is trivial;
and sites read the field but never a partner directly, so the
construction is \emph{table-independent} in the sense of (I) in
\autoref{thm:triangle}. The expressiveness knob is $L$ (the field
truncation cap) and the breadth of $f(\cdot\mid\theta)$ (the
per-field residue stationary spread).

\subsubsection{F81-on-DP: a Dirichlet-process field selector}
\label{sec:f81-dp}

Let the field $\theta(t)$ evolve by F81,
$Q(\theta\!\to\!\theta')=\rho_{\mathrm{chain}}\,\eqm_{\theta'}$ for
$\theta\neq\theta'$, with $\eqm$ the stick-breaking weights of a
Dirichlet process and $\rho_{\mathrm{chain}}$ a scalar rate
multiplier. This is reversible w.r.t.\ $\eqm$ conditional on the DP
draw; F81's defining feature---jump to the destination type
independently of the source---is exactly parent-independent
mutation (PIM), whose stationary allelic partition in the
infinite-alleles model is the Ewens sampling formula equivalent to
the DP / CRP~\citep{ewens1972, Kingman1975, EthierKurtz1981}. It
is the continuous-time cousin of the infinite HMM and HDP-HMM
families~\citep{Beal2002, TehJordanBealBlei2006}. Tractability is
exact: a finite tree occupies finitely many atoms, and F81 is
\emph{strongly lumpable} onto $\{\text{occupied atoms}\}\cup\{\text{one tail block}\}$, so the
occupied-atoms-plus-tail field Felsenstein is an exact collapse of
the infinite-state recursion.

\paragraph{Per-state rate: the reversible form.} A naive per-state
exit-rate multiplier $Q(\theta\!\to\!\theta')=\eta_\theta\,\eqm_{\theta'}$
is \emph{not} reversible (detailed-balance residual
$\eqm_\theta\eqm_{\theta'}(\eta_\theta-\eta_{\theta'})$; the chain
stays $\eqm$-stationary but is irreversible). The reversible way to
vary rate by field-type is a rank-1 symmetric exchangeability
$S_{\theta\theta'}=a_\theta a_{\theta'}$,
$Q(\theta\!\to\!\theta')=a_\theta a_{\theta'}\eqm_{\theta'}$: the
exit rate is $\propto a_\theta$ (a genuine per-state multiplier)
but jumps land in the $a$-tilted $\eqm$
($\propto a_{\theta'}\eqm_{\theta'}$). The clean dichotomy:
constant rate $\Leftrightarrow$ bare-$\eqm$ jump $\Leftrightarrow$
exact Ewens marginal; per-state rate $\Leftrightarrow$ tilted jump
$\Leftrightarrow$ non-Ewens marginal.

\paragraph{Carry vs.\ resample under a modulator jump.} When any
modulator $\theta$ jumps and resets the modulated layer's
generator, the modulated state $x$ is either carried (kept) or
resampled. Detailed balance across the jump leaves the residual
factor $\eqm(x\mid\theta)$ vs.\ $\eqm(x\mid\theta')$: carrying is
reversible iff the modulator preserves the modulated stationary
(changes only \emph{rate}, not \emph{distribution}---the
Tuffley--Steel covarion regime, \citealp{tuffleysteel1998}). If the
modulator changes the equilibrium (the expressive case here),
carrying breaks detailed balance and one must resample $x$ from
$\eqm(\cdot\mid\theta')$. The dynamic-field variant therefore
mandates resampling at each field jump.

\subsubsection{Dynamic site classes: a nested hierarchy}
\label{sec:dynamic-classes}

The same construction layers. A slower per-site reversible DP-CTMC
(an \emph{evolving profile}) sits above the field selector: class
$\to$ field $\to$ residue, each an F81-on-DP modulating the next---a
multi-level covarion / nested-iHMM. By induction the composition
is reversible, table-independent, and de-Finetti-safe under
indels; tractability nests, with the per-node CLV a tensor over
(occupied classes $\times$ occupied fields). The natural timescale
ordering residue $\gg$ field $\gg$ class (point substitution
fastest, coevolutionary regime shifts slower, site identity
slowest) aids identifiability but is not required for
reversibility.

The two dynamic DP layers are independent options: a dynamic
\emph{field} selector alone gives an evolving coevolutionary
coupling over background-composed sites; a dynamic
\emph{site-class} selector alone gives evolving per-site
composition with no coupling; the general dynamic latent-field
variant carries both.

\subsubsection{Archetype vocabulary for per-(class, field) stationaries}
\label{sec:archetype-suppl}

We parameterise the per-(class, field) residue stationary
$\eqm^{(c,\theta)}$ through a shared vocabulary of simplex points
$\{\eqm^{(k)}\}_{k=1}^{K_a}$ (``archetypes''): each $(c, \theta)$ is
mapped to one archetype by a discrete assignment,
\begin{equation}
  \eqm^{(c,\theta)} \;=\; \eqm^{(\arch(c,\theta))},
  \qquad \arch: \{1,\dots,K_c\}\times\{1,\dots,L_{\max}\} \to \{1,\dots,K_a\}.
  \label{eq:arch-mapping}
\end{equation}
This shared-vocabulary constraint is what makes the field selector
$\theta$ act as a \emph{coupling mediator}. If instead each
$\eqm^{(c,\theta)}$ were an unconstrained free simplex point, the model
could pick each independently and a field flip would often produce a
near-identical emission --- no salt-bridge-like covariation would
emerge. A site class is thus a \emph{pattern} of archetype selections
across $\theta$. Some classes ($\arch(c,\cdot) \equiv k$) are
specializers: the field flip is a no-op. Others (e.g.\
$(k_+, k_-, k_+, k_-)$ with $k_+$ positive-charge and $k_-$
negative-charge archetypes) are mediators: a field flip is a discrete
archetype swap that swings the emission across biophysically distinct
simplex points. Because each archetype retains a full 20-simplex mass,
standard GTR substitution physics within an archetype is untouched ---
no expressiveness is sacrificed. Priors: $\eqm^{(k)} \sim
\mathrm{Dirichlet}(\alpha_{\text{prior}}\mathbf 1)$; $\arch(c,\theta)
\sim \mathrm{Cat}(\rho_a)$ with $\rho_a$ a truncated stick-breaking
draw with concentration $\alpha_{\arch}$.

\paragraph{Compound CTMC state and dynamics.}
Consider a cluster of $N$ sites of classes $(c_1,\dots,c_N)$ sharing a
single field selector $\theta$. The compound state $(\theta, x_1, \dots,
x_N)$ evolves according to the following generator:
\begin{align}
  Q_\theta[(\theta,\mathbf{x}), (\theta',\mathbf{x}')]
    &= \rho_{\text{chain}}\,\rho[\theta']\,\prod_{n=1}^{N}\eqm^{(k(c_n,\theta'))}[x'_n],
    &\theta' \ne \theta,
    \label{eq:arch-Q-field}
  \\
  Q^{(n)}[(\theta,\mathbf{x}), (\theta,\mathbf{x}\oplus_n y)]
    &= \exch[x_n, y]\,\eqm^{(k(c_n,\theta))}[y],
    &y \ne x_n,
    \label{eq:arch-Q-res}
\end{align}
where $\mathbf{x}\oplus_n y$ denotes the residue tuple $\mathbf{x}$
with position $n$ set to $y$ and other positions unchanged.
Equation~\eqref{eq:arch-Q-field} is the field-jump generator: when
$\theta$ jumps to $\theta'$, \emph{every} residue is simultaneously
resampled from the new archetype stationary
$\eqm^{(k(c_n,\theta'))}$. This joint resampling is what preserves the
joint stationary $\rho[\theta]\prod_n \eqm^{(k(c_n,\theta))}[x_n]$ and
enforces detailed balance of the compound chain
(cf.~\citealp{tuffleysteel1998} on carrying vs.\ resampling under
covarion dynamics). Equation~\eqref{eq:arch-Q-res} is the standard GTR
residue substitution at fixed field: rate $\exch[x_n, y]\,
\eqm^{(k(c_n,\theta))}[y]$ with the LG08 exchangeabilities $\exch$
carried unchanged from Section~\ref{sec:tkfdp-model}.

\paragraph{Reversibility.}
Both marginals are reversible: $\rho$ is stationary for the F81-on-DP
field chain (Section~\ref{sec:dynamic-suppl}); $\eqm^{(k)}$ is
stationary for GTR with $\exch$ symmetric. The joint
$\rho[\theta]\prod_n \eqm^{(k(c_n,\theta))}[x_n]$ satisfies detailed
balance for the compound generator by direct verification: the
field-jump rate $\rho_{\text{chain}}\rho[\theta']$ times the pre-jump
weight $\rho[\theta]\prod_n \eqm^{(k(c_n,\theta))}[x_n]$ equals the
reverse rate $\rho_{\text{chain}}\rho[\theta]$ times the post-jump
weight $\rho[\theta']\prod_n \eqm^{(k(c_n,\theta'))}[x'_n]$ only when
each per-site resampled distribution $\eqm^{(k(c_n,\theta'))}[x'_n]$
matches the post-jump stationary --- which is exactly the resampling
rule~\eqref{eq:arch-Q-field}. The residue-only
generator~\eqref{eq:arch-Q-res} is reversible w.r.t.\
$\eqm^{(k(c_n,\theta))}$ site-by-site, so is reversible w.r.t.\ the
compound stationary as well.

\paragraph{Rooting via the pulley principle.}
The compound chain is reversible, so Felsenstein's pulley principle
lets us re-root a cherry anywhere without changing the likelihood.
Rooting at leaf $X$ replaces the cherry $\big((X, X_1,\dots,X_N),
(\text{LCA}), (Y, Y_1,\dots,Y_N)\big)$ with a single branch of length
$t$ from the observed $X$-endpoint to the observed $Y$-endpoint.
The $\theta_X$ at the root is a latent draw from $\rho$; site
residues at the root are the observed $X_n$. This eliminates the LCA
state marginalisation used in the cherry-doublet form of
Section~\ref{sec:dynamic-suppl} (the four-case decomposition indexed
by no-jump vs.\ ${\geq}1$-jump on each half-edge) and reduces it to a
two-case (no-jump vs.\ ${\geq}1$-jump) decomposition on the whole
branch, refined below to the three-case $\{0, 1, {\geq}2\}$
decomposition needed to isolate stationary middle segments.

\paragraph{Archetype orbits: a permutation-structured field selector.}
\label{par:archetype-orbits}
The free assignment $\arch(c,\theta)$ of \eqref{eq:arch-mapping} is more
expressive than the covariation it is meant to capture: an unstructured
$(c,\theta)\!\to\!k$ table lets a class drift to an arbitrary archetype at
each field value, so a ``flip'' need not be a \emph{charge-complementary
swap}, and empirically the learned table shows no acid$\leftrightarrow$base
structure above a label-permutation null. We therefore restrict the field
selector to act by \emph{permutation of archetypes within orbits}, which
makes ``flip $=$ swap'' hold by construction and replaces the
$K_c\times L_{\max}$ assignment table by a single partition of the archetype
vocabulary. Take the frozen diagonal base $K_c=K_a=K$, $\arch_0(c)=c$
(class $\equiv$ archetype), and draw a partition
$\mathcal{B}=\{B_1,\dots,B_m\}$ of $\{1,\dots,K\}$ into \emph{archetype
orbits} from a Dirichlet process with concentration $\alpha_{\mathrm{orb}}$;
large $\alpha_{\mathrm{orb}}$ keeps $\mathcal{B}$ near all-singletons, with
occasional small blocks.

\paragraph{Permutation field.} The field alphabet is the product of
symmetric groups on the orbits,
\begin{equation}
  \Theta \;=\; \prod_{b=1}^{m}\mathrm{Sym}(B_b),
  \qquad |\Theta| \;=\; \prod_{b=1}^{m}|B_b|!\,,
  \label{eq:orbit-field-space}
\end{equation}
so a field state is a tuple $\theta=(\sigma_1,\dots,\sigma_m)$ of per-orbit
permutations, and the archetype seen by class $c$ is its base archetype
permuted within its own orbit,
\begin{equation}
  k(c,\theta) \;=\; \sigma_{b(c)}(c),
  \qquad b(c) = \{\,b : c \in B_b\,\}.
  \label{eq:orbit-action}
\end{equation}
This is the constrained form of \eqref{eq:arch-mapping}: as $\theta$ ranges
over $\Theta$, class $c$ visits only the archetypes of its orbit, as a group
action shared across classes. A singleton orbit contributes the trivial
group ($|B_b|!=1$): its classes are \emph{specialisers} with no reachable
flip. A size-2 orbit $B_b=\{c,c'\}$ contributes
$\mathrm{Sym}=\{e,(c\,c')\}$, whose two states leave the pair fixed or
transpose it.

\paragraph{Field dynamics.} Each orbit $B_b$ carries its \emph{own} F81-on-DP
field chain \eqref{eq:arch-Q-field} on $\mathrm{Sym}(B_b)$, evolving
\emph{independently} (asynchronously) of the other orbits: at rate
$\rho_{\text{chain}}\,r$ its permutation resamples from
$\mathrm{Unif}(\mathrm{Sym}(B_b))$, on its own jump times. The \emph{only}
coupling is therefore between columns that \emph{share} an orbit --- a
same-orbit pair is a $2$-cycle ($\{AB,BA\}$, or $\{AA,BB\}$ if both columns
sit at the same pole), the compensatory flip. Columns in \emph{different}
orbits are conditionally independent: the $\{AC,AD,BC,BD\}$ configuration
collapses to two unsynchronised singletons (one oscillating $A\!\leftrightarrow
\!B$, the other $C\!\leftrightarrow\!D$). This keeps a cluster's field
\emph{local to its orbits} --- independent of how the rest of the archetype
vocabulary is partitioned. The single-orbit F81 chain and its reversibility
are exactly \eqref{eq:arch-Q-field}--\eqref{eq:arch-Q-res}; a per-cluster rate
$r$ with an invariant bin $r=0$ gives the $+\Gamma+\mathrm{I}$ field-rate
(singleton-orbit or invariant-bin cluster $=$ static).

\paragraph{Per-cluster locality and the state-space cap.} The cluster
likelihood \emph{factorises} over the distinct orbits its columns touch:
\[
  \log p(\text{data}_C) \;=\; \sum_{b\,\in\,\text{touched}(C)}
    \log p\!\big(\text{cols}_b \,\big|\, \sigma_b\big),
\]
each term a Felsenstein forward over $\mathrm{Sym}(B_b)$ alone. So every forward
runs on $L=|B_b|!\le s_{\max}!=2$ states --- the \emph{maximum orbit size},
never a product --- and neither $|\Theta_C|$ nor the global
$|\Theta|=\prod_b|B_b|!$ enters. For cap-2 clusters a forward is a coupled pair
(two columns sharing an orbit, $L\le2$) or a singleton/covarion single
($L\le2$); a different-orbit pair is two independent single forwards. We cap
$s_{\max}$ (e.g.\ $s_{\max}=2$, transpositions only).

\paragraph{Compensatory flip and its posterior.} A cap-2 cluster
$C=\{i,j\}$ exhibits a compensatory flip iff its columns' classes share a
pair-orbit, $c_i,c_j\in B_b$ with $|B_b|=2$: the transposition
$\sigma_b=(c_i\,c_j)$ swaps their archetypes,
$\eqm^{(k(c_i,\theta))}\leftrightarrow\eqm^{(k(c_j,\theta))}$ --- a
charge-complementary swap when $B_b$ pairs a positive- and a
negative-charge archetype. If the two classes are singletons or lie in
distinct orbits, no field state couples them and $C$ is static. The flip
posterior is then a structural object,
\begin{equation}
  \varphi(C) \;=\; \Pr\!\bigl[\,\sigma_b \ \text{non-identity on some tree
  edge} \ \big|\ \text{data}\,\bigr],
  \label{eq:orbit-flip-posterior}
\end{equation}
identically zero for singleton-orbit clusters rather than an emergent
near-zero --- removing the ``flip vs.\ chance'' ambiguity of the free table.

\paragraph{Inference: split--merge on orbits.} The learned field structure
is now the partition $\mathcal{B}$ of the $K_a$ archetypes, sampled by
Jain--Neal split--merge moves --- as in \autoref{par:arch-split-merge}, but
on archetype \emph{orbits} rather than vocabulary slots: a merge places two
archetypes in a common orbit (creating a candidate flip), a split separates
them. This replaces the $O(K_c L_{\max})$ resampling of the free
$\arch(c,\theta)$ table --- step (3) of the SVI loop below --- with moves on
a partition of $K_a$ elements; the site-class Gibbs ($c_s$) and the
$+\Gamma+\mathrm{I}$ rate updates are unchanged. Acceptance uses the same
class-marginal / field-rate-marginal cluster evidence as the rest of the
loop, summed over the small $\Theta_C$.

\begin{remark}[Exchangeability reading]
Acting on the archetype vocabulary by the subgroup of $S_{K_a}$ that
preserves $\mathcal{B}$ is an exchangeability symmetry: the model is
invariant to relabelling archetypes \emph{within} an orbit, and the field
selects the coset. This is the finite-group analogue of the de-Finetti
coupling of the companion paper, giving the field selector a
representation-theoretic rather than tabular description.
\end{remark}

\paragraph{Inference: soft-EM archetype update with sampled discrete
latents.}
The archetype variant is the sole learned inference target in the
dynamic-latent-field SVI loop (Section~\ref{sec:tkfdp-svi}). Per
outer iter:
\begin{enumerate}
  \item CRP-Gibbs on cluster partitions $\pi$ (unchanged).
  \item Class-Gibbs on per-column labels $c_s$ (unchanged).
  \item Sample $\arch(c,\theta) \sim \mathrm{Cat}\bigl(\rho_a[k]\cdot
    \prod_a \eqm^{(k)}[a]^{N[c,\theta,a]}\bigr)$ from the archetype
    posterior given the current per-$(c,\theta)$ substitution counts
    $N[c,\theta,\cdot]$.
  \item Accumulate Holmes--Rubin sufficient statistics
    (Section~\ref{par:arch-hr}) by soft attribution over the field
    trajectory, given the currently-sampled discrete latents.
  \item M-steps: closed-form Gamma-posterior update on
    $\rho_{\text{chain}}$; GTR fixed-point update on $\eqm^{(k)}$;
    TSB Beta-posterior update on $\rho_a$.
\end{enumerate}
The pattern matches the recommendation of
\citet[Section~4]{bealghahramani2002} for hybrid Gibbs / variational
schemes with discrete-continuous latent decompositions.

\paragraph{Field-chain spectrum: rank-1 collapse.}
The F81-on-DP field-chain generator
$Q_\theta = \rho_{\text{chain}}(\mathbf{1}\rho^\top - I)$
has two distinct eigenvalues:
$\lambda_0 = 0$ (right eigenvector $\mathbf{1}$, left eigenvector
$\rho$) and $\lambda_1 = -\rho_{\text{chain}}$ (multiplicity
$L_{\max}-1$, eigenspace $\{v : \rho^\top v = 0\}$). Spectral
projectors: $P_0 = \mathbf{1}\rho^\top$ (rank $1$, elementwise
$P_0[i,j] = \rho[j]$) and $P_1 = I - P_0$. This gives the closed-form
transition probability
\begin{equation}
  P_\theta(i \to j; \tau) \;=\; \rho[j] \;+\; g_\tau\,(\delta_{ij} - \rho[j])
  \;=\; g_\tau\,\delta_{ij} + (1 - g_\tau)\,\rho[j],
  \label{eq:arch-Ptheta}
\end{equation}
with $g_\tau := e^{-\rho_{\text{chain}}\tau}$. The state-dependent
``no jump on interval $\tau$'' survival probability is a distinct
quantity:
\begin{equation}
  \beta_\tau(i) \;:=\; e^{-\rho_{\text{chain}}(1-\rho[i])\tau}
  \;=\; g_\tau^{\,1-\rho[i]},
  \label{eq:arch-beta}
\end{equation}
smaller than $P_\theta(i \to i; \tau)$ because the latter also credits
trajectories that jump away and return.

\paragraph{Number-of-jumps decomposition.}
With state-dependent exit rate $r'(\theta) = \rho_{\text{chain}}(1-\rho[\theta])$
for the exit from field state $\theta$, decompose the branch $[0, t]$
by the number of real field jumps $M \in \{0, 1, 2+\}$:
\begin{align}
  P(M = 0 \mid \theta_X, t) &= \beta_t(\theta_X),
  \label{eq:arch-P0}
  \\
  P(M = 1, \theta_1 = j, \tau_1 \mid \theta_X, t)\,d\tau_1
    &= \rho_{\text{chain}}\,\rho[j]\,
       \beta_{\tau_1}(\theta_X)\,\beta_{t-\tau_1}(j)\,d\tau_1,
    \quad j \ne \theta_X,
  \label{eq:arch-P1-joint}
  \\
  P(M \geq 2 \mid \theta_X, t)
    &= 1 - P(M=0\mid \theta_X, t) - P(M=1\mid \theta_X, t).
  \label{eq:arch-P2plus}
\end{align}
Marginalising~\eqref{eq:arch-P1-joint} over $\theta_1$ and $\tau_1$
yields $P(M=1 \mid \theta_X, t)$ as a sum over $j \ne \theta_X$ of
\begin{equation}
  \rho[j] \cdot \int_0^t \rho_{\text{chain}}\,\beta_{\tau_1}(\theta_X)\,\beta_{t-\tau_1}(j)\,d\tau_1
  \;=\;
  \rho[j] \cdot
  \begin{cases}
    \dfrac{\beta_t(j) - \beta_t(\theta_X)}{\rho[j] - \rho[\theta_X]},
      & \rho[j] \ne \rho[\theta_X], \\[6pt]
    \rho_{\text{chain}} t\,\beta_t(\theta_X),
      & \rho[j] = \rho[\theta_X].
  \end{cases}
  \label{eq:arch-P1-marg}
\end{equation}
The $\rho[j] = \rho[\theta_X]$ degenerate case is the L'H\^{o}pital
limit of the generic form and appears when two field states have
identical stick-breaking weight.

\paragraph{Sufficient statistics per site.}
\label{par:arch-hr}
For each (cluster, cherry $q$ of length $t_q$, site $n$) we accumulate
the following expectations, weighted by the posterior over
$(\theta_X, \theta_Y, M)$ implied by the site-level cluster likelihood
factorisation:
\begin{align}
  V[k, a]     &:= \text{expected \# of arch-}k\text{ segments starting at residue }a,
  \label{eq:arch-V} \\
  U[k, i, j]  &:= \text{expected \# of }i\to j\text{ substitutions within arch-}k
                  \text{ intervals},
  \label{eq:arch-U} \\
  W[k, i]     &:= \text{expected dwell time in state }i\text{ within arch-}k
                  \text{ intervals},
  \label{eq:arch-W} \\
  N_\theta[q] &:= \text{expected \# of field jumps on cherry }q,
  \label{eq:arch-Ntheta} \\
  T_q         &:= t_q\ \text{(deterministic branch length).}
  \label{eq:arch-Tq}
\end{align}
These aggregate over the corpus by direct summation. Per-case
contributions:

\emph{Case $M = 0$}: $k_X = \arch(c_n, \theta_X)$.
\begin{align}
  V[k_X, X_n]  &\mathrel{+}= P(M=0 \mid \theta_X) \cdot \rho[\theta_X], \\
  W[k_X, i]    &\mathrel{+}= P(M=0 \mid \theta_X) \cdot \rho[\theta_X] \cdot
                             W^{\mathrm{br}}_{k_X}(X_n \to Y_n; i; t), \\
  U[k_X, i, j] &\mathrel{+}= P(M=0 \mid \theta_X) \cdot \rho[\theta_X] \cdot
                             U^{\mathrm{br}}_{k_X, ij}(X_n \to Y_n; t),
\end{align}
where $W^{\mathrm{br}}_k(X \to Y; i; t)$ is the standard GTR bridge
expected dwell in state $i$ given endpoints $X, Y$ and branch length
$t$ under stationary $\eqm^{(k)}$, computed by the standard HR
spectral formula \citep{HolmesRubin2002}. The symmetrisation
$D_k^{1/2} Q_k D_k^{-1/2}$ (with $D_k = \diag(\eqm^{(k)})$; write
$d^k_i := \sqrt{\eqm^{(k)}[i]}$ and let $\xi^k_\alpha, U^k$ be the
eigenvalues/eigenvectors of the symmetrised generator) yields a
real-symmetric matrix whose eigendecomposition serves both the
expected-dwell integral
\[
  W_i^{\mathrm{br}}(X \to Y; t)
  \;=\; \frac{d^k_Y}{d^k_X\, P^k_{XY}(t)}
    \sum_{\alpha, \beta} J^{\alpha\beta}(t)\,
      U^k_{X\alpha}\, U^k_{i\alpha}\, U^k_{i\beta}\, U^k_{Y\beta}
\]
and the expected-transition integral
\begin{equation}
  U_{ij}^{\mathrm{br}}(X \to Y; t)
  \;=\; \frac{\exch[i,j]\, d^k_i\, d^k_j\, d^k_Y}{d^k_X\, P^k_{XY}(t)}
    \sum_{\alpha, \beta} J^{\alpha\beta}(t)\,
      U^k_{X\alpha}\, U^k_{i\alpha}\, U^k_{j\beta}\, U^k_{Y\beta},
  \qquad i \ne j,
  \label{eq:arch-Ubr-spectral}
\end{equation}
by the well-known $J^{\alpha\beta}$ formula (with
$J^{\alpha\beta}(t) = (e^{\xi^k_\alpha t} - e^{\xi^k_\beta t}) /
(\xi^k_\alpha - \xi^k_\beta)$ for distinct eigenvalues and $t\,
e^{\xi^k_\alpha t}$ in the degenerate limit). Note that
$U_{ij}^{\mathrm{br}} \ne Q^k_{ij}\, W_i^{\mathrm{br}}$: the
identity $E[N_{ij}] = Q_{ij}\, E[W_i]$ holds for the
\emph{unconditional} CTMC (Dynkin), but the Doob $h$-transform for
the bridge makes the generator time-inhomogeneous, so the sum-and-
integrate in \eqref{eq:arch-Ubr-spectral} must be carried out with
distinct $\alpha$ (before the jump) and $\beta$ (after the jump)
eigenvalue indices rather than collapsing on the site index.

\emph{Case $M \geq 1$}: the trajectory decomposes into a first
segment $[0, \tau_1]$ under $k_X = \arch(c_n, \theta_X)$, a
(possibly empty) sequence of middle segments under
$\arch(c_n, \theta_m)$ for intermediate field states $\theta_m$, and
a last segment $[\tau_M, t]$ under $k_Y = \arch(c_n, \theta_Y)$.
The three segment kinds have qualitatively different HR structure and
are handled separately below. Contributions integrate over $\tau_1$,
$\theta_1$ (and, for $M \geq 2$, over the intermediate segment
counts) using~\eqref{eq:arch-P1-joint} for the M{=}1 slice and the
residual construction described below for the $M \geq 2$ slice; the
M{=}1 and $M \geq 2$ contributions are combined into a unified
$M \geq 1$ attribution because the first/last-segment formulas take
the same form (with $M \geq 2$ contributing additional middle-segment
terms only).

\emph{First segment (start $X_n$ observed, end resampled from
$\eqm^{(k_1)}$ at $\tau_1^+$).} Because the residue at $\tau_1^+$ is
resampled independently under Interp~1, the first segment is an
\emph{unconditional} CTMC of duration $\tau_1$ starting at $X_n$ ---
its end residue is not conditioned. Dynkin's identity therefore holds:
$E[N_{ij}] = Q^{k_X}_{ij}\, E[T_i]$, and specifically the Doob
$h$-transform correction used in the M{=}0 bridge formula
\eqref{eq:arch-Ubr-spectral} is \emph{not} needed.

\emph{Exact $M \geq 1$ density for $\tau_1$.} We compute the marginal
density of $\tau_1$ conditional on $M \geq 1$, $\theta_X$, $\theta_Y$
in closed form:
$p(\tau_1) \propto r'_X e^{-r'_X \tau_1} \sum_{\theta_1 \ne \theta_X}
[\rho(\theta_1) / (1 - \rho(\theta_X))]\, P_\theta(\theta_1 \to \theta_Y;
t - \tau_1)$.
Using stationarity of $\rho$ under the field CTMC
($\sum_{\theta_1} \rho(\theta_1) P_\theta(\theta_1 \to \theta_Y; s) =
\rho(\theta_Y)$) and the F81-on-DP form $P_\theta(u \to v; s) =
g_s\, \delta_{u,v} + (1 - g_s)\, \rho(v)$ with $g_s = e^{-\rho_{\text{chain}} s}$,
the density collapses to a \emph{two-exponential mixture}:
\begin{equation}
  p(\tau_1 \mid M \geq 1) \;\propto\;
    \beta_1\, e^{\delta_1 \tau_1} + \beta_2\, e^{\delta_2 \tau_1}
  \label{eq:arch-tau1-exact}
\end{equation}
with
\begin{align}
  \delta_1 &= -r'_X, & \delta_2 &= \rho_{\text{chain}} \rho(\theta_X), \\
  \beta_1 &= r'_X\, \rho(\theta_Y), &
  \beta_2 &= r'_X\, \tfrac{\rho(\theta_X)}{1 - \rho(\theta_X)}\,
              \bigl[\rho(\theta_Y) - \delta_{\theta_X, \theta_Y}\bigr]\, g_t.
\end{align}
(For $\theta_X \neq \theta_Y$ both coefficients are positive; for the
$\theta_X = \theta_Y$ self-return case $\beta_2 < 0$, reflecting that
$M = 1$ is impossible and the density is a difference of two
exponentials.) The last-segment density on $t - \tau_M$ is the same
mixture with $\theta_X$ and $\theta_Y$ exchanged (time reversal of the
field bridge).

\emph{Elementary truncated-exp integrals on $[0, t]$.}
For any $\delta \in \mathbb{R}$ and $t > 0$, define the truncated-exp
normalizer, mean, dwell-moment, and transition-average primitives:
\begin{align}
  Z(\delta, t)
  &\;:=\; \int_0^t e^{\delta s}\, ds
  \;=\;
    \begin{cases}
      \dfrac{e^{\delta t} - 1}{\delta}, & \delta \ne 0,\\
      t, & \delta = 0,
    \end{cases}
  \label{eq:arch-Z} \\[4pt]
  \bar\tau(\delta, t)
  &\;:=\; \frac{1}{Z(\delta, t)} \int_0^t \tau\, e^{\delta \tau}\, d\tau
  \;=\;
    \begin{cases}
      \dfrac{t\, e^{\delta t}}{e^{\delta t} - 1} - \dfrac{1}{\delta}, & \delta \ne 0, \\[4pt]
      t/2, & \delta = 0,
    \end{cases}
  \label{eq:arch-Etau} \\[4pt]
  H(\xi; \delta, t)
  &\;:=\; \frac{1}{Z(\delta, t)} \int_0^t \frac{e^{\xi \tau} - 1}{\xi}\,
                                             e^{\delta \tau}\, d\tau \notag \\
  &\;=\;
    \begin{cases}
      \dfrac{Z(\xi + \delta, t) - Z(\delta, t)}{\xi\, Z(\delta, t)},
        & \xi \ne 0,\ \xi + \delta \ne 0, \\[6pt]
      \dfrac{1}{\delta} - \dfrac{t}{e^{\delta t} - 1},
        & \xi + \delta = 0,\ \delta \ne 0, \\[4pt]
      \bar\tau(\delta, t),
        & \xi = 0,
    \end{cases}
  \label{eq:arch-H} \\[4pt]
  T(\xi; \delta, t)
  &\;:=\; \frac{1}{Z(\delta, t)} \int_0^t e^{\xi \tau}\, e^{\delta \tau}\, d\tau
  \;=\; \frac{Z(\xi + \delta, t)}{Z(\delta, t)}, \quad
  T(-\delta; \delta, t) = \frac{t}{Z(\delta, t)}.
  \label{eq:arch-T}
\end{align}
$H$ and $T$ are the two spectral-index primitives appearing in the
free-end dwell and last-jump target-distribution formulas below.

\emph{Combination under the two-exponential mixture.}
For $p(\tau) \propto \beta_1 e^{\delta_1 \tau} + \beta_2 e^{\delta_2 \tau}$
on $[0, t]$, the mixture normalizer, weights, and expectation of any
linear functional $f$ of $\tau$ are
\begin{align}
  Z_{\mathrm{mix}}
  &\;=\; \beta_1\, Z(\delta_1, t) + \beta_2\, Z(\delta_2, t),
  \qquad
  a_i \;=\; \frac{\beta_i\, Z(\delta_i, t)}{Z_{\mathrm{mix}}},
  \qquad a_1 + a_2 = 1,
  \label{eq:arch-mix-weights} \\[4pt]
  E_{\mathrm{mix}}[f(\tau)]
  &\;=\; a_1\, E_{p_{\delta_1}}[f(\tau)] + a_2\, E_{p_{\delta_2}}[f(\tau)].
  \label{eq:arch-mix-linearity}
\end{align}
This linearity lets us evaluate every SS attribution by calling each
single-$\delta$ primitive twice --- once at $\delta_1$, once at
$\delta_2$ --- and taking the $a_1, a_2$-weighted sum. No mixture-
specific closed forms are needed. The reference implementation
(\texttt{gtr\_free\_end\_hr\_avg\_case1\_jax},
\texttt{gtr\_transition\_avg\_case1\_jax}) is machine-precision
verified against the joint (field $\times$ residue) matrix-exp
sufficient statistics
(\texttt{tests/dynfield/test\_hr\_full\_ratematrix.py}, relative
error $\lesssim 10^{-14}$ across the full parameter grid).

\emph{First-segment $W, U$ in terms of $H$.}
The free-end dwell (start at $X_n$, $\tau_1$ drawn from case-1 density with
rate $\delta$) is
\begin{equation}
  W_i^{\mathrm{free}}(X_n; t, \delta)
  \;=\; \frac{d^{k_X}_i}{d^{k_X}_{X_n}}
    \sum_\alpha U^{k_X}_{X_n \alpha}\, U^{k_X}_{i\alpha}\,
      H\!\bigl(\xi^{k_X}_\alpha;\, \delta,\, t\bigr),
  \label{eq:arch-Wfree}
\end{equation}
and the free-end U (Dynkin identity holds because the endpoint is
resampled independently under Interp 1) is
\begin{equation}
  U^{\mathrm{free}}_{ij}(X_n; t, \delta) \;=\;
    Q^{k_X}_{ij}\, W_i^{\mathrm{free}}(X_n; t, \delta),
  \qquad i \ne j.
  \label{eq:arch-Ufree}
\end{equation}
Under the exact $M \geq 1$ density (Eq.~\ref{eq:arch-tau1-exact}), the
first-segment contributions to $\{W, U\}$ at $(c_n, \theta_X, \cdot)$
are the mixture-averaged versions:
\begin{align}
  W^{\text{first}}_i
  &\;=\; a_1^{\text{first}}\, W_i^{\mathrm{free}}(X_n; t, \delta_1^{\text{first}})
    + a_2^{\text{first}}\, W_i^{\mathrm{free}}(X_n; t, \delta_2^{\text{first}}),
  \label{eq:arch-Wfirst-mix} \\
  U^{\text{first}}_{ij}
  &\;=\; a_1^{\text{first}}\, U^{\mathrm{free}}_{ij}(X_n; t, \delta_1^{\text{first}})
    + a_2^{\text{first}}\, U^{\mathrm{free}}_{ij}(X_n; t, \delta_2^{\text{first}})
  \notag \\
  &\;=\; Q^{k_X}_{ij}\, W^{\text{first}}_i,
  \label{eq:arch-Ufirst-mix}
\end{align}
with $(\delta_i^{\text{first}}, \beta_i^{\text{first}})$ from
Eq.~\ref{eq:arch-tau1-exact} and $a_i^{\text{first}}$ from
Eq.~\ref{eq:arch-mix-weights}.

\emph{Last-jump target-distribution via $T$.}
The V attribution at $(c_n, \theta_Y, \cdot)$ receives the expected
residue distribution at the moment of the last field jump. Under
Interp 1 with a resample from $\eqm^{(k_Y)}$, the observed $Y_n$
at time $t$ gives (by reversibility) a bridge-conditional expected
resample distribution which by time-reversal equals
$E_s[P^{k_Y}(Y_n \to a; s)]$ averaged over $s = t - \tau_M$:
\begin{equation}
  P^{\mathrm{res}, k}_a(X; t, \delta) \;:=\; E_s[P^{k}(X \to a; s)]
  \;=\; \frac{d^k_a}{d^k_X}
    \sum_\alpha U^k_{X \alpha}\, U^k_{a\alpha}\,
      T\!\bigl(\xi^k_\alpha;\, \delta,\, t\bigr).
  \label{eq:arch-Pres}
\end{equation}
Under the exact $M \geq 1$ last-segment density (mixture with
$\theta_X, \theta_Y$ exchanged in Eq.~\ref{eq:arch-tau1-exact}), the V
attribution is
\begin{equation}
  V^{\text{last}}_a
  \;=\; a_1^{\text{last}}\, P^{\mathrm{res}, k_Y}_a(Y_n; t, \delta_1^{\text{last}})
    + a_2^{\text{last}}\, P^{\mathrm{res}, k_Y}_a(Y_n; t, \delta_2^{\text{last}}),
  \label{eq:arch-Vlast-mix}
\end{equation}
with $\delta_1^{\text{last}} = -r'_Y$, $\delta_2^{\text{last}}
= \rho_{\text{chain}} \rho(\theta_Y)$, $\beta_1^{\text{last}} = r'_Y \rho(\theta_X)$,
$\beta_2^{\text{last}} = r'_Y (\rho(\theta_Y)/(1 - \rho(\theta_Y)))
(\rho(\theta_X) - \delta_{\theta_X, \theta_Y}) g_t$, and
$a_i^{\text{last}}$ from Eq.~\ref{eq:arch-mix-weights}.

\emph{Middle segments (both endpoints at stationarity), $M \geq 2$
only.} When $M \geq 2$, each middle segment has both endpoints
i.i.d.\ from $\eqm^{(k)}$ (the entering jump resamples the residue
from stationarity; the exiting jump does the same on the next
segment). The segment is a stationary CTMC of expected duration
$\tau_{\mathrm{mid}}$, and Dynkin on the stationary chain gives
$E[T_i] = \tau_{\mathrm{mid}} \eqm^{(k)}[i]$ and
$E[N_{ij}] = \tau_{\mathrm{mid}} \eqm^{(k)}[i]\, \exch[i, j]\,
\eqm^{(k)}[j]$ (using $Q^{k}_{ij} = \exch[i, j]\, \eqm^{(k)}[j]$).
The per-site middle-segment contributions to $(V, W, U)$ are
therefore
\begin{align}
  V[k, a]     &\mathrel{+}= E_{M\geq 1}[\#\ \text{middle arrivals at arch }k]\, \eqm^{(k)}[a],
  \label{eq:arch-mid-V} \\
  W[k, a]     &\mathrel{+}= E_{M\geq 1}[\text{middle time in arch }k]\, \eqm^{(k)}[a],
  \label{eq:arch-mid-W} \\
  U[k, i, j]  &\mathrel{+}= E_{M\geq 1}[\text{middle time in arch }k]\, \eqm^{(k)}[i]\,
                            \exch[i, j]\,\eqm^{(k)}[j].
  \label{eq:arch-mid-U}
\end{align}
The middle-arrival count and middle-time expectations are the
residuals of the whole-branch bridge-conditional quantities after
subtracting the first- and last-segment contributions. Under Interp 1
the first segment sits entirely at $\theta_X$ and the last entirely at
$\theta_Y$; the expected first- and last-segment durations under the
exact $M \geq 1$ mixture densities are
\begin{align}
  \overline{\tau_1}
  &\;=\; a_1^{\text{first}}\, \bar\tau(\delta_1^{\text{first}}, t)
        + a_2^{\text{first}}\, \bar\tau(\delta_2^{\text{first}}, t),
  \label{eq:arch-Etau1-mix} \\
  \overline{t - \tau_M}
  &\;=\; a_1^{\text{last}}\, \bar\tau(\delta_1^{\text{last}}, t)
        + a_2^{\text{last}}\, \bar\tau(\delta_2^{\text{last}}, t).
  \label{eq:arch-EtauM-mix}
\end{align}
(Both use Eq.~\ref{eq:arch-Etau} at the corresponding $\delta$ and the
mixture weights from Eq.~\ref{eq:arch-mix-weights}.) The middle-time
at intermediate field state $\theta_{\text{int}}$ is then
\begin{equation}
  \tau_{\text{mid}}(\theta_{\text{int}})
  \;=\;
  \max\!\left(0,\
    E[T_{\theta_{\text{int}}} \mid \theta_X, \theta_Y, t, M \geq 1]
    - \overline{\tau_1}\, \mathbb{1}[\theta_{\text{int}} = \theta_X]
    - \overline{t - \tau_M}\, \mathbb{1}[\theta_{\text{int}} = \theta_Y]
  \right),
  \label{eq:arch-taumid}
\end{equation}
and the middle-arrival count (excluding the last-arrival at $\theta_Y$
which initiates the last segment, not a middle segment) is
\begin{equation}
  N_{\text{mid}}(\theta_{\text{int}})
  \;=\;
  \max\!\left(0,\
    E[N_{\theta_{\text{int}}} \mid \theta_X, \theta_Y, t, M \geq 1]
    - \mathbb{1}[\theta_{\text{int}} = \theta_Y]
  \right).
  \label{eq:arch-Nmid}
\end{equation}
where $E[T_\theta \mid \cdot]$ and $E[N_\theta \mid \cdot]$ are the
whole-branch bridge-conditional expectations from
Eq.~\ref{eq:arch-Ntheta-closed} and
Eq.~\ref{eq:arch-arrivals-closed} below. The middle-segment SS
attribution~\eqref{eq:arch-mid-V}--\eqref{eq:arch-mid-U} then uses
$\tau_{\text{mid}}$ for the ``middle time in arch $k$'' factor and
$N_{\text{mid}}$ for the ``middle arrivals at arch $k$'' factor. At
$M = 1$ the residuals~\eqref{eq:arch-taumid}--\eqref{eq:arch-Nmid}
vanish identically ($\overline{\tau_1} + \overline{t - \tau_M} = t$
when $M = 1$; the last-arrival subtraction cancels the single arrival
at $\theta_Y$), so a single $M \geq 1$ attribution correctly covers
both the $M = 1$ and $M \geq 2$ slices.

\emph{Last segment (marginal start from $\eqm^{(k_Y)}$, end $Y_n$
observed).} The last segment is a ``reversed bridge'': the forward
start $z$ is drawn from $\eqm^{(k_Y)}$ (the last jump's resample)
and the forward end is observed as $Y_n$. Under GTR reversibility
$\eqm^{(k_Y)}[z]\, P^{k_Y}(z \to Y_n; \tau) = \eqm^{(k_Y)}[Y_n]\,
P^{k_Y}(Y_n \to z; \tau)$, so integrating over $z \sim
\eqm^{(k_Y)}$ the conditional dwell in state $i$ collapses to a
free-end from $Y_n$ in \emph{reversed} time:
\begin{equation}
  W_i^{\mathrm{last}}(Y_n; \tau)
  \;=\; E_{z \sim \eqm^{(k_Y)}}\bigl[
    W^{\mathrm{br}}_i(z \to Y_n; \tau)\bigr]
  \;=\; \int_0^\tau P^{k_Y}(Y_n \to i; s)\, ds
  \;=\; W_i^{\mathrm{free}}(Y_n; \tau).
  \label{eq:arch-last-W}
\end{equation}
For expected transitions the same reversibility argument, applied to
the Doob-$h$-transformed bridge integrand, gives
\begin{equation}
  U^{\mathrm{last}}_{ij}(Y_n; \tau)
    \;=\; \eqm^{(k_Y)}[i]\, \exch[i, j]\, W_j^{\mathrm{free}}(Y_n; \tau),
    \qquad i \ne j,
  \label{eq:arch-last-U}
\end{equation}
which is \emph{not} the Dynkin form $Q^{k_Y}_{ij}\, W_i^{\mathrm{last}}
= \exch[i, j]\, \eqm^{(k_Y)}[j]\, W_i^{\mathrm{free}}(Y_n; \tau)$
(differences: the $i \leftrightarrow j$ swap in both the stationary
factor and the free-end-dwell argument). This is the same phenomenon
as in the M{=}0 bridge --- endpoint conditioning breaks the
unconditional-CTMC Dynkin identity --- but here the closed form is
simpler because one endpoint is stationary rather than observed, and
one 4-index spectral summation collapses to a 2-index (single-$\alpha$)
one. Under the exact $M \geq 1$ last-segment density
(Eq.~\ref{eq:arch-tau1-exact} with $\theta_X, \theta_Y$ exchanged), the
mixture-averaged last-segment attributions to $(W, U)$ at
$(c_n, \theta_Y, \cdot)$ are
\begin{align}
  W^{\text{last}}_j
  &\;=\; a_1^{\text{last}}\, W_j^{\mathrm{free}}(Y_n; t, \delta_1^{\text{last}})
    + a_2^{\text{last}}\, W_j^{\mathrm{free}}(Y_n; t, \delta_2^{\text{last}}),
  \label{eq:arch-Wlast-mix} \\
  U^{\text{last}}_{ij}
  &\;=\; \eqm^{(k_Y)}[i]\, \exch[i, j]\, W^{\text{last}}_j,
  \qquad i \ne j.
  \label{eq:arch-Ulast-mix}
\end{align}
(Note the $i \leftrightarrow j$ swap between $W^{\text{last}}$
indexing and Dynkin's would-be $Q^{k_Y}_{ij} W^{\text{last}}_i$
formula, forced by the reversibility argument.) The V attribution at
$(c_n, \theta_Y, \cdot)$ is $V^{\text{last}}_a$ from
Eq.~\ref{eq:arch-Vlast-mix}.

The M{=}1 attribution follows immediately from
Eqs.~\ref{eq:arch-Wfree}, \ref{eq:arch-Ufree}, \ref{eq:arch-last-W},
\ref{eq:arch-last-U} (first + last only; no middle contribution). For
$M \geq 2$, the first- and last-segment formulas are unchanged and
the middle-segment attribution
\eqref{eq:arch-mid-V}--\eqref{eq:arch-mid-U} accounts for the
intermediate field time. Together with the whole-branch bridge
$E[N_\theta \mid \theta_X, \theta_Y, t]$
(Eq.~\ref{eq:arch-Ntheta-closed}) and the per-destination arrivals
(Eq.~\ref{eq:arch-arrivals-closed}), these formulas express
$(V, U, W, N_\theta)$ as closed-form functions of
$\{\beta_t(i), g_t, t\}$ and the GTR bridge/free-end primitives,
without any integral remaining to be evaluated numerically.

\paragraph{Whole-branch bridge $E[N_\theta \mid \theta_X, \theta_Y, t]$.}
The total field-jump expectation entering the middle-segment residual
construction of the previous paragraph is the bridge-conditional
\begin{equation}
  E[N_\theta \mid \theta_X, \theta_Y, t] \;\cdot\; P_\theta(\theta_X \to \theta_Y; t)
    \;=\; \sum_{i \ne j} Q_\theta[i, j]\,
      \int_0^t P_\theta(\theta_X \to i; s)\,
              P_\theta(j \to \theta_Y; t{-}s)\, ds,
  \label{eq:arch-Ntheta-bridge}
\end{equation}
which is \emph{not} equal to
$\rho_{\text{chain}} \sum_i (1 - \rho[i])\, E_T[i \mid \theta_X, \theta_Y, t]$.
That shortcut multiplies the unconditional exit rate $-Q_\theta[i,i] =
\rho_{\text{chain}} (1 - \rho[i])$ against the bridge-conditional dwell,
and only recovers the correct expectation when the trajectory is at
field-stationarity throughout (endpoints i.i.d.\ from $\rho$, as inside a
case-$(2+)$ middle segment). Under the rank-1 spectral form
\eqref{eq:arch-Ptheta} substitution into~\eqref{eq:arch-Ntheta-bridge}
collapses to a closed form in $\{\rho, \rho_{\text{chain}}, t, g_t,
\sigma\}$ with $\sigma := \sum_i \rho[i]^2$:
\begin{equation}
  E[N_\theta \mid \theta_X, \theta_Y, t] \cdot P_\theta(\theta_X \to \theta_Y; t)
    = \Sigma_1(\theta_Y) + \Sigma_{23}(\theta_X, \theta_Y) + \Sigma_4(\theta_X, \theta_Y),
  \label{eq:arch-Ntheta-closed}
\end{equation}
with
\begin{align}
  \Sigma_1(b)    &= \rho_{\text{chain}}\,t\,\rho[b]\,(1 - \sigma), \\
  \Sigma_{23}(a, b) &= (1 - g_t)\, \rho[b]\, (2\sigma - \rho[a] - \rho[b]), \\
  \Sigma_4(a, b) &= -\rho_{\text{chain}}\, g_t\, t\,
    \bigl(\delta_{a, b}\rho[a] - \rho[a]\rho[b] - \rho[b]^2 + \rho[b]\sigma\bigr).
\end{align}
Deriving this: rewrite $P_\theta(u \to v; \tau) = \rho[v] + g_\tau
(\delta_{u, v} - \rho[v])$, integrate the four cross-products of the
$\{\rho, g_\tau\}$ terms against $Q_\theta[i, j] = \rho_{\text{chain}}
\rho[j]$ ($i \ne j$), and use
$\int_0^t 1\, ds = t$, $\int_0^t g_s\, ds = (1 - g_t)/\rho_{\text{chain}}$,
$\int_0^t g_s\, g_{t-s}\, ds = t\, g_t$. This is verified against
Gillespie in~\texttt{tests/dynfield/test\_hr\_stats\_closed\_form.py}
(test~\texttt{test\_field\_bridge\_EN\_matches\_gillespie}).

\paragraph{Per-destination arrival counts.}
The bridge-conditional expected count of field jumps that \emph{arrive}
at a specific destination $\theta_{\text{int}}$ is needed for the V
attribution (each arrival triggers a residue resample from
$\eqm^{(\arch(c_n, \theta_{\text{int}}))}$). Applying the same
argument as~\eqref{eq:arch-Ntheta-bridge} but marginalising over jump
sources (not destinations):
\begin{equation}
  E[N_{\theta_{\text{int}}} \mid \theta_X, \theta_Y, t] \cdot P_\theta(\theta_X \to \theta_Y; t)
    = \sum_{j \ne \theta_{\text{int}}} Q_\theta[j, \theta_{\text{int}}]
      \int_0^t P_\theta(\theta_X \to j; s)\,
              P_\theta(\theta_{\text{int}} \to \theta_Y; t{-}s)\, ds.
\end{equation}
With $Q_\theta[j, \theta_{\text{int}}] = \rho_{\text{chain}} \rho[\theta_{\text{int}}]$
for $j \ne \theta_{\text{int}}$, factor the destination-independent part
and use $\sum_{j \ne i} P_\theta(a \to j; s) = 1 - P_\theta(a \to i; s)$
to get
\begin{align}
  E[N_i \mid a, b, t] \cdot P_\theta(a \to b; t)
    &= \rho_{\text{chain}}\, \rho[i]\,
       \int_0^t \big(1 - P_\theta(a \to i; s)\big)\,
              P_\theta(i \to b; t{-}s)\, ds
    \label{eq:arch-arrivals-int} \\
    &= \rho_{\text{chain}}\, \rho[i] \Big[
        (1 - \rho[i])\, \rho[b]\, t
        + \big((1 - \rho[i])\, d_{ib} - d_{ai}\, \rho[b]\big) I_1
        - d_{ai}\, d_{ib}\, t\, g_t \Big],
    \label{eq:arch-arrivals-closed}
\end{align}
with $I_1 := (1 - g_t)/\rho_{\text{chain}}$, $d_{ai} := \delta_{a, i} - \rho[i]$,
and $d_{ib} := \delta_{i, b} - \rho[b]$.
Summing~\eqref{eq:arch-arrivals-closed} over $i$ recovers
\eqref{eq:arch-Ntheta-closed} (analytic identity, verified in
\texttt{test\_field\_bridge\_hr\_matches\_direct}).

\paragraph{Case-$M{=}1$ second-segment V attribution via time reversal.}
The case-$M{=}1$ second segment $[\tau_1, t]$ has resampled start
$z \sim \eqm^{(k_1)}$ at $\tau_1^+$ and observed end $Y_n$ in arch $k_1$.
Under residue-endpoint conditioning, $z$ is biased away from
$\eqm^{(k_1)}$ by the Y-endpoint. GTR reversibility gives the
bridge-conditional distribution
$\eqm^{(k_1)}[z]\, P^{k_1}(z \to Y_n; t{-}\tau_1)/\eqm^{(k_1)}[Y_n]
= P^{k_1}(Y_n \to z; t{-}\tau_1)$,
so the second segment's forward-time start $z$ is distributed exactly
as the reversed segment's target from $Y_n$. Averaging over the
case-$M{=}1$ posterior $\tau_1 \sim \exp(\delta \tau_1)/Z_1$ on $[0, t]$
(with $\delta := r'(\theta_1) - r'(\theta_X)$; the residue-endpoint
factor cancels stationarily, cf.\ appendix~\ref{par:arch-hr}),
substitute $\tau' := t - \tau_1$ and apply the GTR spectral form
$P^{k}(X \to a; \tau') = (D^k_{1/2}[a]/D^k_{1/2}[X]) \sum_\alpha
U^k[X, \alpha] U^k[a, \alpha] e^{\xi^k_\alpha \tau'}$:
\begin{equation}
  E\bigl[\, P^{k_1}(Y_n \to a; t{-}\tau_1) \bigm| M{=}1, \theta_X, \theta_1\bigr]
    = \frac{D^{k_1}_{1/2}[a]}{D^{k_1}_{1/2}[Y_n]}
      \sum_\alpha U^{k_1}[Y_n, \alpha]\, U^{k_1}[a, \alpha]\,
      E_\tau[e^{\xi^{k_1}_\alpha \tau'}],
  \label{eq:arch-Vsecond-closed}
\end{equation}
with $E_\tau[e^{\xi \tau'}] = (e^{\delta t}/Z_1) \cdot
(e^{(\xi - \delta) t} - 1)/(\xi - \delta)$ for $\xi \ne \delta$
(and the L'H\^{o}pital limit $(e^{\delta t}/Z_1) \cdot t$ otherwise).
The V second-segment contribution is then
$V[k_1, a] \mathrel{+}=
 P(\text{case-}M{=}1, \theta_X, \theta_1 \mid X_{obs}, Y_{obs}, t)
 \cdot E\bigl[\,P^{k_1}(Y_n \to a; t{-}\tau_1)\bigr]$,
and \emph{not} $\eqm^{(k_1)}[a]$ times the case-1 mass -- that shortcut
ignores the residue-endpoint bias captured by
\eqref{eq:arch-Vsecond-closed} and can be off by
$P^{k_1}(Y_n \to a; \cdot)/\eqm^{(k_1)}[a]$, which is
$\gg 1$ at $a \approx Y_n$ and short remaining branch length. Verified
in \texttt{test\_gtr\_transition\_avg\_case1\_matches\_gillespie}.

\paragraph{HR under soft leaves (PSWM observations).}
\label{par:arch-hr-soft}
The default HR pass of Section~\ref{par:arch-hr} conditions on hard leaf
residues $(X_n, Y_n) \in \{1, \dots, A\}^2$ at every cherry site. In the
PSWM training path each cherry
endpoint carries instead a distribution $\hat p^{A}_n, \hat p^{B}_n \in
\Delta^{A-1}$ derived from LG08 Felsenstein peeling at the parent /
child node of a tree branch. The bridge posterior at that branch
integrates against $\hat p^{A}_n(X_n)\hat p^{B}_n(Y_n)$:
\[
  P^{\mathrm{soft}}_{\theta_X\theta_Y}(t)
    \;=\; \sum_{X_n, Y_n} \hat p^{A}_n(X_n)\,\hat p^{B}_n(Y_n)\,
          P^{k_n}_{X_n Y_n}(t),
\]
and every HR statistic becomes the ratio $\mathcal{N} / P^{\mathrm{soft}}$
where the numerator $\mathcal{N}$ is the same integrand bilinearly weighted
by $\hat p^{A}_n, \hat p^{B}_n$. Under the F81-on-DP field the site
observations are conditionally independent given
$(\theta_X, \theta_Y, M)$, so the cluster-level ratio factorises over
sites $n = 1, \dots, N$ by a
leave-one-out product
\[
  \mathcal{N}_{\mathrm{cluster}}
    \;=\; \sum_n
      \mathcal{N}_n \cdot \!\prod_{m \ne n}\! P^{\mathrm{soft}}_m,
\]
and closed forms for every $\mathcal{N}_n$ follow directly from the
spectral primitives of Section~\ref{par:arch-hr}.

\emph{Bridge primitive.} Define the effective end-state vectors
\begin{equation}
  v^X_\alpha \;:=\; \sum_X \frac{\hat p^{A}_n(X)\, U^{k}_{X\alpha}}{d^k_X},
  \qquad
  v^Y_\alpha \;:=\; \sum_Y \hat p^{B}_n(Y)\, d^k_Y\, U^{k}_{Y\alpha},
  \label{eq:arch-hr-soft-vXvY}
\end{equation}
with $d^k_a := \sqrt{\eqm^{(k)}[a]}$ as in \eqref{eq:arch-Ubr-spectral}.
When $\hat p^{A}_n, \hat p^{B}_n$ are delta functions on $X, Y$ these reduce to
$v^X_\alpha = U^k_{X\alpha}/d^k_X$ and $v^Y_\alpha = d^k_Y U^k_{Y\alpha}$,
matching the hard-observation combination that appears in every
$X, Y$-conditional expectation. The soft bridge probability collapses
to a single spectral sum,
\begin{equation}
  P^{\mathrm{soft}}_n(t) \;=\; \sum_\alpha e^{\xi^k_\alpha t}\,
    v^X_\alpha v^Y_\alpha,
  \label{eq:arch-hr-soft-P}
\end{equation}
and the soft dwell / transition numerators inherit the same rank-1
structure that the hard forms have in $U^k_{X\alpha}, U^k_{Y\beta}$:
\begin{align}
  \mathcal{N}_n^{W, i}
    &= \sum_{\alpha, \beta} J^{\alpha\beta}(t)\,
       v^X_\alpha\, v^Y_\beta\, U^k_{i\alpha}\, U^k_{i\beta},
  \label{eq:arch-hr-soft-Wnumer} \\
  \mathcal{N}_n^{U, ij}
    &= \exch_{ij}\, d^k_i\, d^k_j
       \sum_{\alpha, \beta} J^{\alpha\beta}(t)\,
       v^X_\alpha\, v^Y_\beta\, U^k_{i\alpha}\, U^k_{j\beta},
  \qquad i \ne j.
  \label{eq:arch-hr-soft-Unumer}
\end{align}
The conditional HR statistics are $W_i^{\mathrm{soft}} = \mathcal{N}_n^{W, i} /
P^{\mathrm{soft}}_n$ and $U_{ij}^{\mathrm{soft}} = \mathcal{N}_n^{U, ij} /
P^{\mathrm{soft}}_n$; both reduce to the hard closed forms
\eqref{eq:arch-Ubr-spectral} when $\hat p^{A}_n, \hat p^{B}_n$ are delta.
The per-site compute cost is $O(A^2)$ (one $A \times A$ spectral sum) --
identical to the hard version.

\emph{Case-$M \ge 1$ free-end and second-segment primitives.} The
first-segment substitution factor $E_{\tau_1}[\text{HR at
$X_n$}\mid\delta]$ from Section~\ref{par:arch-hr} is linear in
$U^k_{X_n, :}$, so under a soft start distribution
$\hat p^{A}_n$ we replace the single row $U^k_{X_n, :}/d^k_{X_n}$ with the
weighted average
\[
  \tilde v^X_\alpha \;:=\;
    \sum_X \frac{\hat p^{A}_n(X)\, U^{k}_{X\alpha}}{d^k_X},
\]
(same construction as $v^X$ above, evaluated at the class-index
$k_X$ at $\theta_X$). The free-end dwell and its $U_{ij}$
attribution become
\begin{equation}
  W^{\mathrm{soft}}_i(\delta, t)
    \;=\; d^k_i \sum_\alpha U^k_{i\alpha}\,
       \tilde v^X_\alpha\, E_h(\xi^k_\alpha, \delta, t),
  \quad
  U^{\mathrm{soft}}_{ij}(\delta, t)
    \;=\; Q^k_{\text{off}}[i, j]\, W^{\mathrm{soft}}_i(\delta, t),
  \label{eq:arch-hr-soft-freend}
\end{equation}
with $E_h$ the truncated-exponential moment already used by the hard
case-1 primitive. The soft second-segment
transition average
\eqref{eq:arch-Vsecond-closed} is likewise reweighted:
\[
  P^{\mathrm{soft}}_{\mathrm{avg}}(a; \delta, t)
    \;=\; d^k_a \sum_\alpha U^k_{a\alpha}\,
      \tilde v^Y_\alpha\, E_\tau[e^{\xi^k_\alpha \tau'}],
\]
which is linear in $\hat p^{B}_n$ (so an unnormalised weight
$\hat p^{B}_n \odot \eqm^{(k_Y)}$, needed by the middle-segment attribution
for the joint case-jump likelihood
$L_{\mathrm{res}, j} = \prod_n \langle \hat p^{B}_n, \eqm^{(k_Y(n))}\rangle$,
plugs in unchanged).

\emph{Cluster-level aggregation.} Fix $(\theta_X, \theta_Y)$. Under the
factorisation over sites the leave-one-out identity gives
\begin{align}
  V^{\mathrm{soft}}_{\mathrm{cluster}}
    &\;=\; \rho[\theta_X]\, P^{M=0}_{\theta_X\theta_Y}(t)\,
      \Bigl(\prod_m P^{\mathrm{soft}}_m\Bigr)\,
      \sum_n \frac{V_n^{\mathrm{soft}}}{P^{\mathrm{soft}}_n},
  \label{eq:arch-hr-soft-cluster-V}
\end{align}
and analogously for $W_{\mathrm{cluster}}, U_{\mathrm{cluster}}$ in
case $M=0$; case $M \ge 1$ uses
$L_{\mathrm{res}, j}$ as the leave-one-out product with the same
structural form. In the delta limit $\hat p^{A}_n = \eta_{X_n},
\hat p^{B}_n = \eta_{Y_n}$ every $P^{\mathrm{soft}}_n$ collapses to
$P^{k_n}_{X_n Y_n}(t)$ and the aggregation matches the hard
Section~\ref{par:arch-hr} formulae exactly.

\emph{Cost.} Building $v^X, v^Y$ (and $\tilde v^X, \tilde v^Y$) takes
$O(A^2)$ per site; the spectral sums are $O(A^2)$ per site as in the
hard path. Total additional cost is a small constant factor over the
hard HR pass -- \emph{not} the naive $O(A^4)$ that arises from
enumerating $(X_n, Y_n) \in A^2$ before summing. When soft leaves are
available at every branch of every family the closed forms
\eqref{eq:arch-hr-soft-vXvY}-\eqref{eq:arch-hr-soft-cluster-V}
replace both the E-step (hard-conditional bridge posteriors on
$X_n, Y_n$) and the M-step SS accumulation of Section~\ref{par:arch-hr}
with the marginalised soft-observation analogues; the pi-arch and
$\rho_{\text{chain}}$ M-steps below apply unchanged.

\paragraph{M-steps.}
The M-steps follow the standard GTR / F81-on-DP forms from
\citet{HolmesRubin2002, Bruno1996}.

\emph{Archetype stationaries $\eqm^{(k)}$.} With
$V'[k, a] = V[k, a] + \sum_i U[k, i, a]$ (start-plus-entry count for
state $a$) and $r_k[a] = \sum_{b \ne a} \exch[a, b]\, W[k, b]$
(exchangeability-weighted opportunity), iterate
\begin{equation}
  \eqm^{(k)}[a] \;\leftarrow\; \frac{V'[k, a] + \alpha_{\text{prior}} - 1}
                                    {V_\bullet[k] + \sigma_k \cdot r_k[a]},
  \qquad \sigma_k\ \text{chosen so that } \sum_a \eqm^{(k)}[a] = 1,
  \label{eq:arch-Mstep-pi}
\end{equation}
where $V_\bullet[k] = \sum_a V[k, a]$ and $\sigma_k$ is the unique
positive root of a monotone one-dimensional equation solved by a few
Newton iterations. This is the standard GTR $\eqm$ update with
Dirichlet$(\alpha_{\text{prior}})$ pseudocounts on $\eqm^{(k)}$. The
MAP numerator $V'[k, a] + \alpha_{\text{prior}} - 1$ is negative for
under-observed states when $\alpha_{\text{prior}} < 1$; we clamp both
the numerator and the resulting $\eqm^{(k)}[a]$ elementwise to $10^{-30}$
before renormalising, which preserves sparsity (unused residues land
at $\sim\!10^{-30}$ not $\sim\!1/A$) while keeping
$\sqrt{\eqm^{(k)}[a]} \geq 10^{-15}$ well above the subnormal float64
threshold needed by the spectral decomposition~\eqref{eq:arch-Ubr-spectral}.
The default $\alpha_{\text{prior}} = 0.5$ is a sparse simplex-corner
prior encouraging peaky archetype distributions; setting
$\alpha_{\text{prior}} \geq 1$ recovers a proper MAP with Laplace-style
smoothing.

\emph{Field-chain rate $\rho_{\text{chain}}$} under
Gamma$(\alpha_\rho, \beta_\rho)$ prior:
\begin{equation}
  \rho_{\text{chain}} \;\leftarrow\;
    \frac{\alpha_\rho - 1 + \sum_q N_\theta[q]}
         {\beta_\rho + \sum_q T_q}.
  \label{eq:arch-Mstep-rho}
\end{equation}
Closed form.

\emph{Archetype selector $\rho_a$}: unchanged TSB Beta posterior mean
update on the (soft) archetype assignment counts, mirroring the field
selector $\rho$ update of Section~\ref{sec:dynamic-suppl}.

\paragraph{Rate heterogeneity ($+\Gamma+I$): per-cluster field-rate and
per-site substitution rate.}
\label{par:rate-het}
A single global $\rho_{\text{chain}}$ forces every cluster to share one
field-jump rate. Empirically most structural contacts are
orientation-conserved (static) and only a minority are genuinely
compensatory, so a global rate is dragged toward the static regime and
no cluster can flip decisively. We give each cluster its own field-rate
and each site its own substitution rate, each a discretised
Gamma-plus-invariant mixture in the sense of
\citet{yang1994maximum, yang1996among}.

\emph{Per-cluster field-rate.} Cluster $C$ draws a field-rate multiplier
$r_C$: with probability $p_I$ the \emph{invariant} value $r_C = 0$, else
one of $G$ equiprobable Gamma$(a_f, a_f)$ quantile midpoints
$\{\gamma_1,\dots,\gamma_G\}$ (mean $1$). Its field chain uses
$\rho_{\text{chain}} r_C$ in place of $\rho_{\text{chain}}$ throughout
\eqref{eq:arch-Ptheta}--\eqref{eq:arch-P2plus}, i.e.\ $g^{(C)}_\tau =
e^{-\rho_{\text{chain}} r_C \tau}$. The invariant bin is the crux: $r_C =
0 \Rightarrow g^{(C)}_\tau = 1 \Rightarrow \beta_\tau \equiv 1 \Rightarrow
P(M{\ge}1) = 0$, so $\theta$ is constant on the whole tree and the
cluster collapses to a single field draw $\theta \sim \rho$ with
independent per-column emissions under $\eqm^{(\arch(c_n,\theta))}$ --- a
static class-mixture carrying \emph{no} field flip. Marginalising $r_C$,
\begin{equation}
  L(C) = p_I\, L_0(C) + \frac{1 - p_I}{G} \sum_{g=1}^{G} L\!\left(C \mid r_C = \gamma_g\right),
  \quad
  L_0(C) = \sum_\theta \rho[\theta] \prod_n L_n^{\mathrm{stat}}(\theta),
  \label{eq:rate-het-field-marg}
\end{equation}
with $L_0$ the no-jump (static) evidence. The natural per-cluster flip
statistic is the posterior responsibility of the non-invariant bins,
\begin{equation}
  \varphi(C) \;:=\; P(r_C \neq 0 \mid C)
  \;=\; \frac{\tfrac{1-p_I}{G}\sum_g L(C \mid \gamma_g)}{L(C)},
  \label{eq:flip-posterior}
\end{equation}
replacing hard flip-detection: a cluster flips to the extent its data
prefer a nonzero field rate over the static bin, and the corpus flip
rate is $\mathbb{E}_C[\varphi(C)]$. Conserved contacts have
$\varphi \approx 0$; compensatory (salt-bridge-flip) contacts
$\varphi \to 1$.

\emph{Per-site substitution rate.} Each site $n$ carries a multiplier
$m_n$ on the residue generator~\eqref{eq:arch-Q-res},
$\exch[x_n, y] \mapsto m_n\,\exch[x_n, y]$ (with $m_n = 0$ the invariant
bin, a frozen site), drawn from a discretised Gamma$(a_s, a_s){+}I$
mixture. We fix $\{m_n\}$ \emph{once at initialisation} by empirical
Bayes --- $m_n$ set to its posterior-mean bin under the fixed LG08
per-site likelihoods (per-column conservation) --- and hold it constant.
This absorbs across-site rate variation into the substitution layer so
it is not misattributed to field dynamics (a fast column is explained by
large $m_n$, not by a spurious flip); being fixed, it is not a free
latent during training. In the covarion dichotomy of
Section~\ref{sec:f81-dp} the site-rate is a Tuffley--Steel
\emph{rate-only} modulator (distribution-preserving, residue carried),
whereas the field-rate multiplies the distribution-changing field chain
(residue resampled at each jump); the two heterogeneities act on the two
reversible layers without interfering.

\emph{Evidence unit and cached primitives.} The unit of the class
marginalisation is the per-$(\text{cluster}, \text{class-labeling})$
evidence: $K_c$ labelings for a singleton, $K_c^2$ for a doublet. It
\emph{cannot} be partitioned by archetype. A class is a \emph{sequence}
$\arch(c,\cdot)$ of archetypes across the field states, and as the shared
field $\theta$ evolves on the tree a site of class $c$ is carried through
whichever archetype $\arch(c,\theta(t))$ is live, resampled at each jump
\eqref{eq:arch-Q-field}; the substitution evidence therefore integrates
this archetype-flip trajectory (coupled across the two columns of a
doublet by the \emph{shared} $\theta$). That integration is exactly the
field-HMM forward of \autoref{prop:hier-fels}, and a class's evidence is
\emph{not} the sum of its archetypes' evidences. What does cache are the
\emph{substitution primitives} entering the number-of-jumps
decomposition \eqref{eq:arch-P0}--\eqref{eq:arch-P1-marg}: the
per-$(\text{site } n, \text{archetype } k)$ no-jump branch evidence
$\varepsilon_n^{(k)}$ and the per-$(\text{site } n, \text{ordered
archetype pair } k{\to}k')$ single-jump branch evidence
$\varepsilon_n^{(k,k')}$ (endpoint transitions under $\exch$, $\eqm$,
$m_n$ with a stationary middle segment). With frozen archetypes and
fixed site rates these primitives are \emph{invariant} under every
discrete update; only the $K_a$ archetype rate-matrix eigendecompositions
are ever recomputed (never, here). The per-$(\text{cluster},
\text{class-labeling})$ forward assembles them via the field
transition~\eqref{eq:arch-Ptheta}, at cost $O(L_{\max}^2\cdot\text{edges})$
per labeling --- reusing the cached $\varepsilon$, not re-peeling. Hence:
(i) an $\arch$ move on $(c,\theta)$ leaves the primitives untouched, re-
indexes which primitive each field state selects, and reassembles only
the labelings whose class set contains $c$ ($1$ of $K_c$ per singleton
--- the class-$c$ labeling; $2K_c{-}1$ of $K_c^2$ per doublet --- the
$(c,d)$ and $(d,c)$ labelings for all $d$); every other labeling's cached
forward is reused, so one $\arch$ move triggers only
$N_{\text{sing}} + (2K_c{-}1)N_{\text{pair}}$ field-HMM reassemblies
across the corpus; (ii) field-rate marginalisation~\eqref{eq:rate-het-field-marg}
reassembles from the same $\varepsilon$ with the scalar jump-count
weights $\beta^{(C)}_\tau, P(M\mid r_C)$ rescaled at
$\rho_{\text{chain}} r_C$ ($G{+}1$ assemblies per labeling). The
substitution eigen-cost is thus paid once per archetype and the
per-branch primitive once per $(n,k)$ and $(n,k,k')$; the residual
per-move work is the $O(L_{\max}^2\cdot\text{edges})$ field-HMM
reassembly of the affected labelings. The factorisation is exact --- a
finite discrete mixture over field-rate bins and class labelings --- the
sole approximation being the empirical-Bayes point estimate of $\{m_n\}$,
standard in $+\Gamma+I$ phylogenetics.

\begin{remark}[Exact cap-2 Felsenstein --- future refinement]
\label{rmk:exact-cap2}
Every TKF-DP cluster is size $\le 2$, so the moment-matching /
structured-mean-field forward (\autoref{sec:variational-suppl}) is not
strictly necessary: the cap-2 cluster evidence admits an \emph{exact}
Felsenstein that is also \emph{cheaper} than the naive $A^2$-state
($400$-state) coupled chain. The compound state is $(\theta, x_1, x_2)$,
but conditional on $\theta$ the two columns evolve independently, so a
no-jump branch never forms the dense $A^2 \times A^2$ transition: it
applies two per-column $A\times A$ GTRs, $P^{(\arch(c_1,\theta))} \otimes
P^{(\arch(c_2,\theta))}$, and the field is integrated in closed form via
the rank-1 collapse \eqref{eq:arch-Ptheta} (jumps resample both residues
from the new archetype stationaries, a rank-1 outer product). Tracking
the $(\theta, x_1, x_2)$ message thus costs $O(L_{\max} A^2)$ storage and
$O(L_{\max} A^2)$ per branch --- factored $A\times A$ work, not
$A^2\times A^2$ --- giving the exact cap-2 evidence at
$O(\text{edges}\cdot L_{\max} A^2)$ per $(\text{cluster},
\text{class-labeling})$ with no variational gap. Singletons ($m=1$) are
already exact. This exact pair scorer drops into the same batched
$(\text{cluster}, \text{class-labeling})\times\text{rate-bin}$ enumeration
and $\arch$-move invalidation described above; we note it as a future
refinement and use the moment-matching forward in the current
implementation.
\end{remark}

\emph{M-step.} $\rho_{\text{chain}}$ keeps its Gamma-posterior
update~\eqref{eq:arch-Mstep-rho}, with $\sum_q N_\theta[q]$ and
$\sum_q T_q$ now accumulated under the per-cluster responsibilities over
$r_C$. The invariant weight has a Beta$(a_p, b_p)$-conjugate update
\begin{equation}
  p_I \;\leftarrow\; \frac{a_p - 1 + \sum_C q_I(C)}{a_p + b_p - 2 + |\mathcal{C}|},
  \qquad q_I(C) = P(r_C = 0 \mid C) = 1 - \varphi(C),
  \label{eq:pinv-mstep}
\end{equation}
over the clusters $\mathcal{C}$, the standard $+I$ update. The Gamma
shapes $a_f, a_s$ are fixed hyperparameters (or profiled on a grid).

\paragraph{Full training algorithm.}
\label{par:arch-algorithm}
Each outer iteration of the archetype trainer applies the following
sub-steps in order (Algorithm~\ref{alg:arch-em}). We split them into
the alignment/partition-side Gibbs updates (which sample from the CRP
posterior on cluster partitions and column class labels) and the
archetype/substitution-side updates (which follow an EM-style split with
a deliberate loose coupling between the E-step for the discrete
archetype assignment $\arch$ and the M-step for the continuous
parameters $\eqm^{(k)}, \rho_{\text{chain}}, \rho_a$).

\begin{algorithm}[H]
\caption{One outer iteration of the archetype-variant trainer}
\label{alg:arch-em}
\begin{algorithmic}[1]
\State \textbf{CRP-Gibbs} on cluster partitions
       (Section~\ref{sec:dynamic-suppl}), per MSA.
\State \textbf{Class-Gibbs} on per-column class labels $c_s$, per MSA
       (if $K_c > 1$).
\State \textbf{$\alpha_z$-MH}: random-walk MH on $\log\alpha_z$ under
       a Gamma$(a_\alpha, b_\alpha)$ prior.
\State Extract cluster observations $(\text{classes}, X_{\text{obs}},
       Y_{\text{obs}}, t)$ from all MSAs.
\State \textbf{E-step (arch assignment)}: accumulate the per-$(c, \theta,
       a)$ soft-attribution counts $N[c, \theta, a]$ (residue counts
       attributed to $(c, \theta)$ buckets, marginalised over field-
       trajectory histories under a per-column Categorical emission
       model). Compute the Categorical posterior on $\arch$
\begin{equation}
  q(\arch(c, \theta) = k \mid N) \;\propto\;
    \rho_a[k]\, \prod_a \eqm^{(k)}[a]^{N[c, \theta, a]},
  \label{eq:arch-Estep-cat}
\end{equation}
and hard-assign $\hat\arch(c, \theta) = \arg\max_k q$.
\State \textbf{E-step (HR pass)}: with the fresh $\hat\arch$, accumulate
       per-$(c, \theta)$ HR sufficient statistics
       $\{V_{c, \theta, a}, U_{c, \theta, i, j}, W_{c, \theta, i}\}$
       (Section~\ref{par:arch-hr}) via
       Algorithm~\ref{alg:arch-hr-cluster} below (composition of the
       closed-form primitives \eqref{eq:arch-Ntheta-closed},
       \eqref{eq:arch-arrivals-closed},
       \eqref{eq:arch-Vsecond-closed}). Also accumulate the branch-level
       scalars $\sum_q N_\theta[q]$ and $\sum_q T_q$.
\State \textbf{M-step (aggregate)}: aggregate per-archetype sums
       $V[k, a] = \sum_{c, \theta: \hat\arch(c, \theta) = k}
       V_{c, \theta, a}$ and similarly $U[k, i, j], W[k, i]$.
\State \textbf{M-step ($\eqm^{(k)}$)}: Newton fixed-point update
       \eqref{eq:arch-Mstep-pi} on each $\eqm^{(k)}$.
\State \textbf{M-step ($\rho_{\text{chain}}$)}: closed-form Gamma
       posterior mode \eqref{eq:arch-Mstep-rho}.
\State \textbf{M-step ($\rho_a$)}: TSB posterior mean from soft
       archetype counts $n_k = \sum_{c, \theta} q(\arch(c, \theta) = k
       \mid N)$.
\State Materialise $\eqm^{(c, \theta)} = \eqm^{(\hat\arch(c, \theta))}$.
\end{algorithmic}
\end{algorithm}

\paragraph{Loose coupling between the E-step and M-step.}
Algorithm~\ref{alg:arch-em} runs the E-step for $\arch$ on the
Categorical posterior~\eqref{eq:arch-Estep-cat}, whose sufficient
statistic is the residue-count tensor $N[c, \theta, a]$; the M-step for
$\eqm^{(k)}$ uses the trajectory-based HR sufficient statistics
$\{V, U, W\}$ that carry dwell and transition information that $N$
alone does not. Formally these two updates target related but
non-identical variational bounds --- the Categorical posterior is
correct for the per-column emission model, and the GTR Newton update
is correct for the compound-CTMC trajectory model --- and their
composition is \emph{not} a monotone EM step on any single objective.
The overall procedure has a fixed point at the true maximum-likelihood
solution in the F81 limit (where the residue-count tensor $N$ captures
all the trajectory information the HR sufficient statistics do); away
from F81 the two objectives share stationary points but reach them
along different paths. Empirically the composition converges within
the K=4 top-1000 Pfam configuration used throughout this paper.

\paragraph{Cluster-level HR pass algorithm.}
For completeness, Algorithm~\ref{alg:arch-hr-cluster} states the
per-cluster HR pass invoked in step~6 above; it composes the closed
forms established in the preceding paragraphs.
\begin{algorithm}[H]
\caption{Cluster HR pass: composes the closed-form primitives}
\label{alg:arch-hr-cluster}
\begin{algorithmic}[1]
\Require Cluster observations $(\text{classes}, X_{\text{obs}},
         Y_{\text{obs}})$, branch length $t$, current
         $\{\rho, \rho_{\text{chain}}, \eqm^{(k)}, \arch\}$,
         exchangeabilities $\exch$.
\State Precompute per-archetype eigendecompositions of
       $\Q^{(k)}$ (Section~\ref{par:arch-hr}).
\State For each $(\theta_X, \theta_Y) \in \{1, \ldots, L_{\max}\}^2$:
\State \quad Compute the field-side quantities
       $P_\theta(\theta_X \to \theta_Y; t)$,
       $E[T_\theta \mid \theta_X, \theta_Y, t]$
       (Section~\ref{sec:dynamic-suppl}),
       $E[N_\theta \mid \theta_X, \theta_Y, t]$
       (Eq.~\ref{eq:arch-Ntheta-closed}),
       $E[N_i \mid \theta_X, \theta_Y, t]$
       (Eq.~\ref{eq:arch-arrivals-closed}).
\State \quad \textbf{Case $M = 0$}: for each site $n$, add the GTR
       bridge HR contribution under
       $\arch(c_n, \theta_X)$ to $\{V_{c_n, \theta_X, \cdot},
       W_{c_n, \theta_X, \cdot}, U_{c_n, \theta_X, \cdot, \cdot}\}$,
       weighted by
       $\rho[\theta_X] \cdot \beta_t(\theta_X) \cdot
       \prod_n P^{(\arch(c_n, \theta_X))}(X_n \to Y_n; t)$.
\State \quad \textbf{Case $M \geq 1$}: first-segment free-end HR from
       $X_n$ (Eq.~\ref{eq:arch-Ntheta-bridge}), last-segment
       time-reversed free-end HR from $Y_n$, and last-jump V
       attribution via
       $E[P^{(k)}(Y_n \to a; t - \tau_1)]$
       (Eq.~\ref{eq:arch-Vsecond-closed}). Middle-segment stationary HR
       accounts for the intermediate $\theta$ time.
\State Divide by $P(\text{obs})$ to convert joint sufficient statistics
       to conditional expectations required by the M-step.
\State \Return $\{V_{c, \theta, \cdot}, U_{c, \theta, \cdot, \cdot},
         W_{c, \theta, \cdot}\}, N_\theta[q], T_q, \log P(\text{obs}_q)$.
\end{algorithmic}
\end{algorithm}

\paragraph{Split-merge moves on the archetype block.}
\label{par:arch-split-merge}
The archetype budget $K_a$ is fixed at construction; under truncated
stick-breaking (TSB) the training loop cannot promote a dead slot back
into use on its own, so the effective vocabulary can collapse below
$K_a$. Symmetrically, an overloaded archetype with heavy $\rho_a[k]$
and diffuse $\eqm^{(k)}$ may be absorbing what would otherwise be two
distinct atomic types. We periodically apply hill-climbing
\emph{split-merge} moves to re-allocate archetype slots, following the
Jain-Neal split-merge motif for Dirichlet mixture models
\citep{jainneal2004} but with the full reversible MH acceptance ratio
replaced by a deterministic multinomial-log-likelihood improvement
threshold, consistent with the surrounding stochastic-EM regime.

Let
\[
  \mathcal{L}_{\text{mult}}(\eqm, \arch; N) \;=\;
  \sum_{c, \theta, a} N[c, \theta, a] \,\log\,\eqm^{(\arch[c, \theta])}[a],
  \label{eq:arch-sm-obj}
\]
be the archetype-emission log-likelihood evaluated on the current
residue-count tensor $N[c, \theta, a]$ (the same tensor used in
\eqref{eq:arch-Estep-cat}). A split move is accepted iff
$\mathcal{L}_{\text{mult}}(\eqm', \arch'; N) >
\mathcal{L}_{\text{mult}}(\eqm, \arch; N) + \tau_{\text{split}}$
with $\tau_{\text{split}} = 1$ nat (a minimum improvement threshold to
prevent split-merge from chasing numerical noise). A merge move is
accepted iff
$\mathcal{L}_{\text{mult}}(\eqm', \arch'; N) \geq
\mathcal{L}_{\text{mult}}(\eqm, \arch; N) -
\tau_{\text{merge}}^{\text{rel}} \cdot |\mathcal{L}_{\text{mult}}(\eqm, \arch; N)|$
with $\tau_{\text{merge}}^{\text{rel}} = 10^{-4}$; the tolerance is
\emph{relative} to the current LL magnitude because merging two
near-identical archetypes ($\text{TVD} < \varepsilon_{\text{tvd}}$) is
LL-neutral by construction but may nudge down by numerical
$O(10^{-6})$ due to the $\rho_a$-weighted averaging, and an absolute
tolerance calibrated to one corpus size fails on others. This
asymmetry --- strict for splits, tolerant for merges --- reflects an
implicit MDL/Ockham preference for lower archetype capacity when
the data cannot distinguish between two atomic types.

\emph{Dead slots.} A slot $k_{\text{dead}}$ counts as \emph{dead} and
eligible for revival if either
\begin{enumerate}[label=(\alph*)]
  \item $\rho_a[k_{\text{dead}}] < \varepsilon_{\text{dead}}$ (default
    $10^{-2}$): the slot has essentially no data pulling it anywhere,
    or
  \item $H\bigl(\eqm^{(k_{\text{dead}})}\bigr) < H_{\text{coll}}$
    (default $0.5$ bits): the slot's emission has collapsed onto one
    or two residues under sub-1 Dirichlet prior + SS starvation, so
    observations of any other residue at $(c, \theta)$ positions
    pointing at $k_{\text{dead}}$ contribute
    $\log \eqm^{(k_{\text{dead}})}[a]$ near $\log 10^{-30} \approx -69$,
    a per-residue penalty large enough that even a modest number of
    such observations dominates the joint LL.
\end{enumerate}
The second criterion is essential: rho-only detection misses the
failure mode where a slot has non-trivial usage ($\rho_a > 10^{-2}$)
but has collapsed to a delta-like emission, which is what actually
produces the catastrophic $\Delta \text{LL}$ drops observed in practice.

\emph{Split proposal.} Given a dead slot $k_{\text{dead}}$, choose an
overloaded slot $k_{\text{over}}$ maximising $\rho_a[k]\cdot H(\eqm^{(k)})$
(heavy \emph{and} diffuse) subject to $k_{\text{over}} \notin
\{\text{dead slots}\}$. Collect the per-$(c, \theta)$ residue-count rows
$\{N[c, \theta, \cdot] : \arch[c, \theta] = k_{\text{over}}\}$,
normalise each to the amino-acid simplex, and run weighted
$k=2$-means with $L^1$ (TVD) distance and row-mass weights. The two
centroids $\mathbf{m}_0, \mathbf{m}_1$ become the proposed
$\eqm'^{(k_{\text{over}})}$ and $\eqm'^{(k_{\text{dead}})}$; the
$(c, \theta)$ assignments split accordingly, with $\arch'[c, \theta]$
set to whichever centroid is nearest under TVD.

\emph{Merge proposal.} Enumerate archetype pairs $(k_1, k_2)$ with
$\rho_a[k_i] > \varepsilon_{\text{dead}}$ for both, and choose the pair
minimising the total-variation distance
$\tfrac{1}{2}\|\eqm^{(k_1)} - \eqm^{(k_2)}\|_1$. If the minimum is
below a threshold $\varepsilon_{\text{tvd}}$ (default $0.05$), merge by
setting
\[
  \eqm'^{(k_1)} \;=\;
  \frac{\rho_a[k_1]\,\eqm^{(k_1)} + \rho_a[k_2]\,\eqm^{(k_2)}}
       {\rho_a[k_1] + \rho_a[k_2]},
\]
redirecting all $\arch[c, \theta] = k_2$ to $k_1$, and freeing slot
$k_2$ to a uniform distribution (available for future revival by a
split).

After each accepted move the TSB weights $\rho_a$ and their underlying
Beta sticks $\beta_{a, i}$ are recomputed from the updated $\arch$ by
the same posterior-mean rule used elsewhere in the archetype block
(the archetype-selector line of the M-step in Section~\ref{par:arch-algorithm})
applied to the new occupancy counts $E[n_k]$.

\emph{On the acceptance rule.} A full Jain-Neal MH move would multiply
the $\mathcal{L}_{\text{mult}}$ ratio by (i) the ratio of Dirichlet
priors on $\eqm^{(k)}$ and (ii) the restricted-scan proposal-density
ratio $q(\text{old} | \text{new}) / q(\text{new} | \text{old})$
governing the 2-means initialisation and Gibbs mini-sweep. Both are
close to unity for the $K_a{=}8$ configurations of interest here, and
retaining them would only slow the walk in cases where the raw LL
ratio is already tightly bimodal. In the stochastic-EM setting where
the surrounding updates are not themselves a valid MH kernel, the
combined chain is not sampling from a well-defined posterior anyway,
so the simpler LL-improvement rule is preferred and does not lose
guarantees that are not already forfeit.

\emph{Where in the algorithm.} The split-merge step is invoked after
the atom / archetype M-steps, every $J_{\text{sm}}$ outer iterations.
The choice of $J_{\text{sm}}$ has qualitative consequences.
$J_{\text{sm}}{=}1$ (every iter) is the safe default when using the
sparse Dirichlet prior $\alpha_{\text{prior}} < 1$: on any single iter
the atom M-step can drive an under-observed archetype into an
entropy-collapsed configuration (all mass on one residue), and
split-merge fires at end-of-iter to revive it before the next E-step
propagates the damage. Sparser $J_{\text{sm}}$ (e.g. $J_{\text{sm}}{=}5$)
is only safe when
$\alpha_{\text{prior}} \geq 1$ so that no single M-step can collapse
an archetype past the entropy revival threshold in the
$J_{\text{sm}}{-}1$ intervening iterations. In practice we use
$J_{\text{sm}}{=}1$ and pay the modest overhead of an extra
multinomial-LL scan per iter, which is $O(K_c \cdot L_{\max} \cdot A)$
per proposal and does not depend on corpus size.

Because $\mathcal{L}_{\text{mult}}$ omits the substitution-side dwell
term $-\sum_a r_k[a]\,\eqm^{(k)}[a]$ from the full HR objective
(cf.\ \eqref{eq:arch-Mstep-pi}), an accepted split-merge move is not
guaranteed to improve the joint HR log-likelihood; empirically the
correlation between the two objectives is strong enough that accepted
moves do improve the joint LL after the next atom step re-refines
$\eqm^{(k)}$ from HR sufficient statistics under the new arch.

\paragraph{Alternative training: phylogenetic ELBO on general trees.}
\label{par:arch-phylo-elbo}
The composite-likelihood training regime factorises over independent
cherries: each cherry's parent field state $\theta_P$ is
marginalised independently against the field prior $\rho[\theta_P]$,
so different cherries at the same MSA are free to select different
$\theta_P$ values. This freedom is empirically visible: a diagnostic
run on a trained checkpoint shows a large fraction of size-$\geq\!2$
clusters where the argmax posterior $\theta_P$ varies across
cherries (Shannon entropy $\geq 1$ bit) while the per-cherry
probability of any within-cherry field jump is essentially zero (mean
$\approx 10^{-2}$ under the trained $\rho_{\text{chain}}$). The
cluster ``co-variation'' the model is reporting is thus explained by
cross-cherry $\theta_P$ selection --- a genuine pathway that the
CRP prior does not resist --- rather than by within-cherry
coordinated substitution. We call this the \emph{composite-likelihood
freedom} and it is inherent to any pair-wise cherry training over an
underlying multi-sequence MSA.

The remedy is to train against a joint likelihood that shares field
trajectories across sibling branches emanating from the same internal
node. On a tree $\mathcal{T}$ with branch lengths, the correct joint
Felsenstein-on-DP recursion carries state $(\theta_v, r_1, \ldots,
r_{m})$ at every node $v$ and is exponential in tree depth,
because the per-node conditional-likelihood-vector (CLV) factorisation
\eqref{eq:dynfield-clv-factor} holds only conditional on
$\theta_v$: at an internal node with two or more children, the
field-marginal combine re-couples the columns through the shared
parent field.

We therefore introduce a moment-matching variational family whose
form is preserved by per-edge propagation and closed under
projective combining at internal nodes. At every node $v$ restrict
the joint $(\theta_v, x_v)$ marginal to
\begin{equation}
  q_v(x_v, \theta_v) \;=\; q_v(\theta_v) \,\Bigl[
    \bigl(1 - \lambda_v(\theta_v)\bigr)
        \,\prod_{n \in \mathrm{block}} p_{n,v}(x_{v,n} \mid \theta_v)
    \;+\; \lambda_v(\theta_v)
        \,\prod_{n \in \mathrm{block}} \eqm^{(\arch[c_n, \theta_v])}(x_{v,n})
  \Bigr],
  \label{eq:arch-elbo-family}
\end{equation}
with per-node parameters $q_v(\theta_v) \in \Delta^{L-1}$
(distribution over field states), $p_{n, v}(\cdot \mid \theta_v)
\in \Delta^{A-1}$ (per-site, per-field residue marginal on the
``tracking'' component), and $\lambda_v(\theta_v) \in [0, 1]$
(mixture weight of the ``regenerated'' component: probability that
at least one field event has occurred on the path from $v$ back to
the last ancestor whose residues are still evidence-informed).

The family \eqref{eq:arch-elbo-family} is chosen so that it contains
the exact posterior in three regimes: (i) depth-1 cherries under the
$2 \times 2$ jump-case decomposition of
Section~\ref{par:arch-hr}; (ii) $\rho_{\text{chain}} \to 0$ (no field
events; $\lambda \to 0$ everywhere; $p_{n, v}$ carries the per-site
per-class $\Q$-evolution exactly); (iii) $\rho_{\text{chain}} \to
\infty$ (every edge regenerates; $\lambda \to 1$ everywhere; $q_v \to
\rho$ and $p_{n, v}$ is irrelevant). At intermediate
$\rho_{\text{chain}}$ on deeper trees the ELBO is a controlled
approximation.

The ELBO factorises over the tree:
\begin{equation}
  \mathrm{ELBO}(q) \;=\;
    \sum_{v \in \mathrm{leaves}} \mathbb{E}_{q_v}\bigl[\log P(\mathbf{X}^{\mathrm{obs}}_v \mid x_v, \theta_v)\bigr]
    + \sum_{v \to u \in \mathrm{edges}} \mathbb{E}_q\bigl[\log K_{\mathrm{edge}}(x_u, \theta_u \mid x_v, \theta_v; \tau_{vu})\bigr]
    - \sum_{v} \mathbb{E}_{q_v}\bigl[\log q_v\bigr],
  \label{eq:arch-elbo-total}
\end{equation}
with $K_{\mathrm{edge}}$ the Interp-2 per-edge kernel with
state-dependent no-jump probability $\beta(\theta_v, \tau)$ and jump
weight $W(\theta_v, \theta_u; \tau)$ constructed in the
$2\times 2$-case decomposition. Each term is closed form under the
family \eqref{eq:arch-elbo-family}: the leaf-emission term is
$O(L \cdot m \cdot A)$, the per-edge cross-entropy is $O(L^2
\cdot m \cdot A^2)$, and the per-node entropy is $O(L \cdot
m \cdot A)$.

\emph{Per-edge propagation preserves the form.} The variational
family \eqref{eq:arch-elbo-family} is stated in posterior form; the
corresponding CLV form $F_v(x_v, \theta_v)$ (with $q_v \propto
\rho \prod \eqm \cdot F_v$) is
\begin{equation}
F_v(x_v, \theta_v) \;=\; r_v(\theta_v) \prod_n A_v(x_{v, n} \mid \theta_v, n)
    \;+\; s_v(\theta_v) \cdot 1,
\label{eq:arch-elbo-clv-form}
\end{equation}
with the invariant $\langle \eqm^{(\arch[c_n, \theta])}, A_v(\cdot
\mid \theta, n) \rangle = 1$ and, correspondingly, $\langle 1, \eqm
\rangle = 1$. The scalar component of the CLV is \emph{constant}
in $x_v$ (with weight $s_v$), reflecting that a regenerated ancestor
contributes an $x_v$-independent factor $\prod_n \eqm^{(\arch[c_n,
\theta_p])}(x_v)$ to the observed-evidence integral (already absorbed
by the prior in the CLV normalisation). Substituting
\eqref{eq:arch-elbo-clv-form} into the upward message
$\widetilde U_{u \to p}(x_p, \theta_p) = \sum_{x_u, \theta_u}
K(x_u, \theta_u \mid x_p, \theta_p; \tau)\,F_u(x_u, \theta_u)$
and expanding the two-component $K$
\eqref{eq:arch-elbo-K-decomp} produces four cross-terms, each of
which is already in the target family form: no-jump $\times$ rank-1
$\to$ rank-1 with kernel $\widetilde A(a \mid \theta_p, n) := \sum_y
P^{(c_n, \theta_p)}(a, y; \tau) A_u(y \mid \theta_p, n)$ (invariant
preserved by left-stationarity $\eqm^\top P = \eqm^\top$); no-jump
$\times$ scalar $\to$ scalar with weight $\beta(\theta_p, \tau)
s_u(\theta_p)$ (using $\sum_y P^{(c_n, \theta_p)}(x_p, y; \tau) = 1$);
jump $\times$ rank-1 $\to$ scalar with weight $\sum_{\theta_u}
W(\theta_p, \theta_u; \tau) r_u(\theta_u)$ (using invariant $\langle
\eqm, A_u \rangle = 1$); jump $\times$ scalar $\to$ scalar with weight
$\sum_{\theta_u} W(\theta_p, \theta_u; \tau) s_u(\theta_u)$ (using
$\langle \eqm, 1 \rangle = 1$). The result is exactly in the family
with:
\begin{align}
r_v(\theta_p) &\;=\; \beta(\theta_p, \tau)\,r_u(\theta_p), \\
A_v(a \mid \theta_p, n) &\;=\; \widetilde A(a \mid \theta_p, n), \\
s_v(\theta_p) &\;=\; \beta(\theta_p, \tau)\,s_u(\theta_p)
    \;+\; \sum_{\theta_u} W(\theta_p, \theta_u; \tau)\,\bigl(r_u(\theta_u) + s_u(\theta_u)\bigr).
\end{align}
No projection is needed: the (rank-1 + constant) CLV form is closed
under exact per-edge propagation. The internal-node combine below
does require projection because the exact product of two family
members is not itself in the family.

\emph{Internal-node combining via moment-matching projection.} At a
binary internal node $v$ combining children $u_1, u_2$, the exact
product of two $(\text{rank-1 + scalar})$ CLVs decomposes into three
rank-1 terms and one scalar; the variational projection matches, for
each $\theta_v$: (i) the total mass into $q_v(\theta_v)$; (ii)
the per-site per-$\theta_v$ marginal of $x_{v, n}$ into
$p_{n, v}(\cdot \mid \theta_v)$; (iii) the fraction of the four-term
sum that comes from scalar-involving terms into $\lambda_v(\theta_v)$.
Each projection is closed form, cost $O(m \cdot L \cdot A)$
per combine. Non-binary internal nodes are handled by sequential
binary combines with intermediate projection.

Coordinate ascent alternates post-order updates (leaves up to root
via the combined-then-projected forward messages) and pre-order
updates (root down to leaves via backward messages). The ELBO is
monotone non-decreasing under each half-sweep and typically converges
in $10$--$50$ full sweeps at cost $O(\text{edges} \cdot L^2 \cdot
m \cdot A^2)$ per sweep.

\paragraph{Internal-node combine: moment-matching projection.}
\label{par:arch-phylo-elbo-mmedge}

Per-edge propagation is exact under the CLV form
\eqref{eq:arch-elbo-clv-form} (derived above). The only place a
moment-matching projection is required is at internal nodes with two
or more children, where the exact product of two family CLVs
$F_{u_1} \cdot F_{u_2}$ is not itself in the family.

\emph{Exact product.} Write both children in family form
\eqref{eq:arch-elbo-clv-form}: $F_{u_i}(x, \theta) = r_i \prod A_i + s_i$
for $i \in \{1, 2\}$ (dropping $\theta$-dependence for brevity). The
pointwise product decomposes into four cross-terms:
\begin{equation}
F_{u_1}(x, \theta) F_{u_2}(x, \theta)
    \;=\; \underbrace{r_1 r_2 \prod_n A_1 A_2}_{\text{rank-1, kernel } A_1 A_2}
      + \underbrace{r_1 s_2 \prod_n A_1}_{\text{rank-1, kernel } A_1}
      + \underbrace{s_1 r_2 \prod_n A_2}_{\text{rank-1, kernel } A_2}
      + \underbrace{s_1 s_2}_{\text{constant}}.
\label{eq:arch-mmcombine-exact}
\end{equation}
The three rank-1 terms have distinct per-site kernels
($A_1 A_2$, $A_1$, $A_2$) and cannot be exactly represented by a
single-rank-1 family member. A projection is required.

\emph{Projection.} We project onto the target family
$F_p = r_p \prod A_p + s_p$ by matching (i) the total
equilibrium-weighted mass $\langle \eqm_{\text{all}}, F_p(\cdot,
\theta) \rangle = r_p + s_p$ per $\theta$, and (ii) the per-site
$\eqm$-weighted marginal at each site $n$ under $\langle \eqm,
\cdot \rangle$-weighting at other sites. We assign the scalar
component $s_p$ to the constant term $s_1 s_2$ (which is already in
the family) and put the three rank-1 terms into $r_p A_p$:
\begin{align}
s_p(\theta) &\;=\; s_1(\theta) s_2(\theta), \\
r_p(\theta) &\;=\; r_1 r_2 \prod_n \langle \eqm, A_1 A_2 \rangle_n
    \;+\; r_1 s_2 \;+\; s_1 r_2, \\
A_p(a \mid \theta, n)
    &\;=\; \frac{r_1 r_2 \bigl(\prod_{n' \ne n} \langle \eqm, A_1 A_2 \rangle_{n'}\bigr) A_1(a) A_2(a)
                     + r_1 s_2\, A_1(a) + s_1 r_2\, A_2(a)}
                    {r_p(\theta)}.
\end{align}
Each cross-term contributes to the per-site marginal at site $n$
via its $\prod_{n' \ne n} \langle \eqm, \cdot \rangle$-weighted rest
factor (all equal to $1$ except for the $A_1 A_2$ term, which picks up
$\prod_{n' \ne n} \langle \eqm, A_1 A_2 \rangle_{n'}$). The
$s_1 s_2$ constant term contributes only to $s_p$; it drops out of
$A_p$'s numerator because the family's scalar-per-site marginal $s_p
\cdot \langle \eqm, 1 \rangle_{n'} = s_p$ absorbs it directly. The
family invariant $\langle \eqm, A_p \rangle_n = 1$ follows by direct
computation using $\langle \eqm, A_i \rangle = 1$. This is the
implementation of \texttt{mm\_combine} in \texttt{mm\_clv.py}.

\emph{Cost.} Building per-site inner products $\langle \eqm,
A_1 A_2 \rangle_n$ is $O(A)$ per (site, $\theta$); combining is $O(A)$
per site. Total $O(m L A)$ per combine. Non-binary internal nodes
reduce to sequential binary combines with intermediate projection.

\emph{Approximation quality.} The projection is exact when
$\rho_{\text{chain}} \to 0$ (no field jumps, so $s_i = 0$ throughout
and only the $r_1 r_2 A_1 A_2$ term survives, which is already
rank-1). It is also exact at depth-1 cherries. At depth-2 and deeper,
per-branch $s > 0$ activates the projection and induces the
first-order moment-matching bias. Empirically the residual is
$< 1\%$ in log-likelihood on random depth-2 trees at
$\rho_{\text{chain}} \in [0.1, 1.5]$
(cf.\ \texttt{tests/phylo\_elbo/test\_mm\_clv\_depth2\_exact.py}).

The rest of this paragraph is discarded --- it derived a projection
under an incorrect CLV form that took the scalar as
$\prod_n \eqm^{(\arch)}(x_n)$ rather than the constant $1$.
Under the correct form \eqref{eq:arch-elbo-clv-form}, per-edge
propagation is exact and only the internal-node combine needs the
projection above.

\iffalse

The upward Felsenstein message from $u$ to
its parent $p$ over branch $\tau$ is
\begin{equation}
\widetilde U_{u \to p}(x_p, \theta_p)
    \;=\;
    \sum_{x_u, \theta_u} K(x_u, \theta_u \mid x_p, \theta_p; \tau)
        \,F_u(x_u, \theta_u),
    \label{eq:arch-mmedge-exact-U}
\end{equation}
with $K$ the two-component kernel of \eqref{eq:arch-elbo-K-decomp}.
Substituting the no-jump / jump split of $K$ and the rank-1 / scalar
split of $F_u$ gives four cross-terms $T_A + T_B + T_C + T_D$:
\begin{align}
T_A &\;=\; \beta(\theta_p, \tau)\,r_u(\theta_p) \prod_n
    \widetilde A(x_{p,n} \mid \theta_p, n),
    \quad\text{with}\quad
    \widetilde A(a \mid \theta, n) := \sum_y P^{(c_n, \theta)}(a, y; \tau)\, A_u(y \mid \theta, n),
    \label{eq:arch-mmedge-TA} \\
T_B &\;=\; \beta(\theta_p, \tau)\,s_u(\theta_p) \prod_n
    g(x_{p,n} \mid \theta_p, n),
    \quad\text{with}\quad
    g(a \mid \theta, n) := \sum_y P^{(c_n, \theta)}(a, y; \tau)\, \eqm^{(\arch[c_n, \theta])}(y),
    \label{eq:arch-mmedge-TB} \\
T_C &\;=\; \Bigl[\sum_{\theta_v} W(\theta_p, \theta_v; \tau)\, r_u(\theta_v)\Bigr]\, \prod_n \mathbf{1}(x_{p,n}),
    \label{eq:arch-mmedge-TC} \\
T_D &\;=\; \Bigl[\sum_{\theta_v} W(\theta_p, \theta_v; \tau)\, s_u(\theta_v)
    \prod_n \| \eqm^{(\arch[c_n, \theta_v])}\|_2^2\Bigr]\, \prod_n \mathbf{1}(x_{p,n}).
    \label{eq:arch-mmedge-TD}
\end{align}
Here $\widetilde A$ is the tracking marginal transported to $p$ by the
substitution kernel (forward action, $A_u \mapsto P \cdot A_u$ at each
site); $g$ is the \emph{transported equilibrium}, which is \emph{not}
equal to $\eqm^{(\arch[c_n, \theta_p])}$ in general. The failure is
elementary: while stationarity gives $\eqm^\top P = \eqm^\top$, i.e.\
$\sum_y \eqm(y)\,P(y, a) = \eqm(a)$ (summing $y$ into the
\emph{source} slot of $P$), the transported equilibrium sums $y$ into
the \emph{destination} slot, $\sum_y P(a, y)\,\eqm(y)$, which is a
different quantity. It is a distribution over $a$ only if $P$ is also
column-stochastic ($\sum_a P(a, y) = 1$ for every $y$); for a
reversible $P$, this holds iff $\eqm$ is uniform. In particular,
\begin{equation}
\sum_a\sum_y P(a, y)\,\eqm(y)
    \;=\; \sum_y \eqm(y) \sum_a P(a, y),
    \label{eq:arch-mmedge-g-failure}
\end{equation}
where the column sums $\sum_a P(a, y)$ are generally $\neq 1$, so $g$
is not even a probability distribution over $a$. Consequently
$g \neq \eqm$ in general.

The $x_p$-independence of $T_C$ and $T_D$ reflects that the jump
component of $K$ carries a source-side factor $\prod_n
\eqm^{(\arch[c_n, \theta_v])}(x_{v, n})$ evaluated at $x_v$ (the
\emph{child}), not $x_p$; after summing $x_v$ against the child's
rank-1 or scalar structure using $\langle \eqm, A_u \rangle = 1$ and
$\langle \eqm, \eqm \rangle = \|\eqm\|_2^2$ respectively, the residual
depends only on $\theta_p$.

\emph{The family is not closed.} The exact message
$\widetilde U_{u \to p}$ is a sum of three distinct per-site rank-1
factors $\{\widetilde A, g, \mathbf{1}\}$ plus no
$\prod_n \eqm^{(\arch[c_n, \theta_p])}$-scalar term. The
(rank-1 + scalar) family carries a single tracking $A$ and a single
$\prod_n \eqm(x_p)$ scalar term, so we cannot express
$\widetilde U_{u \to p}$ in the family without approximation.

\emph{Moment-matching projection.} We project onto the family
\begin{equation}
U_{u \to p}(x_p, \theta_p) \;=\; r_v(\theta_p) \prod_n A_v(x_{p,n} \mid \theta_p, n)
    \;+\; s_v(\theta_p) \prod_n \eqm^{(\arch[c_n, \theta_p])}(x_{p,n})
    \label{eq:arch-mmedge-U-form}
\end{equation}
by matching, for each $\theta_p$: (i) total $\theta_p$-mass
$\langle \eqm, \widetilde U(\cdot, \theta_p) \rangle$; (ii) per-site
$x_{p, n}$-marginals $\langle \eqm_{\backslash n}, \widetilde U(\cdot,
\theta_p) \rangle_{x_p \setminus x_{p, n}}$ under the equilibrium
weighting that gives $\langle \eqm, \cdot \rangle = 1$ for a family
member. The projection is:
\begin{itemize}[leftmargin=1.6em]
\item \emph{No-jump-derived mass $\to$ rank-1.} Assign the total
    $\theta_p$-mass of $T_A + T_B$ to $r_v(\theta_p)$ and set
    $A_v(\cdot \mid \theta_p, n)$ to the corresponding weighted mixture
    of $\widetilde A$ and $g$:
    \begin{align}
      r_v(\theta_p)
        &\;=\; \beta(\theta_p, \tau)\,\bigl[r_u(\theta_p) + s_u(\theta_p)\,\gamma_g(\theta_p)\bigr],
        \label{eq:arch-mmedge-rv} \\
      A_v(a \mid \theta_p, n)
        &\;=\; \frac{r_u(\theta_p)\,\widetilde A(a \mid \theta_p, n)
                     + s_u(\theta_p)\,\gamma_g(\theta_p)\,\widehat g(a \mid \theta_p, n)}
                    {r_u(\theta_p) + s_u(\theta_p)\,\gamma_g(\theta_p)},
        \label{eq:arch-mmedge-Av}
    \end{align}
    where $\gamma_{g, n}(\theta_p) := \langle \eqm^{(\arch[c_n, \theta_p])}, g(\cdot \mid \theta_p, n) \rangle
    = \|\eqm^{(\arch[c_n, \theta_p])}\|_2^2$ (site-local; independent of
    $A_u$), $\widehat g := g / \gamma_{g, n}$ is the normalised
    transported equilibrium (so $\langle \eqm, \widehat g \rangle = 1$),
    and $\gamma_g(\theta_p) := \prod_n \gamma_{g, n}(\theta_p)$ is the
    cluster-level correction. The cluster-level $\gamma_g$ appears in
    both numerator and denominator because $T_B$'s per-site marginal at
    site $n$ picks up $\prod_{n' \ne n} \langle \eqm, g \rangle_{n'} =
    \gamma_g / \gamma_{g, n}$ from the other sites and $\gamma_{g, n}$
    from the local $g = \gamma_{g, n} \widehat g$, whose product is
    $\gamma_g$ independent of $n$. By construction $\langle \eqm, A_v
    \rangle = 1$ (weighted average of two unit-invariant factors), and
    the per-site $\eqm$-weighted second moment
    $\sum_a \eqm(a) A_v(a \mid \theta_p, n)^2$ is the closest family
    match to the exact $T_A + T_B$ under this per-site marginalisation.

\item \emph{Jump-derived mass $\to$ scalar-$\eqm$.} Assign the total
    equilibrium-weighted mass of $T_C + T_D$ (constant in $x_p$) to
    $s_v(\theta_p)$ under the family's $\prod_n \eqm^{(\arch[c_n,
    \theta_p])}$ representation. The projected term $s_v(\theta_p)
    \prod_n \eqm^{(\arch[c_n, \theta_p])}(x_p)$ has equilibrium-weighted
    total mass $s_v(\theta_p) \prod_n \|\eqm^{(\arch[c_n, \theta_p])}\|_2^2
    = s_v(\theta_p)\,\gamma_g(\theta_p)$, so mass matching gives
    \begin{equation}
      s_v(\theta_p) \;=\; \frac{1}{\gamma_g(\theta_p)}
        \sum_{\theta_v} W(\theta_p, \theta_v; \tau)
        \Bigl[r_u(\theta_v) + s_u(\theta_v)\,\gamma_\eqm(\theta_v)\Bigr],
        \qquad
        \gamma_\eqm(\theta) := \prod_n \|\eqm^{(\arch[c_n, \theta])}\|_2^2.
      \label{eq:arch-mmedge-sv}
    \end{equation}
    ($\gamma_\eqm(\theta)$ and $\gamma_g(\theta)$ are the same function
    of $\theta$; two names are used only to disambiguate the source
    and destination archetype indices in the sum.) With this choice
    the total equilibrium-weighted $\theta_p$-mass matches
    $\widetilde U$ exactly.
\end{itemize}

\emph{Preservation and approximation summary.} The projection
\eqref{eq:arch-mmedge-rv}--\eqref{eq:arch-mmedge-sv} exactly preserves:
(P1) the family invariant $\langle \eqm^{(\arch[c_n, \theta_p])},
A_v \rangle = 1$ (verified below); (P2) the total $\theta_p$-marginal
$\langle \eqm, U_{u \to p} \rangle = \langle \eqm, \widetilde U \rangle$;
(P3) the per-site $\eqm$-weighted first moment
$\langle \eqm A_v(\cdot \mid \theta_p, n) \rangle = \langle \eqm,
\text{no-jump per-site marginal at site } n \rangle$. Two
approximations are introduced:
\begin{enumerate}[label=(A\arabic*),leftmargin=2.0em]
\item \emph{No-jump joint-marginal collapse.} The exact no-jump
    component is $r_u \prod_n \widetilde A + s_u \prod_n g$, a
    \emph{sum} of two rank-1 factors; the projection replaces it with
    $\prod_n [\text{mixture of } \widetilde A, g]$, a \emph{product}
    of per-site mixtures. This is a first-order moment-matching
    approximation; it is exact when either $s_u = 0$ (only tracking
    mass) or $\widetilde A = g$ at every site (equivalent to $A_u = \eqm$
    at every site, which occurs in the ``all sites regenerated''
    limit).
\item \emph{Jump-scalar reshaping.} The exact jump-derived mass is
    constant in $x_p$; the projection reshapes it as
    $s_v \prod_n \eqm(x_p)$, which peaks at
    $x_p = \arg\max_x \eqm(x)$. Since the family's scalar term will be
    consumed together with a prior $\prod_n \eqm(x_p)$ factor at the
    posterior evaluation (equation \eqref{eq:arch-elbo-post}), the
    approximation shifts only the higher-order moments beyond
    $\langle \eqm, \cdot \rangle$, and it is exact for the total mass
    and vanishing per-site variance in the $\rho[\theta] \to
    \delta_{\theta_0}$ single-field limit.
\end{enumerate}

\emph{Verification of the family invariant $\langle \eqm, A_v \rangle
= 1$.} By linearity and $\langle \eqm, A_u \rangle = 1$,
$\langle \eqm, \widetilde A \rangle
    = \sum_x \eqm(x) \sum_y P(x, y) A_u(y)
    = \sum_y A_u(y) \sum_x \eqm(x) P(x, y)
    = \sum_y A_u(y) \eqm(y) = 1$
using $\eqm^\top P = \eqm^\top$ (left-invariance of the stationary).
For $g$, $\langle \eqm, g \rangle
    = \sum_x \eqm(x) \sum_y P(x, y) \eqm(y)
    = \sum_y \eqm(y) \sum_x \eqm(x) P(x, y)
    = \sum_y \eqm(y)^2 = \|\eqm\|_2^2$; so
$\langle \eqm, \widehat g \rangle = 1$ by construction. The numerator
of \eqref{eq:arch-mmedge-Av} thus has $\eqm$-weighted total
$r_u \cdot 1 + s_u \gamma_g \cdot 1 = r_u + s_u \gamma_g$, which is the
denominator, so $\langle \eqm, A_v \rangle = 1$. $\square$

\emph{Cost.} Building $\widetilde A$ costs $O(A^2)$ per (site,
$\theta$) via one $P^{(c_n, \theta)}$ mat-vec; $g$ is precomputable
from the substitution eigen-cache and reused across all edges of the
same $\tau$ if branch lengths are binned. The mixture combine
\eqref{eq:arch-mmedge-Av} is $O(A)$ per site. The jump-scalar update
\eqref{eq:arch-mmedge-sv} is $O(L)$ per $\theta_p$ once $\gamma_\eqm$
and the $\theta_v$-marginals of $r_u, s_u$ are cached. Aggregate:
$O(m \cdot L \cdot A^2 + L^2)$ per edge, matching the informal cost
declared for \eqref{eq:arch-elbo-total}.

\fi

\paragraph{Alternative training: tied-$\theta$ MCMC on whole-tree
histories.}
\label{par:arch-lg08-is}
A third route that side-steps both the variational family of
Section~\ref{par:arch-phylo-elbo} and the soft closed forms of
Section~\ref{par:arch-hr-soft} is to sample the tree history end to
end and compute HR sufficient statistics on the fully-observed
trajectory. The core insight is that composite training treats each
branch as an independent ``cherry'' with its own marginalised parent
field state $\theta_P$, yet all branches emanating from an internal
node $V$ physically share the same $\theta_V$: composite likelihood
is deliberately breaking a constraint the tree imposes. Tying
$\theta$ at every internal node undoes the break directly, and does
so \emph{without} importance-weight variance, because with tied $\theta$
the sample is drawn from (an approximation to) the tree-joint TKF-DP
posterior and every HR estimator is just a Monte-Carlo average.

The scheme is a two-block Gibbs step nested in the outer SVI iter:
$X \sim p(X \mid y, \mathcal{T}, \Theta)$ followed by
$\theta \sim p(\theta \mid X, y, \mathcal{T}, \Theta)$ per cluster, plus a
binary branch-jump indicator $M_v$ for $\rho_{\text{chain}}$
identifiability. Each block is a Felsenstein pass over a discrete
state (residue or field), so nothing expensive appears.

\emph{Notation.} Fix a family with guide tree $\mathcal{T}$, leaf set
$\mathcal{L}$ with observed residues $y_l^s$ for leaf $l$ and site $s$,
and branch lengths $\{\tau_v\}_{v \ne \text{root}}$. Let
$\eqm^{(c, \theta)}$ be the per-(class, field) stationary, and $Q^{(c,
\theta)}_{ij} = \exch_{ij}\,\eqm^{(c, \theta)}(j)$ for $i \ne j$ the
associated GTR generator. Write $X = (x_v^s)$ for the residues at every
internal-node/site pair. Cluster assignments $\mathcal{C}$ partition
the $L$ columns of the alignment; each cluster carries its own field
trajectory $\theta_v$ on $\mathcal{T}$.

\emph{Block~1: sample $X$ (residues).} Use the standard bottom-up /
top-down Felsenstein sampler on the tree with the class-conditional
field-marginal transition $P^{(c_s)}_{\mathrm{marg}}(\tau) =
\sum_\theta \rho[\theta] \exp(Q^{(c_s, \theta)} \tau)$ as the per-site
kernel. This is a proper row-stochastic matrix (a mixture of
row-stochastic matrices), even though it is not the exponential of a
single generator, so peeling and stochastic mapping are unmodified.
The draw preserves branch-length correlations exactly and delivers a
whole-tree residue sample per site in $O(|\mathcal{T}| \cdot A^2)$ per
class per branch-length bin.

\emph{Block~2: sample $\theta$ (field, per cluster).} Given $X$, the
per-branch cluster emission under Interp~2 factorises in a form
convenient for Felsenstein-on-$\theta$. Define
$\beta(\theta) := \exp\!\bigl(-\rho_{\mathrm{chain}}(1 - \rho[\theta])
\tau_v\bigr)$ (per-branch survival probability of $\theta$),
\begin{align*}
  \Sigma_{\mathrm{prod}}(\theta)
    &\;:=\; \prod_{s \in \mathcal{C}}
       \bigl[\exp(Q^{(c_s, \theta)} \tau_v)\bigr]_{x_p^s, x_c^s}, \\
  \pi_{\mathrm{prod}}(\theta)
    &\;:=\; \prod_{s \in \mathcal{C}} \eqm^{(c_s, \theta)}(x_c^s).
\end{align*}
Then the per-branch cluster emission conditional on both endpoints'
field values is
\begin{equation}
  E_{\mathcal{C}, v}(\theta_p, \theta_c)
    \;=\;\underbrace{\beta(\theta_p)\,\Sigma_{\mathrm{prod}}(\theta_p)\,
                     \delta_{\theta_c = \theta_p}}_{M = 0\ \mathrm{term}}
      \;+\;\underbrace{(1 - \beta(\theta_p))\,\rho[\theta_c]\,
                       \pi_{\mathrm{prod}}(\theta_c)}_{M \ge 1\ \mathrm{term,\ Interp~2\ resample}}.
  \label{eq:arch-lg08is-branch-emission}
\end{equation}
The factorisation $E = A(\theta_p)\,\delta + B(\theta_p)\,C(\theta_c)$ with
$A = \beta \cdot \Sigma_{\mathrm{prod}}$, $B = 1 - \beta$,
$C = \rho \cdot \pi_{\mathrm{prod}}$ makes the Felsenstein-on-$\theta$
child update at a node $v$
\begin{equation}
  \phi_{v \leftarrow c}(\theta_v)
    \;=\; A_{v\to c}(\theta_v)\,\mathrm{CLV}_c(\theta_v)
      \;+\; B_{v\to c}(\theta_v)\,\bigl\langle C_{v\to c},\,\mathrm{CLV}_c\bigr\rangle,
  \label{eq:arch-lg08is-clv-theta}
\end{equation}
which costs $O(L)$ (not $O(L^2)$) per branch per cluster. The full CLV
pass is $O(m \cdot L \cdot |\mathcal{T}|)$ per cluster for building the
$A, B, C$ caches plus $O(L \cdot |\mathcal{T}|)$ for combining
messages. Top-down sampling of $\theta_v$ uses \eqref{eq:arch-lg08is-branch-emission}
directly.

\emph{Block~3: sample $M_v$ per branch.} Once $(\theta_p, \theta_c)$ is
drawn at branch $v$'s endpoints,
\begin{equation}
  P(M_v = 0 \mid \theta_p, \theta_c, X, \tau_v)
    \;=\; \begin{cases}
       0 & \theta_p \ne \theta_c, \\
       \dfrac{A_v(\theta_p)}{A_v(\theta_p) + B_v(\theta_p) \, C_v(\theta_p)}
         & \theta_p = \theta_c,
    \end{cases}
  \label{eq:arch-lg08is-M-binary}
\end{equation}
gives a Bernoulli sample of $M_v \in \{0, 1\}$. The Rao-Blackwellised
$\rho_{\mathrm{chain}}$ SS increment on branch $v$ is
\begin{equation}
  \Delta N_\theta^v \;=\; M_v \cdot E[N_\theta \mid \theta_p, \theta_c, \tau_v, M \geq 1],
  \label{eq:arch-lg08is-N-inc}
\end{equation}
where the case-$M \geq 1$ expectation
\begin{equation}
  E[N_\theta \mid \theta_p, \theta_c, \tau_v, M \geq 1]
    \;=\; \frac{E[N_\theta \mid \theta_p, \theta_c, \tau_v]\,
                 P_\theta(\theta_p \to \theta_c; \tau_v)}
                {P_\theta(\theta_p \to \theta_c; \tau_v)
                  - \beta(\theta_p)\,\delta_{\theta_p, \theta_c}}
  \label{eq:arch-lg08is-EN-M1}
\end{equation}
divides the whole-branch bridge form
\eqref{eq:arch-Ntheta-closed} by the $M \geq 1$ mass $P_\theta - \beta
\delta$. Summing $\Delta N_\theta^v$ over branches gives
$N_\theta^{\mathrm{sum}}$ for the $\rho_{\mathrm{chain}}$
Gamma-posterior M-step. A bare $\sum_v M_v$ (as in an early
implementation) undercounts jumps on long branches where the
$M \geq 1$ posterior concentrates mass above 1, and self-reinforcingly
drives $\rho_{\mathrm{chain}} \to 0$ (fewer $M = 1$ samples, lower
estimate, fewer samples, \ldots); the closed-form expectation avoids
this collapse. A finer sample of the exact jump count under the
truncated-Poisson posterior would refine the tail but is unnecessary
once~\eqref{eq:arch-lg08is-EN-M1} is used.

\emph{HR sufficient statistics.} For each branch $v$ in each cluster
$\mathcal{C}$ with the sampled $(\theta_p, \theta_c, M_v)$ and the
observed residues $(x_p^s, x_c^s)_{s \in \mathcal{C}}$, we accumulate:
\begin{itemize}
  \item[$M_v = 0$] (so $\theta_p = \theta_c = \theta$): per site $s$,
    standard GTR bridge HR at archetype $k = \mathrm{arch}(c_s, \theta)$:
    increment $V[c_s, \theta, x_p^s]$ by 1 and add the closed-form
    dwell and transition contributions $W_i^{\mathrm{br}},
    U_{ij}^{\mathrm{br}}$ from \eqref{eq:arch-Ubr-spectral} at endpoints
    $(x_p^s, x_c^s)$ to $W[c_s, \theta, \cdot]$ and $U[c_s, \theta, \cdot, \cdot]$.
  \item[$M_v = 1$] (means $M \geq 1$ field jumps on the branch;
    intermediate trajectory unobserved, but $(\theta_p, \theta_c)$ are
    fixed by Block 2 and $(x_p^s, x_c^s)$ by Block 1): per site $s$
    apply the case-$M \geq 1$ closed forms of
    Section~\ref{par:arch-hr} with $(\theta_X, \theta_Y) = (\theta_p,
    \theta_c)$ \emph{fixed}, i.e.\ skipping the $L \times L$
    marginalisation over $(\theta_X, \theta_Y)$ that appears in the
    composite path. The three-segment attribution is:
    \begin{itemize}
      \item First segment, archetype $k_p := \mathrm{arch}(c_s,
        \theta_p)$, forward end $x_p^s$, duration $\tau_1$ under the
        exact mixture density~\eqref{eq:arch-tau1-exact}. Add
        the start-boundary $V^{\text{first}}[a] = \delta_{a, x_p^s}$
        to $V[c_s, \theta_p, \cdot]$, $W^{\text{first}}_i$ from
        \eqref{eq:arch-Wfirst-mix} to $W[c_s, \theta_p, \cdot]$, and
        $U^{\text{first}}_{ij}$ from~\eqref{eq:arch-Ufirst-mix} to
        $U[c_s, \theta_p, \cdot, \cdot]$.
      \item Middle segments (residual field states $\theta_{\text{int}}
        \in \{1, \dots, L\}$) at their respective archetypes
        $k_{\text{int}} := \mathrm{arch}(c_s, \theta_{\text{int}})$.
        Compute the bridge-conditional-on-$M \geq 1$ quantities
        \begin{align*}
          E[T_{\theta_{\text{int}}} \mid \theta_p, \theta_c, \tau_v, M \geq 1]
            &\;=\; \frac{E[T_{\theta_{\text{int}}} \mid \theta_p, \theta_c, \tau_v]\,
                          P_\theta(\theta_p \to \theta_c; \tau_v)
                          - \beta(\theta_p) \tau_v\, \delta_{\theta_p, \theta_c}\,
                            \delta_{\theta_{\text{int}}, \theta_p}}
                         {P_\theta(\theta_p \to \theta_c; \tau_v)
                          - \beta(\theta_p)\, \delta_{\theta_p, \theta_c}}, \\
          E[N_{\theta_{\text{int}}} \mid \theta_p, \theta_c, \tau_v, M \geq 1]
            &\;=\; \frac{E[N_{\theta_{\text{int}}} \mid \theta_p, \theta_c, \tau_v]\,
                          P_\theta(\theta_p \to \theta_c; \tau_v)}
                         {P_\theta(\theta_p \to \theta_c; \tau_v)
                          - \beta(\theta_p)\, \delta_{\theta_p, \theta_c}},
        \end{align*}
        from~\eqref{eq:arch-Ntheta-closed}
        and~\eqref{eq:arch-arrivals-closed}. Subtract first- and
        last-segment residuals~\eqref{eq:arch-taumid}--\eqref{eq:arch-Nmid}
        to get $\tau_{\text{mid}}(\theta_{\text{int}})$ and
        $N_{\text{mid}}(\theta_{\text{int}})$, then add
        \eqref{eq:arch-mid-V}--\eqref{eq:arch-mid-U} to
        $V[c_s, \theta_{\text{int}}, \cdot]$, $W[c_s, \theta_{\text{int}},
        \cdot]$, $U[c_s, \theta_{\text{int}}, \cdot, \cdot]$.
      \item Last segment, archetype $k_c := \mathrm{arch}(c_s,
        \theta_c)$, observed end $x_c^s$, duration $\tau_v - \tau_M$
        under the mixture density with $\theta_X \leftrightarrow
        \theta_Y$ exchanged in~\eqref{eq:arch-tau1-exact}. Add
        $V^{\text{last}}_a$ from~\eqref{eq:arch-Vlast-mix},
        $W^{\text{last}}_j$ from~\eqref{eq:arch-Wlast-mix}, and
        $U^{\text{last}}_{ij}$ from~\eqref{eq:arch-Ulast-mix} to
        $V[c_s, \theta_c, \cdot]$, $W[c_s, \theta_c, \cdot]$,
        $U[c_s, \theta_c, \cdot, \cdot]$.
    \end{itemize}
    Also increment $N_\theta^{\mathrm{sum}}$ by $\Delta N_\theta^v$
    from~\eqref{eq:arch-lg08is-N-inc}--\eqref{eq:arch-lg08is-EN-M1}.
\end{itemize}
Sum over branches, clusters, and families gives the same per-$(c_s,
\theta)$-indexed $(V, U, W)$ tensors and scalar
$(N_\theta^{\mathrm{sum}}, T^{\mathrm{sum}})$ statistics the composite
pass yields; downstream aggregation to per-archetype $(V_k, U_k, W_k)$
via $\mathrm{arch}(c_s, \theta) = k$ and the $\eqm^{(k)}$ /
$\rho_{\mathrm{chain}}$ M-steps of Section~\ref{par:arch-hr} apply
unchanged.

\emph{Why no importance weights.} We are sampling directly from an
approximation of the tree-joint TKF-DP posterior (Block~1 uses the
class-marginal transition; Block~2 is exact given $X$ under Interp~2
with Interp~1 resample). The M-step estimator is
\begin{equation}
  \widehat{E}_\Theta\!\bigl[\mathrm{SS}_\mathcal{C}\bigr]
    \;=\; \mathrm{SS}_\mathcal{C}\!\bigl(X, \theta, M, y\bigr),
  \label{eq:arch-lg08is-SS-tied}
\end{equation}
a plain Monte-Carlo evaluation on the sampled trajectory. There is no
target/proposal ratio, no ESS collapse, and no per-cluster or
per-cluster-size weight variance. Cross-cluster $\theta$-independence
is respected by construction (each cluster gets its own draw);
within-cluster $\theta$ is tied at every internal node by construction
too.

\emph{Failure modes of the intermediate designs, for the record.} An
earlier version of this section proposed
\emph{importance sampling} against the composite target with a
class-marginal Felsenstein proposal. Three attempts, in order, each
diagnosed and abandoned:
\begin{itemize}
  \item[(i)] Store per-node marginals $p_v(a) \propto \eqm(a)\,L_v(a)$
    and draw each branch endpoint independently. Discards the parent-child
    correlation carried by $P(\tau_v)$; empirically the sampled tuples
    are $15$--$25\%$ less conserved than LG08 predicts and
    $\rho_{\mathrm{chain}}$ diverges to $> 1$ within an iteration.
  \item[(ii)] Store CLVs and draw whole-tree histories under a
    single-stationary LG08 proposal. Fixes (i)'s branch-correlation
    problem but the per-branch log-ratio between LG08 and the
    class-marginal target accumulates to $\pm 20$--$50$ nats over the
    $\sim\!2N{-}1$ nodes of a deep Pfam tree, collapsing
    $\mathrm{ESS}/K$ to $\sim\!6 \times 10^{-3}$ at $K_c = 4$.
  \item[(iii)] Switch the proposal to class-conditional GTR and hope
    the residual is $O(\tau^2)$. It is, but the numerator of the
    correction weight ratio was still the composite target
    $\Sigma_{\mathrm{cherry}}(X, Y; \tau) = \eqm(X) \cdot P(X \to Y; \tau)$
    while the denominator was the sampler marginal $q_{\mathrm{tree}}(X_p)
    \cdot P(X_p \to X_c)$. Under GTR reversibility this leaves an
    unavoidable $\log \eqm^{(c)}(X_p) - \log q_{\mathrm{tree}}(X_p)$ per
    site, which at leaves (where $q_{\mathrm{tree}}$ collapses to a
    delta on the observed residue) is $\approx -\log(1/A) \approx 3$
    nats per site and accumulates linearly in cluster size $m$. That
    means larger clusters get exponentially down-weighted in the M-step
    for the wrong reason (tree-conditional vs stationarity mismatch,
    not field-coupling signal). Measured
    $\log w /m \approx -2.85$ constant across $m = 1, \dots, 8$ at
    iter 1, confirming the log-$\eqm$ baseline is dominant.
\end{itemize}
Tied-$\theta$ MCMC dodges all three failure modes because it does not
need to reweight anything: the SS is computed on the trajectory that
was drawn, whose target the sampler already approximates.

\emph{Preprocessing footprint.} Preprocessing stores only the tree
topology (a $2N{-}1$ int parent map), branch lengths (a
$2N{-}1$ float64 vector) and the observed leaf residues (an
$N \times L$ int8 matrix). Storage is $O(N L)$ per family, dominated by
the leaf-residue matrix and totalling only $\sim\!10$ MB for the
top-1000 corpus. The CLVs for Block~1 are recomputed per iter from the
current $\eqm^{(c)}$; the $A, B, C$ tables for Block~2 are recomputed
per iter per cluster from the sampled $X$ and the current
$\eqm^{(c, \theta)}$. Both fit comfortably in RAM at Pfam scale.

\paragraph{Backward recursion for the moment-matching family.}
\label{par:arch-phylo-elbo-backward}

We now derive the second half of the message-passing structure needed
to extract per-branch pair posteriors: the \emph{downward} conditional
partials $D_v(x_v, \theta_v)$ and their propagation through the
variational family \eqref{eq:arch-elbo-family}. We give the exact
recursion first (linear-algebra form), then show it preserves the
(rank-1 + scalar) CLV form up to a moment-matching projection at
sibling-merge points, and finally state the concrete update rule.

\emph{Setup.} On a rooted tree $\mathcal{T}$ with the compound-CTMC
generative model (root state drawn from
$\rho_\theta \prod_n \eqm^{(\arch[c_n, \theta])}(x_n)$; each branch
of length $\tau$ transports $(x, \theta)$ under the joint kernel of
\eqref{eq:dynfield-transition}), define the standard Felsenstein
partials
\begin{align}
F_v(x_v, \theta_v) &\;\equiv\; P(\text{obs at leaves under } v \mid x_v, \theta_v)
    &&\text{(upward)} \\
D_v(x_v, \theta_v) &\;\equiv\; P(\text{obs at leaves outside subtree at } v \mid x_v, \theta_v)
    &&\text{(downward)}.
\end{align}
Posterior at $v$:
\begin{equation}
q_v(x_v, \theta_v) \;=\;
    \frac{ P^{\mathrm{prior}}_v(x_v, \theta_v) \cdot F_v(x_v, \theta_v) \cdot D_v(x_v, \theta_v)}
             {P(\text{observations})},
    \qquad
    P^{\mathrm{prior}}_v = \rho_{\theta_v} \prod_n \eqm^{(\arch[c_n, \theta_v])}(x_{v, n}),
    \label{eq:arch-elbo-post}
\end{equation}
since the compound CTMC's joint stationary factorises as
$\rho \cdot \prod_n \eqm^{(\arch)}$ at every node (a consequence of
the F81-on-DP field being at stationary throughout the tree).

\emph{Exact downward recursion.} At the root $r$, no observations lie
outside the subtree, so $D_r \equiv 1$. For a non-root $v$ with parent
$p$ and siblings $\{s_i\}$ (children of $p$ other than $v$),
condition on parent state and sum:
\begin{equation}
D_v(x_v, \theta_v) \;=\;
    \sum_{x_p, \theta_p}
        K(x_v, \theta_v \mid x_p, \theta_p; \tau_v) \;
        D_p(x_p, \theta_p) \;
        \prod_i U_{s_i \to p}(x_p, \theta_p),
    \label{eq:arch-elbo-D-exact}
\end{equation}
where $U_{s_i \to p}(x_p, \theta_p) := \sum_{x_{s_i}, \theta_{s_i}}
K(x_{s_i}, \theta_{s_i} \mid x_p, \theta_p; \tau_{s_i}) \, F_{s_i}$
is precisely the upward message the forward sweep already computes at
$p$. Write
\begin{equation}
M_v(x_p, \theta_p) \;\equiv\; D_p(x_p, \theta_p) \prod_i U_{s_i \to p}(x_p, \theta_p)
    \label{eq:arch-elbo-Mv}
\end{equation}
for the aggregated ``outside'' partial at $v$'s parent; then
\eqref{eq:arch-elbo-D-exact} is a single edge propagation of $M_v$
through the branch $p \to v$, summed on the source side (parent) rather
than the destination side (as in the forward sweep).

\emph{Kernel decomposition.} Under Interp 2, the joint kernel factors
into no-jump and $\geq\!1$-jump components (equations
\eqref{eq:dynfield-transition}--\eqref{eq:dynfield-weights}):
\begin{equation}
K(x_v, \theta_v \mid x_p, \theta_p; \tau) =
    \underbrace{\beta(\theta_p, \tau) \delta_{\theta_v = \theta_p}
                    \prod_n P^{(c_n, \theta_p)}(x_{p, n}, x_{v, n}; \tau)}_{\text{no-jump}}
    + \underbrace{W(\theta_p, \theta_v; \tau) \prod_n \eqm^{(\arch[c_n, \theta_v])}(x_{v, n})}_{\text{jump}},
    \label{eq:arch-elbo-K-decomp}
\end{equation}
with $\beta(\theta_p, \tau) = e^{-\rho_{\text{chain}}(1 - \rho_{\theta_p}) \tau}$
and $W(\theta_p, \theta_v; \tau) = K_{\mathrm{field}}(\theta_v \mid \theta_p; \tau)
- \delta_{\theta_v = \theta_p}\, \beta(\theta_p, \tau)$.

Substituting into \eqref{eq:arch-elbo-D-exact} and evaluating the
$\delta_{\theta_v = \theta_p}$ constraint in the no-jump branch:
\begin{align}
D_v(x_v, \theta_v)
    &\;=\;
    \beta(\theta_v, \tau)
        \sum_{x_p} \Bigl[\prod_n P^{(c_n, \theta_v)}(x_{p, n}, x_{v, n}; \tau)\Bigr]
        M_v(x_p, \theta_v)
    \nonumber \\
    &\quad
    + \Bigl[\prod_n \eqm^{(\arch[c_n, \theta_v])}(x_{v, n})\Bigr]
        \sum_{\theta_p} W(\theta_p, \theta_v; \tau)
        \sum_{x_p} M_v(x_p, \theta_p).
    \label{eq:arch-elbo-D-two-terms}
\end{align}

The first term (no-jump) is a per-site product of transposed
substitution kernels applied to $M_v$'s rank-1 content at the diagonal
$\theta_p = \theta_v$; the second term (jump) is a scalar in $x_v$
times the per-site stationary product at $\theta_v$. This is the same
(rank-1 + scalar) template as the forward message, but with two
critical differences:
\begin{enumerate}[label=(\alph*),leftmargin=1.6em]
\item Kernel direction: $\beta$ is evaluated at $\theta_v$ (destination
    in the downward direction) via the $\delta$-collapse, whereas in
    the forward pass $\beta$ is evaluated at the parent (also the
    ``destination'' from the upward perspective). Under reversibility
    of the field CTMC these two conventions differ by no factor;
    $\beta(\theta_v)$ is the correct multiplier in both directions
    once labels are matched to the resulting message's source axis.
\item Jump-weight orientation: $W(\theta_p, \theta_v)$, not
    $W(\theta_v, \theta_p)$ --- the summation variable is the parent's
    field, i.e., the \emph{source} of the field's forward transition.
    Using detailed balance $\rho_{\theta_p} W(\theta_p, \theta_v) =
    \rho_{\theta_v} W(\theta_v, \theta_p)$ (for $\theta_p \neq
    \theta_v$; equal on the diagonal), we may rewrite
    \begin{equation}
      \sum_{\theta_p} W(\theta_p, \theta_v; \tau)\, \bar M_v(\theta_p)
      \;=\;
      \frac{1}{\rho_{\theta_v}}
      \sum_{\theta_p} W(\theta_v, \theta_p; \tau)\,
      \bigl[\rho_{\theta_p} \bar M_v(\theta_p)\bigr],
      \label{eq:arch-elbo-W-flip}
    \end{equation}
    where $\bar M_v(\theta_p) := \sum_{x_p} M_v(x_p, \theta_p)$ is the
    theta-slice mass of $M_v$. Identity \eqref{eq:arch-elbo-W-flip}
    reduces the backward jump-weight sum to the same $W(\theta_v,
    \theta_p)$-structure as the forward mm-edge, provided we replace
    $\bar M_v(\theta_p) \mapsto \rho_{\theta_p} \bar M_v(\theta_p)$ on
    the input and divide by $\rho_{\theta_v}$ on the output.
\end{enumerate}

\emph{Backward propagation in the moment-matching family.} Suppose
$M_v$ is in the (rank-1 + scalar) form
\begin{equation}
M_v(x_p, \theta_p) = r_M(\theta_p) \prod_n A_M(x_{p, n} \mid \theta_p)
    + s_M(\theta_p) \prod_n \eqm^{(\arch[c_n, \theta_p])}(x_{p, n}),
    \label{eq:arch-elbo-M-form}
\end{equation}
with the normalisation
$\langle \eqm^{(c_n, \theta)}, A_M(\cdot | \theta) \rangle = 1$
inherited from the forward sweep. Substituting into
\eqref{eq:arch-elbo-D-two-terms} and using the per-site GTR
reversibility identity
$\sum_{y} P^{(c_n, \theta)}(y, x; \tau)\, \eqm^{(c_n, \theta)}(y)
= \eqm^{(c_n, \theta)}(x)$
(direct consequence of detailed balance on the substitution chain),
the two terms of \eqref{eq:arch-elbo-D-two-terms} become:
\begin{align}
D_v^{\text{no-jump}}(x_v, \theta_v)
    &\;=\;
    \beta(\theta_v, \tau)\,
    \Bigl[
      r_M(\theta_v)
      \prod_n \sum_y P^{(c_n, \theta_v)}(y, x_{v, n}; \tau)\, A_M(y \mid \theta_v)
      \nonumber \\
    &\qquad\qquad
      + s_M(\theta_v)
      \prod_n \eqm^{(\arch[c_n, \theta_v])}(x_{v, n})
    \Bigr], \\[0.4em]
D_v^{\text{jump}}(x_v, \theta_v)
    &\;=\;
    \Bigl[\prod_n \eqm^{(\arch[c_n, \theta_v])}(x_{v, n})\Bigr]
    \sum_{\theta_p} W(\theta_p, \theta_v; \tau)\,\bar M_v(\theta_p).
\end{align}

Combining, $D_v$ has the (rank-1 + scalar) form:
\begin{equation}
D_v(x_v, \theta_v) = r_D(\theta_v) \prod_n A_D(x_{v, n} \mid \theta_v)
    + s_D(\theta_v) \prod_n \eqm^{(\arch[c_n, \theta_v])}(x_{v, n}),
    \label{eq:arch-elbo-D-form}
\end{equation}
with
\begin{align}
r_D(\theta_v) &= \beta(\theta_v, \tau)\, r_M(\theta_v), \\
A_D(x \mid \theta_v)_n &= \sum_y P^{(c_n, \theta_v)}(y, x; \tau)\, A_M(y \mid \theta_v)
    \qquad\text{(per site $n$; note transposed convention)}, \\
s_D(\theta_v) &= \beta(\theta_v, \tau)\, s_M(\theta_v)
    + \sum_{\theta_p} W(\theta_p, \theta_v; \tau)\,\bar M_v(\theta_p).
    \label{eq:arch-elbo-sD}
\end{align}

The three closed-form updates constitute the backward analog of the
forward \texttt{mm\_edge}. Compared to the forward propagation:
\begin{itemize}[leftmargin=1.6em]
\item The rank-1 magnitude $r_M(\theta_v) \to r_D(\theta_v)$ transforms
    identically ($\beta$-scaling at the diagonal $\theta_p = \theta_v$).
\item The per-site rank-1 marginal $A_M \to A_D$ is propagated by the
    \emph{transposed} substitution kernel $y \to x$ rather than
    $x \to y$; algebraically, $A_D = P^{(c_n, \theta_v)} \cdot A_M$ if
    we treat $A$ as a column vector, whereas the forward update was
    $A_v = (P^{(c_n, \theta_v)})^{\!\top} \cdot A_c$. Under GTR
    reversibility these two operations are related by
    $A_D = D_{\pi}^{-1} A_v^{\text{fwd}} D_\pi$ where $D_\pi = \diag(\eqm)$,
    but the two are numerically distinct in general.
\item The scalar $s_M(\theta_v) \to s_D(\theta_v)$ picks up a
    $\rho$-scaled jump-weight sum with source $\theta_p$
    \eqref{eq:arch-elbo-W-flip}, differing from the forward
    $\sum_{\theta_u} W(\theta_v, \theta_u) \bar F(\theta_u)$ by the
    $\rho$-swap of source and destination in $W$.
\end{itemize}
The normalisation invariant $\langle \eqm, A_D \rangle = 1$ is
preserved by the transposed propagation, by the same argument as the
forward: $\langle \eqm, P^{(c_n, \theta_v)} A_M \rangle_x =
\langle A_M, P^{(c_n, \theta_v)\!\top} \eqm \rangle_y =
\langle A_M, \eqm \rangle = 1$ using
$(P^{(c_n, \theta_v)})^{\!\top} \eqm = \eqm$ (right-invariance of the
stationary under reversibility).

\emph{Sibling merge (moment-matching at the aggregator $M_v$).} The
aggregate $M_v = D_p \cdot \prod_i U_{s_i \to p}$ is a product of
several (rank-1 + scalar) forms and must be projected back onto the
family before the edge propagation above can be applied. Under the
moment-matching rule of \texttt{mm\_combine}, each pairwise product
of terms $[r_1 R_1 + s_1] \cdot [r_2 R_2 + s_2]$ decomposes into
$r_1 r_2 R_1 R_2 + r_1 s_2 R_1 + s_1 r_2 R_2 + s_1 s_2$; the scalar
component $s_1 s_2$ is kept as the projected scalar, and the three
rank-1 residuals are merged into a single per-site marginal by
matching the per-site $\langle \eqm, \cdot \rangle$-moments. This is
the same projection used in the forward sweep at internal-node
combines; it is exact when there is only one contributing sibling
(so $M_v = D_p \cdot U_{s \to p}$ is a single pairwise product) and
introduces a controlled bias for higher-arity nodes, matching the
forward's cost/quality tradeoff.

\emph{Root and leaf boundary.} At the root, $M_v = D_p \cdot U_{s
\to p}$ reduces to $\prod_{i \neq v} U_{s_i \to \text{root}}$ (no
$D_p$ factor); the root's ``prior'' contribution $\rho \cdot \prod
\eqm$ enters only when computing the posterior via \eqref{eq:arch-elbo-post}
so that $D_r \equiv 1$ needs no representation. At a leaf $\ell$,
$F_\ell(x, \theta) = \delta(x = x_\ell^{\mathrm{obs}})$ and
$D_\ell(x, \theta)$ is computed by exactly one backward mm-edge from
its parent; the leaf posterior over $\theta_\ell$ then reduces to
$q_\ell(\theta) \propto \rho_\theta \, \eqm^{(\arch[\ldots, \theta])}(x_\ell^{\mathrm{obs}}) \cdot D_\ell(x_\ell^{\mathrm{obs}}, \theta)$.

\emph{Coordinate-ascent convergence.} A full backward sweep after each
forward sweep gives one ELBO coordinate-ascent iteration; the ELBO
\eqref{eq:arch-elbo-total} is non-decreasing per half-sweep and
converges in $10$--$50$ full sweeps in practice (matched to
\citep{evolmoves}). At convergence, the forward and backward CLVs at
every node are jointly consistent with the marginal posterior
\eqref{eq:arch-elbo-post}, and \eqref{eq:arch-elbo-D-exact} is
satisfied to floating-point precision on cherries and to the
moment-matching tolerance on deeper trees.

\paragraph{Bridge posterior on a branch.}
\label{par:arch-phylo-elbo-bridge}

For extraction of HR sufficient statistics we need not the marginal
posterior at a node but the \emph{joint posterior on a branch} ---
the pair distribution over $((x_v, \theta_v), (x_u, \theta_u))$ at
the two ends of the branch $v \to u$ of length $\tau_{vu}$:
\begin{equation}
q_{vu}(x_v, \theta_v, x_u, \theta_u)
    \;\propto\;
    P^{\mathrm{prior}}_v(x_v, \theta_v) \cdot D_v(x_v, \theta_v)
    \cdot K(x_u, \theta_u \mid x_v, \theta_v; \tau_{vu})
    \cdot F_u(x_u, \theta_u),
    \label{eq:arch-elbo-bridge}
\end{equation}
where $F_u$ and $D_v$ are the converged CLVs and the normalising
constant is the tree log-likelihood $P(\text{obs})$.

Substituting the two-term $K$ \eqref{eq:arch-elbo-K-decomp} and the
(rank-1 + scalar) forms of $F_u$ and $D_v$ into
\eqref{eq:arch-elbo-bridge}, and marginalising over positions
$(x_v, x_u)$ at all sites but one, each site's pair distribution
$q_{vu, n}(x_{v, n}, x_{u, n} \mid \theta_v, \theta_u)$ is a mixture
over $2 \times 2 = 4$ terms (no-jump vs.\ jump, tracking vs.\ regenerated
sources at each end), each of which factorises across sites given
$(\theta_v, \theta_u)$. The concrete formulas appear in
Algorithm~\ref{alg:arch-elbo-hr} below.

\paragraph{Field-jump-count decomposition of a branch.}
\label{par:arch-phylo-elbo-jumps}

The branch kernel $K(x_u, \theta_u \mid x_v, \theta_v; \tau)$ marginalises
over the entire compound trajectory on the branch, but the HR sufficient
statistics require attribution of dwell time and substitution counts
to specific archetypes, and archetypes change with $\theta$. We therefore
decompose the branch by the number $M$ of field jumps that occurred on
the interval $(0, \tau)$:
\begin{align}
K(x_u, \theta_u \mid x_v, \theta_v; \tau) \;=\;
    \underbrace{K^{(M=0)}(x_u \mid x_v, \theta_v; \tau)\,\delta_{\theta_u = \theta_v}}_{\text{no jump}}
    +\underbrace{\sum_{M \geq 1} K^{(M)}(x_u, \theta_u \mid x_v, \theta_v; \tau)}_{\text{$M \geq 1$ jumps}}.
\end{align}

\emph{$M = 0$: no field jump.} The path is a pure GTR evolution under
the fixed archetype $\arch[c_n, \theta_v]$ at each site over the whole
branch of length $\tau$. Probability of the joint (no-jump event and
observed residues):
\begin{equation}
K^{(M=0)}(x_u \mid x_v, \theta_v; \tau)
    \;=\; \beta(\theta_v, \tau) \prod_n P^{(c_n, \theta_v)}(x_{v, n}, x_{u, n}; \tau),
    \label{eq:arch-elbo-KM0}
\end{equation}
with $\beta(\theta_v, \tau) = e^{-\rho_{\text{chain}}(1 - \rho_{\theta_v})\tau}$.

\emph{$M \geq 1$: at least one field jump.} Under Interp 2 the
field-CTMC generator has state-dependent exit rate $\rho_{\text{chain}}
(1 - \rho_{\theta})$; conditional on any jump occurring, the destination
distribution is $\rho_{\theta_u} / (1 - \rho_{\theta_v})$ for
$\theta_u \neq \theta_v$. Aggregating over $M \geq 1$ we may write:
\begin{equation}
\sum_{M \geq 1} K^{(M)}(x_u, \theta_u \mid x_v, \theta_v; \tau)
    \;=\;
    W(\theta_v, \theta_u; \tau)
    \, \prod_n \eqm^{(\arch[c_n, \theta_u])}(x_{u, n}),
    \label{eq:arch-elbo-KM1}
\end{equation}
i.e., the ``regenerated'' component of the branch kernel. The
per-site product $\prod_n \eqm^{(\arch[c_n, \theta_u])}(x_{u, n})$
expresses the collapse assumption that after any field jump the
substitution CTMC is at stationary conditional on the current
archetype (the ``carry'' vs.\ ``resample'' choice of Section~\ref{par:arch-carry-resample}),
which under Interp 2 corresponds to a full re-equilibration per jump.

\emph{Case-conditional HR primitives.} On the branch, for each pair
$(\theta_v, \theta_u)$, the case-$M{=}0$ contribution to HR statistics
attributes all dwell time and transitions to archetype $\arch[c_n,
\theta_v]$ for site $n$. This is exactly the standard \emph{whole-branch
GTR bridge} HR pass at site $n$: given endpoint residues $(x_{v, n},
x_{u, n})$ and rate $Q^{(c_n, \theta_v)}$, the per-site sufficient
stats $(W_n^{(0)}, U_n^{(0)})$ are computed by the reversible-bridge
closed forms of Section~\ref{par:arch-hr}, with the ``segment start''
count $V_n^{(0)}$ receiving a $\delta_{a = x_{v, n}}$ contribution
(since the segment starts at the observed root-side residue).

The case-$M \geq 1$ contribution is more intricate: the branch is
partitioned into a first segment $[0, \tau_1]$ under archetype
$\arch[c_n, \theta_v]$, a middle chain of jumps at random times with
free intermediate archetypes drawn from
$\arch[c_n, \theta_{\mathrm{int}}]$ with $\theta_{\mathrm{int}} \sim
\rho$, and a last segment $[\tau_M, \tau]$ under archetype
$\arch[c_n, \theta_u]$. The whole-branch HR contributions decompose
as (dropping the site index $n$ for readability):
\begin{align}
V^{(M \geq 1)}_a
    &\;=\; \delta_{a = x_v}
    + \Bigl(\sum_{\theta_{\mathrm{int}}} \mathrm{Arr}(\theta_{\mathrm{int}}; \tau)\Bigr)
        \cdot \eqm^{(\arch[c_n, \theta_{\mathrm{int}}])}(a)
    + P_{\mathrm{avg}}^{\mathrm{last}}(x_u; a; \tau),
    \label{eq:arch-elbo-Vjump}\\
W^{(M \geq 1)}_i
    &\;=\; W_i^{\mathrm{first}}(x_v; \theta_v; \tau_1)
    + \mathbb{E}[\tau_{\mathrm{mid}}] \cdot \eqm^{(\arch[c_n, \theta_{\mathrm{int}}])}(i)
    + W_i^{\mathrm{last}}(x_u; \theta_u; \tau - \tau_M),
    \label{eq:arch-elbo-Wjump}\\
U^{(M \geq 1)}_{i, j}
    &\;=\; U_{i, j}^{\mathrm{first}}(x_v; \theta_v; \tau_1)
    + \mathbb{E}[\tau_{\mathrm{mid}}] \cdot \eqm^{(\arch[c_n, \theta_{\mathrm{int}}])}(i)
        \cdot S_{i, j}
        \cdot \eqm^{(\arch[c_n, \theta_{\mathrm{int}}])}(j)
    + U_{i, j}^{\mathrm{last}}(x_u; \theta_u; \tau - \tau_M),
    \label{eq:arch-elbo-Ujump}
\end{align}
with $\mathrm{Arr}$, $P_{\mathrm{avg}}^{\mathrm{last}}$, and the
$\mathrm{first}/\mathrm{last}$-segment HR primitives all computed by
the case-1 closed forms of Sections~\ref{par:arch-hr}--\ref{par:arch-hr-cluster}
(specifically \eqref{eq:arch-Wfirst-mix}--\eqref{eq:arch-Ulast-mix}
for $W^{\mathrm{first}, \mathrm{last}}, U^{\mathrm{first}, \mathrm{last}}$,
\eqref{eq:arch-Etau1-mix}--\eqref{eq:arch-EtauM-mix} for
$E[\tau_1], E[\tau - \tau_M]$, and \eqref{eq:arch-Nmid} for
$\mathbb{E}[\tau_{\mathrm{mid}}]$). The middle-segment attribution
uses the stationary of the intermediate archetype for both dwell
and transitions, weighted by the expected number of arrivals at that
$\theta_{\mathrm{int}}$.

The V-attribution term $P_{\mathrm{avg}}^{\mathrm{last}}(x_u; a; \tau)$
is the expected starting-residue distribution of the last segment,
computed by the ``transition-avg'' closed form
\eqref{eq:arch-Pres}: it is not $\delta_{a = x_u}$ because the last
segment starts \emph{internally} on the branch (not at the observed
$x_u$), so the segment-start residue is averaged over the
last-jump-time distribution.

\emph{Spectral acceleration.} The per-$(c, k)$ HR primitives
$W^{\mathrm{first}}$, $W^{\mathrm{last}}$, $U^{\mathrm{first}}$,
$U^{\mathrm{last}}$ have closed forms in the eigenspace of the GTR
generator: given the eigendecomposition
$Q^{(c_n, \arch[c_n, \theta])} =
D_{\eqm^{(c_n, \theta)}}^{1/2} U_{c_n, \theta} \Xi_{c_n, \theta}
U_{c_n, \theta}^{\top} D_{\eqm^{(c_n, \theta)}}^{-1/2}$ (Section~\ref{par:arch-hr}),
the per-segment integrals in \eqref{eq:arch-elbo-Wjump}--\eqref{eq:arch-elbo-Ujump}
factor into a sum over eigenpair-index pairs $(a, b)$ of the elementary
$\mathrm{H}$ and $\mathrm{T}$ integrals \eqref{eq:arch-H}--\eqref{eq:arch-T}.
The per-branch cost is thus $O(L^2 \cdot m \cdot A^2)$ (assembling the
per-site pair distribution) plus $O(L^2 \cdot m \cdot A)$ (spectral
integration), rather than the naive $O(A^{2m})$ or even $O(A^4)$
per-site cost that a direct enumeration of the $(x_v, x_u)$ pairs
would require. This mirrors the ``accumulate sums in eigenspace''
trick of the composite HR pass (Section~\ref{par:arch-hr-cluster}).

\paragraph{Obtaining HR sufficient statistics from the ELBO'd tree.}
\label{par:arch-phylo-elbo-hr}
The M-steps for the archetype parameters \eqref{eq:arch-Mstep-pi} and
$\rho_{\text{chain}}$ \eqref{eq:arch-Mstep-rho} consume aggregated HR
sufficient statistics
$\{V_{c, \theta, \cdot}, U_{c, \theta, \cdot, \cdot}, W_{c, \theta,
\cdot}\}$, $\sum_q N_\theta[q]$, $\sum_q T_q$. Under composite training
these are accumulated over the cluster HR pass of
Algorithm~\ref{alg:arch-hr-cluster}, one HR pass per cherry. Under
tree training the same statistics are accumulated one HR pass per
\emph{branch} of the tree, with the branch endpoints' $(\theta,
r_1, \ldots, r_{m})$ distributions read from the converged
variational marginal $q$.

For each branch $v \to u$ of length $\tau_{vu}$, given the converged
forward CLVs $F_v, F_u$ and backward CLVs $D_v, D_u$
(Section~\ref{par:arch-phylo-elbo-backward}):

\begin{enumerate}
\item Compute the joint pair posterior $q_{vu}(x_v, \theta_v, x_u,
    \theta_u)$ via \eqref{eq:arch-elbo-bridge}. Under the moment-matching
    family it factorises across sites $n$ conditional on $(\theta_v,
    \theta_u)$; extract the four-term per-site pair distribution
    $q_{vu, n}(x_{v, n}, x_{u, n} \mid \theta_v, \theta_u)$ as a
    mixture over $\{$no-jump/jump$\} \times \{$tracking$_v$/regenerated$_v\}
    \times \{$tracking$_u$/regenerated$_u\}$ (some cells vanish under
    $\delta_{\theta_v = \theta_u}$ on the no-jump branch), giving the
    field-pair weight
    \begin{equation}
      w^{vu}(\theta_v, \theta_u)
        \;=\;
        \sum_{x_v, x_u} q_{vu}(x_v, \theta_v, x_u, \theta_u).
    \label{eq:arch-elbo-wpair}
    \end{equation}
    The residue pair distribution conditional on $(\theta_v, \theta_u)$
    factorises across sites.

\item For each pair $(\theta_v, \theta_u)$ and each site $n$, compute
    the per-site HR contributions using the case-conditional primitives
    of Section~\ref{par:arch-phylo-elbo-jumps}:
    \begin{itemize}
    \item \emph{$\theta_v = \theta_u$ (no jump possible along with the
        $\delta$-collapse, weight $\beta \cdot q_{vu, n}$)}: apply the
        whole-branch GTR bridge HR primitive of Section~\ref{par:arch-hr}
        at archetype $\arch[c_n, \theta_v]$, endpoints
        $(x_{v, n}, x_{u, n})$ integrated against the per-site pair
        distribution.
    \item \emph{$M \geq 1$-jump contribution (any $\theta_v, \theta_u$)}:
        apply \eqref{eq:arch-elbo-Vjump}--\eqref{eq:arch-elbo-Ujump},
        integrated against the same per-site pair distribution but with
        the residue-marginal at $x_u$ taken from
        $\eqm^{(\arch[c_n, \theta_u])}$ (regenerated component).
    \end{itemize}

\item Weight the per-$(\theta_v, \theta_u)$ contributions by
    \eqref{eq:arch-elbo-wpair} and sum across all field-pair contexts,
    all sites $n$, and all branches, indexing into the archetype bins by
    $k = \arch[c_n, \theta]$ (with $\theta = \theta_v$ for first-segment
    attribution, $\theta_u$ for last-segment, and $\theta_{\mathrm{int}}$
    weighted by $\mathrm{Arr}(\theta_{\mathrm{int}})$ for middle-segment):
    \begin{align}
      V[k, a]
        &\;\mathrel{{+}{=}}\;
        \sum_{vu, n, \theta_v, \theta_u}
          w^{vu}(\theta_v, \theta_u)\,
          \bigl[\arch[c_n, \theta_v] = k\bigr] \,
          V^{(vu, \theta_v, \theta_u)}_{n, k, a}, \\
      U[k, i, j], W[k, i]
        &\;\text{analogously.}
    \end{align}
    The field-jump counts and time sums receive
    \begin{equation}
      \sum_q N_\theta[q] \;\mathrel{{+}{=}}\;
        \sum_{vu, \theta_v, \theta_u} w^{vu}(\theta_v, \theta_u)
        \,\mathbb{E}[N_\theta \mid \theta_v, \theta_u, \tau_{vu}],
      \qquad
      \sum_q T_q \;\mathrel{{+}{=}}\;
        \sum_{vu} \tau_{vu},
      \label{eq:arch-elbo-Nsum-Tsum}
    \end{equation}
    with $\mathbb{E}[N_\theta \mid \theta_v, \theta_u, \tau_{vu}]$ from
    the bridge-conditional field expectation \eqref{eq:arch-Nbridge}.
\end{enumerate}

Under the moment-matching family the per-branch cost is dominated by
the eigenspace expansion of the per-$(\theta_v, \theta_u)$ HR
primitives at each of $m$ sites: $O(L^2 \cdot m \cdot A^2)$ per
branch, matching the forward and backward mm-edge cost. The overall
per-iter cost per cluster is thus
$O(\text{edges} \cdot L^2 \cdot m \cdot A^2)$ times a small
coordinate-ascent factor.

The archetype $\eqm^{(k)}$ Newton MAP step \eqref{eq:arch-Mstep-pi} and
the $\rho_{\text{chain}}$ closed-form step \eqref{eq:arch-Mstep-rho}
are then unchanged: both consume the aggregated $(V, U, W, N_\theta,
T)$, whose \emph{source} has moved from independent cherries to
tree-branch expectations under the variational posterior. The M-step
targets an ELBO-consistent lower bound on the true tree likelihood
(rather than the composite likelihood), so the fixed points of the
resulting SEM loop are stationary points of the tree-ELBO objective.

\emph{Reduction to composite training.} If the tree is a forest of
independent depth-1 cherries with no shared internal ancestors, the
variational family \eqref{eq:arch-elbo-family} is exact (each cherry
runs its own two-leaf projection), and the branch-summed HR statistics
above reduce identically to the cherry HR statistics of
Algorithm~\ref{alg:arch-hr-cluster}. Composite training is therefore
the depth-1 special case of tree-ELBO training. Any Pfam MSA
represented as a star tree (a single internal ancestor with $Q$ leaf
cherries) already forbids the composite-likelihood freedom, since all
$Q$ cherries must share the same $\theta$ posterior at the root
node.

\emph{Class assignments.} The per-site class $c_n$ is fixed within a
tree-ELBO sweep, following the composite training regime
(Section~\ref{par:arch-algorithm}, sampled by the outer Gibbs sweep).
A mean-field per-site class marginal
$q_{n, v}(c_n) \in \Delta^{K_c - 1}$ could be added to the per-node
variational state at cost $O(K_c \cdot L \cdot m \cdot A^2)$
per edge, but at training corpus scale the sampled hard assignment
avoids the diffuse-marginal failure mode observed in strong-signal
regimes with $(c_i^*, c_j^*) \leftrightarrow (c_j^*, c_i^*)$
symmetry, and the class Gibbs is more cost-effective.

\emph{Trigger for tree-ELBO training.} The composite (cherry) regime
is adequate when (a) the corpus is presented as a set of aligned pairs
without tree structure, or (b) the composite-likelihood artifact
diagnostics (cross-cherry $\theta_P$ entropy vs.\ mean within-cherry
jump probability) show no artifact. When either condition fails --- as
in the training runs of Section~\ref{sec:experiments} where the
diagnostic reports $\geq 65\%$ of clusters in the artifact quadrant ---
tree-ELBO training becomes the preferred regime.

\paragraph{Rate heterogeneity: discretised-$\Gamma$ plus invariant mixture on
the field-flip rate.}
\label{par:arch-gamma-plus-I}
The base model of the previous paragraphs assumes a single global
$\rho_{\text{chain}}$ governs the field-flip rate for every cluster.
Adapting the classical Yang~\citep{yang1994maximum, yang1996among}
$+\Gamma+I$ rate-heterogeneity framework to the archetype trainer, we
introduce a per-cluster latent
$b_q \in \{\mathrm{inv}\} \cup \{1, 2, \dots, K_r\}$
that scales the flip rate:
\begin{equation}
  \rho^{(q)}_{\text{chain}} \;=\;
    \begin{cases}
      0, & b_q = \mathrm{inv},\\
      \rho_{\text{chain}} \cdot m_{b_q}, & b_q \in \{1, \dots, K_r\},
    \end{cases}
  \label{eq:arch-rate-het}
\end{equation}
where $\{m_1, \dots, m_{K_r}\}$ are the $K_r$ equal-mass quantile means
of a $\Gamma(\alpha_\gamma, \alpha_\gamma)$ distribution
(so $E[m] = 1$; $\alpha_\gamma$ controls the variance across bins). The
mixing weights are $\Pr[b_q = \mathrm{inv}] = p_{\text{inv}}$ and
$\Pr[b_q = k] = (1 - p_{\text{inv}})/K_r$ for the $\Gamma$ bins,
mirroring Yang's $+\Gamma+I$ default.

\emph{Interpretation.} The invariant bin ($b_q = \mathrm{inv}$) models
clusters whose field-selector trajectory is frozen for the entire
branch --- structurally invariant sites in Yang's original
substitution-rate application transfer directly to structurally
non-switching cluster contexts here. The $K_r$ $\Gamma$ bins give the
model flexibility to have both slow-mixing clusters (low $m$, mostly
one archetype for the branch) and fast-mixing clusters (high $m$,
walking through multiple archetypes) share a single $\rho_{\text{chain}}$
base rate, with the shape parameter $\alpha_\gamma$ learned from data.

\emph{Standard choice.} $K_r = 4$, following~\citet{yang1996among};
this and the single $p_{\text{inv}}$ / $\alpha_\gamma$ scalar pair adds
just three global parameters to the model, without altering the
archetype vocabulary $\{\eqm^{(k)}\}$ or field-state marginals $\rho$.
The version implemented in this work uses $K_r = 4$ Gamma bins plus one
invariant bin ($4\Gamma+I$).

\emph{E-step.} For each cluster $q$, evaluate the marginal likelihood
under each of the $K_r + 1$ candidate rates (invariant + Gamma bins),
either by running the HR pass $K_r + 1$ times with different
$\rho^{(q)}_{\text{chain}}$ or by conditioning $\rho_{\text{chain}}$
inside a vectorised HR call. The soft posterior is
\begin{equation}
  q_q(b) \;\propto\; \Pr[b_q = b] \cdot P(\text{obs}_q \mid b_q = b,
    \eqm, \arch, \rho, \rho_{\text{chain}}),
  \label{eq:arch-rate-Estep}
\end{equation}
computable in closed form from the HR pass at each rate.

\emph{M-step.} The HR sufficient statistics
$\{V_{c, \theta, \cdot}, U_{c, \theta, \cdot, \cdot}, W_{c, \theta, \cdot}\}$
aggregate over clusters with weights $q_q(b)$ (equivalently: each
cluster contributes $K_r + 1$ weighted copies of its stats). The
Newton update~\eqref{eq:arch-Mstep-pi} on $\eqm^{(k)}$ proceeds
unchanged. The Gamma-posterior $\rho_{\text{chain}}$
update~\eqref{eq:arch-Mstep-rho} becomes
\begin{equation}
  \rho_{\text{chain}} \;\leftarrow\;
    \frac{\alpha_\rho - 1 + \sum_q \sum_k q_q(k)\, N_\theta[q, k]}
         {\beta_\rho + \sum_q \sum_k q_q(k)\, m_k\, T_q},
  \label{eq:arch-Mstep-rho-het}
\end{equation}
where the sum over $k$ omits the invariant bin (which contributes zero
to both numerator and denominator). The $p_{\text{inv}}$ update is the
Beta-conjugate closed form
$p_{\text{inv}} \leftarrow (\sum_q q_q(\mathrm{inv}) + a_p - 1) /
(Q + a_p + b_p - 2)$
under a $\mathrm{Beta}(a_p, b_p)$ prior; $\alpha_\gamma$ is updated by
Newton on the marginal log-likelihood, given the fixed bin weights and
$\{N_\theta[q, k], T_q\}$ sufficient statistics.

\emph{Expected impact.} Structural specialisers under the base model
(classes whose $\arch(c, \theta)$ is constant across $\theta$) currently
occupy the low-$\rho_{\text{chain}}$ tail; under
\eqref{eq:arch-rate-het} they will preferentially assign to the
invariant or lowest-$m$ bin. Dynamic mediators (classes that visit
multiple archetypes across $\theta$) will migrate to the higher-$m$
bins, sharpening both the archetype-usage separation and the
compensatory-flip signal (Section~\ref{par:arch-algorithm}).

\paragraph{Per-site $\Gamma+I$ rate heterogeneity on the substitution
rate.}
\label{par:arch-gamma-plus-I-persite}
Alongside the per-cluster $\Gamma+I$ mixture on the field-flip rate,
we add a second, independent $\Gamma+I$ mixture at the per-site level
that scales the archetype substitution rate. This mirrors Yang's
original~\citep{yang1994maximum} +$\Gamma$+I on the GTR substitution
matrix and is the standard treatment of positional rate heterogeneity
in phylogenetics.
%
For each column $s$ in the corpus, introduce a latent
$b_s \in \{\mathrm{inv}\} \cup \{1, \dots, K_r^{\mathrm{site}}\}$ that
scales the site's substitution rate:
\begin{equation}
  Q^{(k, s)} \;=\;
    \begin{cases}
      \mathbf{0}, & b_s = \mathrm{inv},\\
      m^{\mathrm{site}}_{b_s} \cdot Q^{(k)}, & b_s \in \{1, \dots, K_r^{\mathrm{site}}\},
    \end{cases}
  \label{eq:arch-persite-rate-het}
\end{equation}
where $k = \arch(c_s, \theta_s)$ is the archetype visible at column
$s$ when the field selector is $\theta_s$, and
$m^{\mathrm{site}}_{1}, \dots, m^{\mathrm{site}}_{K_r^{\mathrm{site}}}$
are the quantile means of a
$\Gamma(\alpha^{\mathrm{site}}_\gamma, \alpha^{\mathrm{site}}_\gamma)$
distribution (mean~$1$). Bin priors are
$\Pr[b_s = \mathrm{inv}] = p^{\mathrm{site}}_{\mathrm{inv}}$ and
$\Pr[b_s = k] = (1 - p^{\mathrm{site}}_{\mathrm{inv}})/K_r^{\mathrm{site}}$.
%
\emph{Composition.} The per-site and per-cluster mixtures act on
different rate constants (site-substitution versus field-flip) and
are prior-independent; the effective per-(cluster, site) rate is
$m^{\mathrm{site}}_{b_s} \cdot Q^{(k)}$ for the residue CTMC and
$\rho^{(q)}_{\text{chain}}$ for the field CTMC.
%
\emph{Invariant behaviour (witness-site convention).} When
$b_s = \mathrm{inv}$, the site has no continuous substitution
($Q^{(k, s)} = \mathbf{0}$), so between field jumps the residue is
preserved: the Case~0 (no-jump) emission collapses to
$[\![ Y_s = X_s ]\!]$ via the $m^{\mathrm{site}}_{b_s} \xi^{(k)} = 0$
eigenvalue limit. \emph{At} each field jump, however, the residue
resamples from the stationary at the new archetype
$\eqm^{(k(c_s, \theta_Y))}$, exactly like a variant-rate site under
Interpretation~1. The $M{\geq}1$ (jump) emission is therefore the
same for all bins: $\eqm^{(k(c_s, \theta_Y))}_{Y_s}$.
%
This is a deliberate departure from Yang's~\citeyearpar{yang1994maximum}
original ``$+I$ = no evolution at all'' semantics. Two motivations:
\emph{(i) covariation-signal clarity} --- for an invariant-rate site
whose observed residue nonetheless changes between $X$ and $Y$, the
only remaining explanation is a field jump; if multiple invariant
sites in the same cluster co-change, that is unambiguous evidence of
a coordinated field transition (they act as ``witness sites'' for
the field). Under the old ``never changes'' convention this signal
was folded into the likelihood-zero regime and lost; and
\emph{(ii) mechanism separation} --- 100\%-conserved sites in the
witness-site convention are captured through a distinct mechanism:
a singleton cluster whose per-cluster rate bin
(Eq.~\ref{eq:arch-rate-het}) is invariant, so its field never jumps.
This factorises "site rate" and "cluster field-flip rate" instead of
conflating them.
%
Reversibility is preserved: an invariant bin is trivially reversible
w.r.t. any stationary, and the Case~$\geq$1 formula is the same as
for variant bins.

\emph{Efficient inference by rate-eigenvalue scaling.} Since
$Q^{(k, s)}$ shares the eigenvectors $U^{(k)}, D^{(k)}$ of $Q^{(k)}$
and differs only by scaling the eigenvalues by $m^{\mathrm{site}}_{b_s}$,
every closed-form HR primitive (bridge \eqref{eq:hr-bridge},
free-end \eqref{eq:hr-freeend}, case-$1$ transition average
\eqref{eq:hr-caseavg}) can be evaluated at any bin by substituting
$\xi^{(k)} \mapsto m^{\mathrm{site}}_{b_s} \cdot \xi^{(k)}$
inside the per-site call. No re-diagonalisation is required; the
$K_a$ archetype eigendecompositions computed once per outer iter
(cached until the next $\arch$~M-step) are reused across all
$(K_r^{\mathrm{site}} + 1)$ bins and all sites. Numerically this makes
per-site bin evaluation $O(A^2)$ per (bin, site), a factor of $A$
faster than a naive per-bin re-diagonalisation.

\emph{Per-site update strategy.} The per-site bin
$b_s$ is treated as persistent state; the HR pass applies each
site's rate multiplier $m^{\mathrm{site}}_{b_s}$ to $\xi^{(k)}$ at
every iteration without additional per-bin evaluation cost. Two
update regimes are supported (controlled by the trainer
\texttt{--site-resample-every R} flag):
\begin{enumerate}
  \item \emph{$\Gamma$-only init, no resample ($R = 0$).} Bins are drawn
    uniformly over the $\Gamma$ quantile bins ($1 \ldots K_r^{\mathrm{site}}$)
    at training start and held fixed. The invariant bin ($b_s = \mathrm{inv}$)
    is \emph{not} used at initialisation: assigning bin~0 to a column
    whose observed $(X, Y)$ is variant ($X \neq Y$ at any cherry)
    zeroes the Case~0 factor at that cluster on the very first HR
    pass; the resulting drop in Case-0 log-likelihood is severe
    enough (though not $-\infty$: Case~$\geq 1$ still contributes
    $\eqm^{(k(c, \theta_Y))}_Y > 0$) to make the atom step at iter~1
    diverge in practice. A safe uniform $\Gamma$ init sidesteps this.
  \item \emph{$\Gamma$-init + infrequent Gibbs ($R > 0$).} Same
    $\Gamma$-only init; then every $R$ outer iters, bin
    $b_s$ is refreshed by a partial Gibbs step using the Rao-Black\-well\-ised
    conditional (Eq.~\ref{eq:arch-persite-Gibbs}). The Gibbs
    conditional evaluates data-conditional per-site log-likelihoods
    over all $K_r^{\mathrm{site}} + 1$ bins and downweights bin~0 on
    variant columns (where Case~0 pays a large penalty) while
    upweighting bin~0 on invariant-observation columns in slow
    clusters (where Case~0~$= 1$ dominates). Bin~0 is thus
    introduced only where the data supports it, avoiding the
    initialisation blow-up of a naïve $+\Gamma+I$ prior draw.
\end{enumerate}

\emph{Historical note on the invariant-bin encoding.} An earlier
version of this section used the strict Yang~'94 semantics for
invariance (residue preserved in \emph{both} no-jump and jump cases).
That model class made a per-cherry Gibbs step unsafe: a column whose
observed $(X, Y)$ was $X = Y$ in one cherry and $X \neq Y$ in another
could not be assigned $b_s = \mathrm{inv}$ without zeroing the cluster
likelihood on the second cherry. The witness-site convention above
removes this obstruction (invariant bin now emits
$\eqm^{(k(c, \theta_Y))}_Y > 0$ in the jump case) and enables
$+\Gamma + I$ under the per-cherry Gibbs directly. The pre-witness-site
$\Gamma$-only sub-model remains available as an A/B baseline through
the \texttt{allow\_invariant=False} option of the Gibbs sampler.
%
The Gibbs conditional is Rao-Blackwellised over the field
trajectory:
\begin{equation}
  \tilde q_s(b) \;\propto\; \Pr[b_s = b] \cdot
    \!\!\sum_{\theta_X, \theta_Y}\!\!
       w^{(q)}_{\theta_X,\theta_Y}
       \cdot L^{(b)}_{\text{site}}(s; \theta_X, \theta_Y),
  \label{eq:arch-persite-Gibbs}
\end{equation}
where $q$ is the cluster containing $s$,
$w^{(q)}_{\theta_X, \theta_Y}$ is the pair weight
$\rho(\theta_X)\,P_{\theta_X \to \theta_Y}(\tau_q)$ produced by the
HR pass, and
\[
  L^{(b)}_{\text{site}}(s; \theta_X, \theta_Y) =
    \begin{cases}
      \mathbb{1}[X_s = Y_s], & b = \mathrm{inv},\\
      P^{k_s, m^{\mathrm{site}}_b}_{X_s Y_s}(\tau_q),
        & b = 1..K_r^{\mathrm{site}} \text{ and } \theta_X = \theta_Y,\\
      \eqm^{(k_s(\theta_Y))}_{Y_s},
        & b = 1..K_r^{\mathrm{site}} \text{ and } \theta_X \neq \theta_Y,
    \end{cases}
\]
with $k_s(\theta) = \arch(c_s, \theta)$. This
approximation ignores between-site coupling under the field
posterior (which is weak under Interpretation~1: sites are
conditionally independent given the field trajectory, and
between-site coupling only arises through the pair-averaged
posterior over $(\theta_X, \theta_Y)$).

\emph{M-step for $p^{\mathrm{site}}_{\mathrm{inv}}$,
$\alpha^{\mathrm{site}}_\gamma$.} Under a $\mathrm{Beta}(a_p, b_p)$
prior:
$p^{\mathrm{site}}_{\mathrm{inv}} \leftarrow
(\sum_s \mathbb{1}[b_s = \mathrm{inv}] + a_p - 1) /
(S_{\text{total}} + a_p + b_p - 2)$,
where $S_{\text{total}}$ is the total number of columns across the
corpus. $\alpha^{\mathrm{site}}_\gamma$ is held fixed at its
initialisation (Yang's~\citeyearpar{yang1996among} original
recommendation; can be swept externally).

\emph{Summary of cost.} Marginalising per-site bins at every outer
iter would multiply HR-pass work by $(K_r^{\mathrm{site}} + 1)$; the
persistent-plus-infrequent-Gibbs scheme reduces this to
$1 + K_r^{\mathrm{site}}/R \approx 1.8$~× amortised for
$K_r^{\mathrm{site}} = 4$, $R = 5$, giving essentially the runtime of
the base model plus a constant every $R$ iters for the resample.

\subsubsection{Variational ELBO for general phylogenies}
\label{sec:variational-suppl}

\Cref{prop:hier-fels} is tight at the cap-2 cluster geometry that
TKF-DP training uses: the per-node conditional likelihood at a
cherry parent has the closed-form decomposition
\begin{equation}
  F^v(x_v,\theta_v) = \beta_v(\theta_v)\,\prod_{n=1}^{m}
                       \alpha_n^v(x_{v,n},\theta_v) + \gamma_v(\theta_v)
  \label{eq:rank1plusscalar}
\end{equation}
--- a per-site rank-$1$ term (the cumulative no-event branch) plus
a residue-constant scalar (the regenerated branch); per-edge
propagation preserves this form. On phylogenies of depth $>1$, the
combine at an internal node $v$ with children $u_1,u_2$ multiplies
the rank-$1$ count, $K_v = K_{u_1}K_{u_2} + K_{u_1} + K_{u_2}$,
giving cost super-exponential in tree depth. A structured
mean-field variational restriction recovers linear cost while
remaining exact at cap-2.

\paragraph{Variational family.} At every tree node $v$ restrict the
joint $(x_v,\theta_v)$ marginal to the per-$\theta$ mixture of a
per-site factored marginal and the dynfield stationary:
\begin{equation}
  q_v(x_v,\theta_v) = q_v(\theta_v)\,\Big[(1-\lambda_v(\theta_v))
                            \prod_{n=1}^{m} p_n^v(x_{v,n}\!\mid\theta_v)
                       + \lambda_v(\theta_v)
                            \prod_{n=1}^{m} \eqm^{(c_n,\theta_v)}(x_{v,n})\Big].
  \label{eq:varform}
\end{equation}
Per-node variational parameters: $q_v(\theta_v)\in\Delta^{L-1}$,
$\lambda_v(\theta_v)\in[0,1]^L$, and $p_n^v(\cdot\!\mid\theta_v)\in
\Delta^{k-1}$ for $n=1,\dots,m$, total $O(L\cdot m\cdot k + 2L)$
per node. The interpretation of $\lambda_v(\theta_v)$ is the
posterior probability that at least one field event has occurred
on the path from $v$ back to the last ``tracking'' ancestor; in the
regenerated component residues sit at the model's joint stationary.

\paragraph{ELBO.}
\begin{align}
  \mathcal{L}(q) &= \!\sum_{\text{leaves } v}\!\mathbb{E}_{q_v}[\log p(\text{obs}_v\!\mid\!x_v,\theta_v)]
                    \nonumber \\
                 &\quad + \!\sum_{\text{edges } u\to v}\!\mathbb{E}_{q_v,q_u}[\log K(\phi_v,\phi_u;\tau)]
                    - \!\sum_{\text{nodes } v}\!\mathbb{E}_{q_v}[\log q_v(\phi_v)]
                    + \log Z_{\text{prior}}.
  \label{eq:elbo}
\end{align}
Each term has a closed form under \eqref{eq:varform}: leaf-emission
and node-entropy at $O(L\cdot m\cdot k)$ (the mixture entropy uses
the standard component-wise upper bound, tight when $\lambda\in\{0,1\}$);
per-edge cross-entropy at $O(L^2\cdot m\cdot k^2)$ with the per-site
$\Q$-evolution and field-transition kernel both factored over sites
given the endpoint fields $(\theta_v,\theta_u)$.

\paragraph{Per-edge propagation.} For each parent-child pair
$u\to v$ of length $\tau$, the exact propagation
\eqref{eq:rank1plusscalar} of one $(\text{rank-}1+\text{scalar})$
form to another is closed form; under the variational family the
propagation simplifies to per-site $\Q$-matrix actions on
$p_n^u(\cdot\!\mid\theta)$ followed by a renormalisation. Cost per
edge: $O(m\cdot L\cdot k^2 + L^2)$.

\paragraph{Internal-node projection (M-projection: moment matching).}
At an internal node $v$ with children $u_1,u_2$, the exact
multiplicative combine of two $(\text{rank-}1+\text{scalar})$ child
contributions yields four components ($3$ rank-$1$ $+$ $1$ scalar).
The variational M-projection restores the form \eqref{eq:varform}
as the $\KL(p_v\,\|\,q_v)$ minimiser:
\begin{itemize}\itemsep0pt
  \item $q_v(\theta_v)$ to the $\theta_v$-marginal of the exact
        four-term product;
  \item $p_n^v(x_{v,n}\!\mid\theta_v)$ to the per-site, per-$\theta_v$
        residue marginal of the exact product;
  \item $\lambda_v(\theta_v)$ to the proportion of total mass in
        terms involving at least one child's scalar component.
\end{itemize}
All updates are closed form; $O(m\cdot L\cdot k)$ per binary combine.
$k$-ary internal nodes are handled as $k-1$ sequential binary
combines with projection after each.

\paragraph{Internal-node projection (I-projection: KL minimisation).}
M-projection produces unbiased marginals but can systematically
overestimate $\lambda_v(\theta_v)$ when the exact $p_v$ has strong
residue-level structure (matching scalar mass via mass proportion
alone ignores the residue-level information that distinguishes the
rank-$1$ mode from the regenerated mode), loosening the ELBO and
risking pathological over-clustering bias downstream. The tighter
alternative is I-projection: minimise $\KL(q_v\,\|\,p_v)$ directly.
Writing
$\mathcal{T}(x_v\!\mid\theta_v) := (1-\lambda_v(\theta_v))\prod_n p_n^v(x_{v,n}\!\mid\theta_v)
                                  + \lambda_v(\theta_v)\prod_n \eqm^{(c_n,\theta_v)}(x_{v,n})$
(so $q_v(x_v,\theta_v)=q_v(\theta_v)\mathcal{T}(x_v\!\mid\theta_v)$),
the ELBO at node $v$ decomposes as
\begin{equation}
  \mathrm{ELBO}_v = \sum_{\theta_v} q_v(\theta_v)\,\big[\,\mathbb{E}_{\mathcal{T}(\cdot|\theta_v)}[\log p_v(\cdot,\theta_v)] - \log q_v(\theta_v) + \mathrm{H}[\mathcal{T}(\cdot\!\mid\theta_v)]\,\big].
  \label{eq:elbo-node}
\end{equation}
The per-parameter coordinate-ascent gradients (with simplex
multipliers $\mu$ for $\sum q_v = \sum p_n^v = 1$ suppressed):
\begin{align}
  \frac{\partial\,\mathrm{ELBO}_v}{\partial\,q_v(\theta_v)}
    &= \mathbb{E}_{\mathcal{T}(\cdot|\theta_v)}[\log p_v(\cdot,\theta_v)] - \log q_v(\theta_v) - 1 + \mathrm{H}[\mathcal{T}(\cdot\!\mid\theta_v)],
    \label{eq:dql-dq}\\
  \frac{\partial\,\mathrm{ELBO}_v}{\partial\,\lambda_v(\theta_v)}
    &= q_v(\theta_v)\,\big[\,\mathbb{E}_{\boldsymbol{\eqm}^{\theta_v}}[\log p_v/\mathcal{T}](\cdot,\theta_v) - \mathbb{E}_{\mathbf{p}^v}[\log p_v/\mathcal{T}](\cdot,\theta_v)\,\big],
    \label{eq:dql-dlambda}\\
  \frac{\partial\,\mathrm{ELBO}_v}{\partial\,p_n^v(a\!\mid\theta_v)}
    &= q_v(\theta_v)\,(1-\lambda_v(\theta_v))\,\mathbb{E}_{\mathbf{p}_{-n}^v(\cdot|\theta_v)}\!\big[\log\!\tfrac{p_v(a,x_{-n},\theta_v)}{\mathcal{T}(a,x_{-n}\!\mid\theta_v)}\big],
    \label{eq:dql-dp}
\end{align}
where $\boldsymbol{\eqm}^{\theta_v}(x) := \prod_n \eqm^{(c_n,\theta_v)}(x_n)$,
$\mathbf{p}^v(x\!\mid\theta_v) := \prod_n p_n^v(x_n\!\mid\theta_v)$, and
$\mathbf{p}_{-n}^v$ omits site $n$. The $q_v(\theta_v)$ stationarity
\eqref{eq:dql-dq} gives the closed-form update
\begin{equation}
  q_v(\theta_v) \propto \exp\!\big(\mathbb{E}_{\mathcal{T}(\cdot|\theta_v)}[\log p_v(\cdot,\theta_v)] + \mathrm{H}[\mathcal{T}(\cdot\!\mid\theta_v)]\big).
  \label{eq:q-update}
\end{equation}
The $\lambda_v$ stationarity \eqref{eq:dql-dlambda} is the
``equal expected log-evidence ratio under both mixture components''
condition: at the stationary $\lambda$, the regenerated and
tracking components contribute the same average $\log p_v/\mathcal{T}$.
No closed form in general; solve by bisection on
$\lambda_v(\theta_v)\in[0,1]$ (the expression is monotone in
$\lambda$ once $\mathcal{T}$ is updated coherently). The $p_n^v$
stationarity \eqref{eq:dql-dp} is an exponential-family fixed
point: $\log p_n^v(a\!\mid\theta_v)$ equals the
$\mathbf{p}_{-n}^v$-averaged log of the $(p_v/\mathcal{T})$ ratio
at site $n$, up to a constant.

The mixture entropy $\mathrm{H}[\mathcal{T}]$ in
\eqref{eq:elbo-node} has no closed form in general but admits the
Jensen lower bound
$\mathrm{H}[\mathcal{T}(\cdot\!\mid\theta_v)] \geq (1-\lambda_v)\mathrm{H}[\mathbf{p}^v(\cdot\!\mid\theta_v)] + \lambda_v\mathrm{H}[\boldsymbol{\eqm}^{\theta_v}]$,
tight at $\lambda_v\in\{0,1\}$. Substituting the bound makes
\eqref{eq:elbo-node} and the gradients fully closed form at the
cost of a slightly looser ELBO. In practice M-projection
\eqref{eq:varform} provides a fast warm start; I-projection
fixed-point updates tighten the bound at additional per-sweep cost
and can be invoked when downstream behaviour indicates pathological
over-clustering bias. The two projections agree exactly at the
cap-2 cherry parent: the exact $F^v$ is already in the variational
family, so all four gradient expressions vanish at the exact
solution under both projections.

\paragraph{Coordinate ascent.} Standard message-passing: post-order
sweep (combine $+$ project at internal nodes), pre-order sweep
(downward updates), iterate. Per sweep $O(N\cdot m\cdot L\cdot k^2)$.
The ELBO is non-decreasing; typical convergence in $10$--$50$
sweeps.

\paragraph{Certified bound from the moment-matching marginals, without
iteration.} Because $\mathcal{L}(q)$ in \eqref{eq:elbo} is a valid
lower bound on $\log P(\text{data})$ for \emph{any} $q$ in the family
\eqref{eq:varform}, the M-projected marginals
$\{q_v(\theta),\,p_n^v(\cdot\!\mid\!\theta),\,\lambda_v(\theta)\}$ from
a single upward moment-matching pass already define such a $q$: plugging
them into \eqref{eq:elbo} yields a certified bound with \emph{no}
coordinate-ascent iteration. The essential caveat is that this is the
ELBO functional \emph{evaluated at} the marginals --- one extra tree
pass accumulating the leaf-emission terms, the per-edge cross-entropies
$\mathbb{E}_{q_u,q_v}[\log K]$, the per-node entropies, and $\log
Z_{\text{prior}}$ --- and is \emph{not} the log-normaliser $\log
Z_{\mathrm{MM}}$ that the moment-matching forward itself accumulates
while combining and propagating messages. The two are distinct
quantities:
\begin{equation}
  \mathcal{L}(q_{\mathrm{MM}}) \;\le\; \log P(\text{data})
  \quad\text{always, whereas}\quad
  \log Z_{\mathrm{MM}} \;\lessgtr\; \log P(\text{data}),
  \label{eq:mm-not-a-bound}
\end{equation}
since $\log Z_{\mathrm{MM}}$ is an expectation-propagation--style
marginal-likelihood \emph{estimate} that can fall on either side of the
truth. Only the functional $\mathcal{L}(q_{\mathrm{MM}})$ is a bound.
The evaluation pass costs the same $O(N\cdot m\cdot L\cdot k^2)$ as one
message-passing sweep, and bound validity is preserved under the Jensen
entropy lower bound of \eqref{eq:elbo-node} (lower-bounding $\mathrm{H}$
only lowers $\mathcal{L}$). At the cap-$2$ cherry geometry
$\mathcal{L}(q_{\mathrm{MM}}) = \log P(\text{data})$ exactly (no
projection fires and $q_{\mathrm{MM}}=q^\star$); at depth $>1$ the
M-projection's tendency to overestimate $\lambda_v$ merely loosens the
bound without invalidating it, so the certified moment-matching bound is
the recommended default and the I-projection sweeps
\eqref{eq:q-update}--\eqref{eq:dql-dp} are reserved for tightening only
when downstream over-clustering bias indicates the $\lambda$ slack
matters.

\paragraph{Tree-structured refinement: parent--child correlation at
mean-field cost.} The family \eqref{eq:varform} is mean-field
\emph{over nodes}: the edge term $\mathbb{E}_{q_u,q_v}[\log K]$ in
\eqref{eq:elbo} averages over \emph{independent} node marginals,
discarding the parent--child correlation. One would expect a
tree-structured $q$ that instead carries branch conditionals
$q(x_v,\theta_v\!\mid\!x_u,\theta_u)$ to reinstate the
super-exponential rank growth $K_v = K_{u_1}K_{u_2}+K_{u_1}+K_{u_2}$
that the mean-field restriction was introduced to tame. It does not,
because the field couples the $m$ sites \emph{only} through a single
per-branch event. Decompose each branch by its field-jump indicator
$\Delta\in\{0,1\}$ (\eqref{eq:arch-P0}--\eqref{eq:arch-P2plus}, with
$\Delta=1$ collecting $M\ge1$): conditioned on
$(\theta_u,\theta_v,\Delta)$ the sites evolve independently, so the
cheapest correlation-preserving branch conditional is the two-component
mixture
\begin{equation}
  q(x_v,\theta_v\!\mid\!x_u,\theta_u)
     = \!\!\sum_{\Delta\in\{0,1\}}\!\! q(\Delta,\theta_v\!\mid\!\theta_u)\,
        \prod_{n=1}^{m} q_\Delta(x_{v,n}\!\mid\!x_{u,n},\theta_u,\theta_v),
  \label{eq:tree-branch-cond}
\end{equation}
with a per-site tracking kernel $q_0(\cdot\!\mid\!x_{u,n})$ for the
no-jump component ($\Delta=0$, forcing $\theta_v=\theta_u$) and the
field-renewed stationary
$q_1(x_{v,n})=\eqm^{(c_n,\theta_v)}(x_{v,n})$, independent of
$x_{u,n}$, for the regenerated component ($\Delta=1$). This is
exactly the rank-$1$-plus-scalar form \eqref{eq:rank1plusscalar} read
as a variational \emph{branch conditional}: the $\Delta=0$ term is the
per-site rank-$1$ carry, the $\Delta=1$ term the residue-constant
shared renewal. Its cost is $O(m\,k^2)$ per branch --- the same order
as the mean-field edge term --- because $\Delta$ refactorises the
branch over sites; the parent--child correlation is carried entirely
by the shared $\Delta$ (all sites track together, or all renew
together, the ``simultaneous across-sites renewal'' of the
dynamic-field jump). Internal-node combines project the four-component
product back onto the two-$\Delta$-component form
\eqref{eq:tree-branch-cond} --- an M- or I-projection of the
\emph{conditional} rather than of the marginal --- so the rank stays
bounded at $2$ instead of growing with depth. Because
\eqref{eq:tree-branch-cond} recovers the node mean-field
\eqref{eq:varform} when the tracking kernel $q_0$ is frozen to the
product marginal, its optimised ELBO is never looser and generically
strictly tighter, at the same asymptotic cost, and remains exact on
cap-$2$ cherries (where no projection fires). The rank-$1$-plus-scalar
propagation is thus not merely a cap-$2$ computational shortcut but the
exact functional form of the cheapest correlation-preserving
variational branch conditional for the dynamic field.

\paragraph{Field-exact structured mean-field: the tight, tree-correlated
bound.} The node mean-field \eqref{eq:elbo} and the tree-structured branch
conditional above still \emph{approximate the field}; but the field carries
only $L$ states, so it is cheap to keep \emph{exact} while mean-fielding only
the $\mathcal{A}^m$ residues. Take
\begin{equation}
  q(\boldsymbol{\theta},\mathbf{x}_{\mathrm{int}})
    = q_\Theta(\boldsymbol{\theta})\,
      \prod_{v\ \mathrm{internal}}\prod_{n=1}^{m}
        q_{v,n}(x_{v,n}\mid\theta_v),
  \label{eq:field-exact-family}
\end{equation}
where $q_\Theta$ is an \emph{unrestricted tree-structured}
(Markov-on-the-tree) distribution over the whole field configuration ---
carried by its node marginals $q_v(\theta_v)$ and edge marginals
$q_{uv}(\theta_u,\theta_v)$ --- and $q_{v,n}(\cdot\mid\theta_v)$ is the
mean-field residue marginal at internal node $v$, site $n$, conditioned on the
\emph{local} field state (leaves fix $x$ to the observation). Crucially
$q_\Theta$ is \emph{not} factorised across nodes: the parent--child field
correlation that the node mean-field discards is retained in full.

Writing $\mu_{v,n}(\cdot\mid\theta):=q_{v,n}(\cdot\mid\theta)$ at internal
nodes and $\mu_{v,n}:=\delta_{x^{\mathrm{obs}}_{v,n}}$ at leaves, the ELBO is
\begin{align}
  \mathcal{L}
    &= \sum_{\theta} q_r(\theta)\Big[\log\rho(\theta)
         + \textstyle\sum_n \sum_a \mu_{r,n}(a\mid\theta)\log\eqm^{(c_n,\theta)}(a)\Big]
       \nonumber\\
    &\ + \sum_{u\to v}\sum_{\theta_u,\theta_v} q_{uv}(\theta_u,\theta_v)
         \Big[\log P_\theta(\theta_u\!\to\!\theta_v;\tau)
         + \textstyle\sum_{\Delta} r(\Delta\!\mid\!\theta_u,\theta_v)
             \sum_n \mathbb{E}_{\mu_{u,n},\mu_{v,n}}
                    [\log P^{(\Delta)}_{\mathrm{sub},n}]\Big]
       \nonumber\\
    &\ + \mathrm{H}[q_\Theta]
       + \sum_{v\ \mathrm{internal}}\sum_n\sum_{\theta} q_v(\theta)\,
             \mathrm{H}[q_{v,n}(\cdot\mid\theta)],
  \label{eq:field-exact-elbo}
\end{align}
with the field-only responsibility
$r(\Delta\!\mid\!\theta_u,\theta_v)
= P_{\mathrm{field}}(\theta_v,\Delta\mid\theta_u)/P_\theta(\theta_u\!\to\!\theta_v)$
(so the $\Delta$-sum inside $\log K$ is Jensen-bounded but its field factor
collapses to the exact $\log P_\theta$, cf.\ the certified-bound remark), and
the tree-Markov entropy in closed form,
\begin{equation}
  \mathrm{H}[q_\Theta] = \sum_{u\to v}\mathrm{H}[q_{uv}]
      - \sum_v (\deg v - 1)\,\mathrm{H}[q_v],
  \label{eq:field-tree-entropy}
\end{equation}
which is \emph{exact} for a tree-structured $q_\Theta$ --- on a tree the
junction-tree entropy is the true entropy, no Bethe approximation.

\emph{Coordinate ascent.} Two blocks. (i)~Fixing $q_\Theta$, each residue
marginal has the closed-form exponential-family update
$\log q_{v,n}(a\mid\theta)\propto$ [incident-edge and emission expected
log-potentials at $(v,n,\theta,a)$], a per-$(v,n,\theta)$ softmax over the
messages from $v$'s parent and children --- the residue mean-field of the
earlier paragraphs, now conditioned on $\theta$. (ii)~Fixing the residues,
\eqref{eq:field-exact-elbo} is, in $q_\Theta$, the free energy of an $L$-state
tree Markov random field with node potentials
$\varphi_v(\theta)=\exp(\text{root-prior}+\text{residue-emission expectations})$
and edge potentials
$\psi_{uv}(\theta_u,\theta_v)=P_\theta(\theta_u\!\to\!\theta_v)
\exp(\sum_\Delta r\sum_n\mathbb{E}[\log P^{(\Delta)}_{\mathrm{sub},n}])$;
because the graph is a tree, a single exact Felsenstein/belief-propagation pass
returns the \emph{exact} maximising node and edge marginals $q_v,q_{uv}$ ---
the field is never factorised. Cost per sweep
$O(N(L^2 + m\,L\,\mathcal{A}^2))$: the field block is $L^2$ per edge, the
residues mean-field.

\emph{Exactness (corrected).} At a single-branch tree and at a cap-$2$ cherry
the residue mean-field is exact (the rank-$1$+scalar form is closed under a
lone edge and under a single two-leaf combine) and the field BP is exact (a
two- or three-node MRF), so the maximised \eqref{eq:field-exact-elbo} equals
$\log P(\mathrm{data})$: its maximiser reproduces the moment-matching forward
wherever that forward is itself exact. This is the precise sense in which the
cap-$2$ ELBO is the exact log-likelihood --- it holds for the
field-\emph{correlated} family \eqref{eq:field-exact-family}. The
node-factorised family \eqref{eq:varform} does \emph{not} share it: factorising
the field across nodes leaves a strictly positive gap even at a cherry (it
drops the endpoint field correlation), so its ELBO is a valid but loose floor.
The field-exact family removes exactly that gap while keeping the
$\mathcal{A}^m$ residues mean-fielded.

\emph{Where this mean-field is exact.} Like the residue-correlated family
analysed below, \eqref{eq:field-exact-elbo} keeps the field factor tree-Markov,
so it is exact only when the residue-integrated field posterior is tree-Markov:
at a single branch (two field states, no separator to violate) it is exact for
any $m$, and on every tree with three or more field states --- cherries included
--- it is a strict bound, because the shared or chained residues couple field
states across their tree separators (the paths made explicit under
\eqref{eq:res-corr-branch}) and a tree-Markov field cannot hold them.
Node-factorising the residues on top of this drops the along-branch residue
correlation as well, so \eqref{eq:field-exact-elbo} is strictly looser than the
residue-correlated bound \eqref{eq:res-corr-branch} at the same
$O(N(L^2+m\mathcal{A}^2))$ cost; it is the floor. For $m=1$ the compound
$(\theta,x)$ chain has only $L\mathcal{A}$ states, so one keeps that joint and is
exact at any depth rather than mean-fielding at all --- the $m=1$ instance of the
joint form to which cherry- and singleton-exactness properly belong (below).

\paragraph{Residue-correlated tree-structured family: the headline
bound.} The residue mean-field's two remaining gaps --- singletons at
depth ${\ge}2$ and the $m\ge2$ interior --- both come from dropping the
\emph{along-branch} residue correlation. Restore it without leaving the
tractable class by making the residues, too, Markov on the tree
\emph{conditioned on the field and its jumps}. Promote the per-branch jump
indicator $\Delta\in\{0,1\}$ to a variational latent and take the branch
conditional
\begin{equation}
  q(\mathbf{x}_v\mid\mathbf{x}_u,\theta_u,\theta_v,\Delta)
    = \mathbb{I}[\Delta{=}0]\,\mathbb{I}[\theta_v{=}\theta_u]
        \prod_{n=1}^{m} q_{n}^{uv}(x_{v,n}\mid x_{u,n})
    + \mathbb{I}[\Delta{=}1]\prod_{n=1}^{m} r_{n}^{v}(x_{v,n}\mid\theta_v),
  \label{eq:res-corr-branch}
\end{equation}
so the full $q$ is Markov-on-the-tree over $(\boldsymbol\theta,\boldsymbol\Delta,
\mathbf{x})$: a tracking component ($\Delta{=}0$, no field jump) that carries a
free per-site parent$\to$child kernel $q_n^{uv}$ --- the along-branch
correlation the mean-field dropped --- and a renewal component ($\Delta{=}1$)
in which the sites are drawn from a free per-site node marginal $r_n^v$, not
the prior stationary, so it absorbs downward evidence. The mean-field
\eqref{eq:field-exact-family} is the special case $q_n^{uv}(x'\mid x)=q_{v,n}(x')$.
Equation~\eqref{eq:res-corr-branch} \emph{specifies} $q$ in the rooted
conditional parametrisation (a product of branch conditionals $q(x_v\mid x_u)$
off a root factor); the entropy below is written in the equivalent
\emph{moment} parametrisation, whose edge marginal is
$q_{uv}(x_u,x_v)=q_u(x_u)\,q(x_v\mid x_u)$ (node marginal $\times$ branch
conditional). The two describe the same tree-Markov family --- conditionals are
convenient for specifying and sampling $q$, marginals for the junction-tree
entropy --- and give the same $\mathrm{H}[q]$
($\mathrm{H}[q_{uv}]=\mathrm{H}[q_u]+\mathrm{H}[x_v\mid x_u]$).

\emph{Tractability is unchanged.} The decisive property survives: conditioned
on $(\theta_u,\theta_v,\Delta)$ the residues \emph{factorise over sites}, so the
edge pairwise belief is a product of $m$ blocks of size $\mathcal{A}\times
\mathcal{A}$, never the $\mathcal{A}^{2m}$ joint. With $\Delta$ in $q$ the energy
is a plain sum (no $\log\!\sum_\Delta$), and the junction-tree entropy
\begin{equation}
  \mathrm{H}[q] = \sum_{u\to v}\mathrm{H}[q_{uv}] - \sum_v(\deg v-1)\,\mathrm{H}[q_v],
  \qquad
  \mathrm{H}[q_{uv}] = \mathrm{H}[q_{uv}^{\Theta\Delta}]
     + \!\!\sum_{\theta_u,\theta_v,\Delta}\!\! q_{uv}(\theta_u,\theta_v,\Delta)
       \sum_n \mathrm{H}[q_n^{uv}(\cdot,\cdot\mid\cdot)],
  \label{eq:res-corr-entropy}
\end{equation}
is computed per site; cost stays $O\!\big(N(L^2 + m\,\mathcal{A}^2)\big)$ per
sweep --- the \emph{same class} as the residue mean-field. Coordinate ascent
alternates the per-site kernel updates (a closed-form
$\log q_n^{uv}\!\propto\!$ incident-edge $+$ emission expected log-potentials)
with an exact field-and-jumps belief propagation over the induced $(\theta,
\Delta)$ tree MRF; for $m\ge2$ the two-child combine, which produces three
distinct rank-$1$ residue vectors whose sum does not factorise over sites, is
projected back onto \eqref{eq:res-corr-branch} (the sole approximation).

\emph{Where the bound is strict, and where exactness lives instead.}
Equation~\eqref{eq:res-corr-branch} keeps $q_{\Theta\Delta}$ tree-Markov and
couples the residues to it only through the ELBO's \emph{expectations} --- the
field-jump law $q(\theta_v,\Delta\mid\theta_u)$ does not depend on the residue
\emph{values} $x_u$. So the maximised ELBO equals $\log P(\mathrm{data})$
exactly when the true field-jump posterior
$P(\boldsymbol\theta,\boldsymbol\Delta\mid\mathrm{data})
 \propto P_{\mathrm{field}}\prod_n L_n(\boldsymbol\theta,\boldsymbol\Delta)$
is tree-Markov. Each per-site evidence
$L_n=\sum_{\mathbf{x}_n}\prod_e T^{\Delta_e}_\theta$ is a whole-tree Felsenstein
contraction, so integrating the residues out couples field states through paths
that run \emph{through the residues}: on a chain
$\theta_1\!-\!x_1\,\text{---}\,\theta_2\!-\!x_2\,\text{---}\,\theta_3\!-\!x_3$
the route $\theta_1\!-\!x_1\!-\!x_2\!-\!x_3\!-\!\theta_3$ bypasses the separator
$\theta_2$, and at a cherry the shared root residue couples the two leaf field
states past $\theta_r$. A tree-Markov $q_{\Theta\Delta}$ cannot carry these.
Hence:
\begin{itemize}\itemsep0pt
  \item \emph{$\le 2$ field nodes (a single branch), any $m$: exact.} Two field
        states have no separator to violate, so $q_{\Theta\Delta}$ is
        unconstrained and the bound is tight.
  \item \emph{$\ge 3$ field nodes (cherries and every deeper tree): a strict
        bound}, for $m\ge2$ \emph{and} $m=1$ alike. The residue-induced field
        couplings above are dropped. The gap is the residue$\to$field feedback
        --- second-order and small when the field drives the residues more than
        the residues identify the field ($\sim0.06$--$0.5$ nat on cherries and
        depth-$2$ trees, vanishing as $\rho_{\mathrm{chain}}\!\to\!0,\infty$).
        Variationally this is the forward's combine blow-up
        ($K_v=K_{u_1}K_{u_2}+K_{u_1}+K_{u_2}$).
\end{itemize}
The exactness that \emph{is} available for cherries at any $m$, for singletons
at any depth, and at $\rho_{\mathrm{chain}}=0$ belongs instead to the
\emph{joint} form that keeps $(\theta,x)$ in one message --- the rank-$1$+scalar
CLV \eqref{eq:rank1plusscalar}, in which the field jump is coupled to the actual
residue values. Its \emph{forward} is exact through $|C|=2$ and at cherries; but
its ELBO \emph{functional} needs $\mathbb E_q[\log F]$ over that joint message,
and $\log$ of an $\mathcal{A}^m$-site product does not factorise, so the
certified joint bound costs $O(\mathcal{A}^m)$ --- cheap only at $m=1$
($L\mathcal{A}$) or $m\le2$. The residue-correlated family
\eqref{eq:res-corr-branch} recovers most of that tightness at
$O(N(L^2+m\mathcal{A}^2))$: it strictly \emph{dominates} the residue mean-field
\eqref{eq:field-exact-family} (restoring the along-branch residue correlation)
at the same asymptotic cost, and is the recommended scalable bound;
\eqref{eq:field-exact-family} is retained only as the cheaper floor. Its
M-projected single forward pass is a valid certified bound provided it is
evaluated through \eqref{eq:res-corr-entropy}, not its own log-normaliser.

\paragraph{Tight regimes.} The maximised ELBO of the node-factorised
family \eqref{eq:varform} equals $\log P(\text{data})$ only as
$\rho_{\mathrm{chain}}\!\to\!\infty$: every edge regenerates
($\lambda_v\equiv 1$, $q_v(\theta_v)\to\eqm_\theta$), the field states become
independent across nodes, and factorising $q_\Theta$ across nodes then costs
nothing. At the cap-2 cherry and as $\rho_{\mathrm{chain}}\!\to\!0$ the exact
object is instead the field-correlated family \eqref{eq:res-corr-branch}
(equivalently the moment-matching forward's $\log Z_{\mathrm{MM}}$, exact
wherever no combine-projection fires): the cherry couples the root and leaf
field states, and at $\rho_{\mathrm{chain}}\!\to\!0$ the branch factor forces
$\theta_v=\theta_u$ --- neither of which the node-factorised $q_\Theta$ can
represent, so \eqref{eq:varform} keeps a strictly positive gap there. The
TKF-DP training loop scores cap-2 cherries by that forward; the
residue-correlated bound \eqref{eq:res-corr-branch} carries the exactness into
depth-$>1$ inference (ancestor reconstruction, tree-level held-out validation)
while remaining a certified lower bound at intermediate $\rho_{\mathrm{chain}}$
on deeper trees.

\paragraph{Per-site class assignments under the variational family.}
The form \eqref{eq:varform} treats the per-site class $c_n$ as a
fixed attribute of site $n$, not a variable to be integrated. At
the cap-2 cluster geometry the class assignments are marginalised
exactly by the $K_c^2$-term sum
$\sum_{c_1, c_2} \pi_c(c_1)\pi_c(c_2)\,P^{(c_1, c_2, \theta_P)}$,
which is bounded (at $K_c = 8$, $K_c^2 = 64$). At general cluster
size $m$ the analogous marginalisation costs $K_c^m$, destroying
the linear-in-$m$ scaling that the dynamic field was designed to
deliver. Two paths preserve linearity:
\begin{enumerate}\itemsep0pt
  \item \emph{Sampling.} Per-site class is a sampled hard
        assignment under a separate Gibbs sweep on $c_s$ given
        the cluster's coupling-side observations. This is the
        path taken by the TKF-DP training loop, where cap-2
        cherries are the unit and per-cherry signal is too weak
        to drive a mean-field variational marginal to a hard
        fixed point reliably.
  \item \emph{Mean-field per-site marginals.}
        $q_n^v(c_n) \in \Delta^{K_c-1}$ is added to the per-node
        variational state and marginalised inside the cluster
        emission,
        $\sum_{c_1, \dots, c_m} \prod_n q_n^v(c_n)\,
         P^{(c_1, \dots, c_m, \theta_v)}$.
        Cost stays $O(L\cdot m\cdot k^2)$ per cluster.
\end{enumerate}

Path~2 captures anti-correlation between coupled sites \emph{only
at strong-signal coordinate-ascent fixed points} where each
$q_n^v$ collapses to a Dirac on $c_n^*$, recovering the
hard-assignment emission $P^{(c_1^*, \dots, c_m^*, \theta_v)}$ that
exposes the per-class compositional bias responsible for
anti-correlation. The hard-assignment limit is contained in the
mean-field family as a degenerate special case, and at
$\KL$-strong evidence the coordinate-ascent fixed point is
provably the Dirac. At weak signal (few observations per cluster,
$K_c$ large, or two symmetric modes such as $(c_1^*, c_2^*)
\leftrightarrow (c_2^*, c_1^*)$ contributing equal likelihood),
the mean-field marginals stay diffuse, the emission averages over
$K_c^m$ class assignments---including the $K_c^{m-1}$ same-class
assignments that are typically neutral on anti-correlation---and
the variational emission posterior loses the anti-correlation
structure that motivates the per-site class DP in the first place.

Mitigations of escalating richness, applicable when downstream
metrics show the diffuse-marginal failure mode:
\emph{structured-mean-field per cluster pair}, replacing
$\prod_n q_n^v$ with
$\prod_{(i, j) \in \text{pairs}(C)} q_{ij}^v(c_i, c_j)$ (adds
$K_c^2$ per pair, $K_c^2 \cdot \binom{m}{2}$ per cluster, recovers
$(c_1^*, c_2^*) \leftrightarrow (c_2^*, c_1^*)$ symmetry exactly);
\emph{warm-start} the coordinate ascent from a sampled hard
$c_s$ so the descent begins inside the basin of the relevant
Dirac; \emph{entropy tempering} on $q_n^v(c_n)$ that pushes
toward hard assignments early in the optimisation and relaxes
once the rest of the variational state has settled. The TKF-DP
training loop sidesteps the issue entirely via path~1 (sampling);
the mean-field class extension is needed only when extending the
ELBO to general phylogenies, where aggregated per-tree signal at
each site is stronger than the per-cherry signal that defeats it
at training scale.

\paragraph{Parity with the Potts variant.}
\Cref{prop:hier-fels} together with this ELBO extension puts the
dynamic-field variant on equal structural footing with the Potts
variant's Cohn/Girsanov bound (\autoref{sec:class-elbo}): both are
exact at the cap-2 cluster training geometry and both admit a
principled variational extension to deeper trees. Combined with the
no-Sinkhorn-needed property of the field-marginal joint stationary
under the dynamic-field variant---the joint
$J^{(c_1,c_2)}(a,b)=\sum_\theta \eqm_\theta\,\eqm^{(c_1,\theta)}(a)\,\eqm^{(c_2,\theta)}(b)$
is marginal-consistent by construction, so \autoref{thm:revcond} is
automatically satisfied---the dynamic-field variant is a drop-in
substitute for the Potts variant throughout the inference stack.

\subsection{Survival of the approximate evaluators under A1 and under the dynamic field}
\label{sec:approx-suppl}

The TKF-DP inference loop evaluates the substitution likelihood by
three techniques; each carries over under A1 and under the
dynamic-field variant of \autoref{sec:dynamic-suppl}.

\paragraph{Cohn/Girsanov class-level variational ELBO.} The bound
on the coupled-CTMC path measure on matched (both-present) branches
via a constant-rate variational family with closed-form pairwise
bridge expectations (\autoref{sec:class-elbo}) survives unchanged
in form: a matched branch is still a coupled reversible CTMC, now
with the Sinkhorn-corrected $\Q_{\mathrm{pair}}$; the
marginal-consistent stationary is plugged in directly, and the
$H=0$ exactness is preserved. A1 changes only the indel-seam
scoring (the lone-branch stationary and the conditional-Boltzmann
insertion), which is outside the coupled-interior bound. Under the
dynamic-field variant the bound becomes moot at the cap-2 cluster
geometry: sites are conditionally independent given the field, the
re-equilibrating dynamic field model is not a CTBN, and both are
handled by \Cref{prop:hier-fels} at depth $1$. For inference on
general phylogenies under the dynamic-field variant
(ancestor-residue reconstruction at internal nodes, tree-level
held-out validation), the structured mean-field ELBO of
\autoref{sec:variational-suppl} restores tractable
linear-in-cluster-size cost while remaining exact at cap-2 and in
the rate-limit regimes; it plays the same structural role for the
dynamic field as Cohn/Girsanov does for the Potts coupling.

\paragraph{Complete-data conjugacy.} The secret-destination
Dirichlet conjugacy on $\eqm^{(c)}$ and the Gamma--Poisson on rates
survive, and A1 makes them cleaner. Under A1 the pair-marginal of
class $c$ is exactly $\eqm^{(c)}$ regardless of partner, so the
per-class sufficient-statistic counts aggregate without
partner-dependent reweighting (under the free-$h$ A3 model they
would not, since $c$'s effective composition varies per partner).
The side potentials $h$, now deterministically Sinkhorn-determined
by $(\eqm,H)$, are demoted from sampled parameters to a
deterministic readout, so no prior or update is needed for them.
Gamma--Poisson for the indel and rate constants is untouched.
Under the dynamic-field variant the per-(class, field) stationary
$\eqm^{(c,\theta)}$ admits the same Dirichlet conjugacy on sufficient
statistics gathered per $(c,\theta)$ bin; the field selector itself
is updated by a separate truncated-stick-breaking Beta posterior on
the per-cluster field assignment counts.

\paragraph{Laplace for the Potts atoms $H$.} $H$ remains the free,
real-valued, non-conjugate coupling parameter under A1, so the
Gaussian (Hessian-based) Laplace approximation still applies. The
only change is that because $h=\mathrm{Sinkhorn}(\eqm,H)$ is now an
inner deterministic step, the Laplace objective is
$H\mapsto h(H)\mapsto \pj\mapsto$ likelihood, and the Hessian
acquires the Sinkhorn-map Jacobian $\partial h/\partial H$
(well-defined and differentiable by the implicit function theorem).
The multi-seed Laplace mixture for multimodal $H$ posteriors
carries over verbatim. Under the dynamic-field variant the
$H$-parameter is replaced by the per-(class, field) stationary
$\eqm^{(c,\theta)}$, which is conjugate to its Dirichlet base
measure; no Laplace approximation is required on the
substitution-side parameters of the dynamic field.

\paragraph{Summary.} A1 preserves all three evaluators: the ELBO is
unchanged on matched branches; conjugacy is cleaner; Laplace
acquires an extra Jacobian. The dynamic-field variant replaces the
coupled-CTMC ELBO with the exact hierarchical recursion of
\Cref{prop:hier-fels} at cap-2 and its variational extension of
\autoref{sec:variational-suppl} on deeper trees, retains conjugacy
on the per-(class, field) stationary, and obviates Laplace on the
substitution side entirely.

\subsection{Augmented indel histories via a time-indexed pair SCFG}
\label{sec:gravestone-scfg}

The class-level ELBO of \S\ref{sec:class-elbo} was derived under
fixed class composition over a constant-composition segment of a
branch. To compose with the tree we need the timed indel history on
each branch, including \emph{transient} insertions: under coupling,
an insertion that is born and then deleted on the same branch is not
visible in the alignment but does change the rate of every other
site in its DP class during its lifetime. The naive scheme of
sampling alignment-visible indel timings from their TKF92 prior is
inadequate for two distinct reasons: first, transient lineages are
missed; second, even for visible events, the conditional density
given the alignment depends on substitution evidence and is not what
the alignment marginal alone provides.

We resolve both issues by augmenting the latent state with
\emph{gravestones}. The augmentation does not change the model. It
changes the latent representation: every fragment that ever existed
during the branch --- alive at start, alive at end, deleted on the
branch, or transiently inserted then deleted --- is represented as
a labeled position in the augmented branch history. With gravestones
the alive count $N_a(s)$ is bounded by the visible count $N_v(s)$,
so the indel CTMC bridge has a finite state space and admits a
closed-form recursive sampler. The gravestone counts at internal
nodes of the tree become latent variables resampled jointly with the
rest of the augmented history.

\subsubsection*{The branching process and its grammar}

We first present the TKF91 case (single residue per fragment) where
the branching structure is cleanest, then read off the TKF92
expansion. Let $L(t)$ denote a single link observed $t$ time units
after its birth. Under TKF91 dynamics, $L(t)$ has lifetime
$\tau \sim \mathrm{Exp}(\mu)$, and during its lifetime spawns child
links at rate $\lambda$. Each child is itself an $L$ of the
appropriate residual age. The branching process is a
continuous-time Galton--Watson process expressible as a stochastic
context-free grammar with productions indexed by continuous time:
\begin{align*}
L(t) &\;\to\; A \cdot D(t,\, t)
& &\text{at weight } e^{-\mu t}, \\
L(t) &\;\to\; G_\tau \cdot D(t,\, \tau)
& &\text{at density } \mu e^{-\mu \tau}\, d\tau, \;\; \tau \in (0, t), \\
D(t,\, T) &\;\to\; L(t - u) \cdot D(t,\, T)
& &\text{at density } \lambda \, du, \;\; u \in (0, T), \\
D(t,\, T) &\;\to\; \varepsilon
& &\text{at weight } e^{-\lambda T}.
\end{align*}
The terminals $A$ and $G_\tau$ stand for an alive residue and a
gravestone residue that died at relative time $\tau$, respectively,
each carrying a DP key, a site class, and an amino-acid trajectory
over its lifetime as governed by \S\ref{sec:class-elbo}. The
non-terminal $D(t, T)$ generates a list of child links born during
a window of size $T$ inside an observation interval of size $t$.
The reading is: a link either survives to observation (probability
$e^{-\mu t}$) and can have spawned children at any time in
$(0, t)$, or it died at some $\tau \in (0, t)$ (density
$\mu e^{-\mu \tau}$) and can only have spawned children during
$(0, \tau)$. The empty-list weight $e^{-\lambda T}$ is the
probability of zero events from a Poisson process of rate $\lambda$
over a window of $T$.

\paragraph{TKF92 expansion.} The generalization to TKF92 replaces
each single-residue terminal with a fragment of geometrically-distributed
length, leaving the fragment-level branching structure entirely
unchanged:
\begin{align*}
A &\;\to\; \mathcal{R}^A \cdot A^*,
& A^* &\;\to\; \mathcal{R}^A \cdot A^* \;\;\text{at weight } r,
& A^* &\;\to\; \varepsilon \;\;\text{at weight } 1 - r, \\
G_\tau &\;\to\; \mathcal{R}^G_\tau \cdot G_\tau^*,
& G_\tau^* &\;\to\; \mathcal{R}^G_\tau \cdot G_\tau^* \;\;\text{at weight } r,
& G_\tau^* &\;\to\; \varepsilon \;\;\text{at weight } 1 - r,
\end{align*}
where $\mathcal{R}^A$ is a single alive residue and
$\mathcal{R}^G_\tau$ a single gravestone residue with shared death
time $\tau$. Fragment death is collective: all residues of a
fragment share its lifetime $[0, \tau]$ and die together at $\tau$.
Key-DP labels $z_s$ and class-DP labels $c_s$ are drawn
independently per residue at fragment birth from their respective
Chinese restaurant predictives, even within a single fragment, so a
fragment can carry residues of several key classes and several site
classes. The expansion is invisible to the indel skeleton --- the
bivariate Riccati of the next subsection lives at the fragment level
--- but enters the substitution likelihood and the DP-class
composition through the per-residue payload.

Pairing this with an ancestor sequence yields the pair version: each
ancestor fragment present at branch start is an independent root of
the recursion above, with $t = t_b$ the branch length; insertions
appearing in the children list have no ancestor counterpart and
carry only a descendant payload. The pair grammar generates an
aligned pair of fragment sequences (with alive/gravestone labels on
the descendant side) jointly with the parentage tree.

\subsubsection*{Inside generating function and the bivariate Riccati}

Let $f_t(z, w) = \mathbb{E}[z^{N_a(t)} w^{N_g(t)}]$ where $N_a, N_g$
count alive and gravestone descendants of one root fragment of age
$t$. The inside generating function of the grammar satisfies the
Volterra fixed point
\[
f_t(z, w)
\;=\; z\, e^{-\mu t}\, \exp\!\Bigl(\lambda \!\int_0^t (f_s(z, w) - 1)\, ds\Bigr)
\;+\; w\!\int_0^t \mu e^{-\mu \tau}\, \exp\!\Bigl(\lambda \!\int_{t-\tau}^t (f_s(z, w) - 1)\, ds\Bigr)\, d\tau,
\]
which is equivalent to the forward equation
\[
\partial_t f_t(z, w) \;=\;
\bigl[\lambda z^2 - (\lambda + \mu) z + \mu w\bigr] \, \partial_z f_t(z, w),
\qquad f_0(z, w) = z.
\]
The right-hand side is quadratic in $z$, so the characteristic
equation $\dot z = -[\lambda z^2 - (\lambda + \mu) z + \mu w]$ is
Riccati. Define
\[
\Delta(w) \;=\; \sqrt{(\lambda - \mu)^2 + 4 \lambda \mu (1 - w)},
\qquad
z_\pm(w) \;=\; \frac{(\lambda + \mu) \pm \Delta(w)}{2 \lambda}.
\]
The Riccati flow
$(z(t) - z_+) / (z(t) - z_-) = (z(0) - z_+)/(z(0) - z_-) \cdot e^{-\Delta(w) t}$
inverts to give $f_t(z^*, w)$ in closed form, hence closed-form
pointwise transition probabilities
$P\!\bigl((N_a, N_g)(t)\,\big|\,(N_a, N_g)(0)\bigr)$ by coefficient
extraction. Setting $w = 1$ recovers the standard linear
birth-death PGF for $N_a$ alone.

\subsubsection*{Recursive traceback sampler}

Given an augmented history with endpoint counts
$(a_0, g_0) = (N_a(0), N_g(0))$ and
$(a^*, g^*) = (N_a(t_b), N_g(t_b))$, the within-branch indel
history is sampled by recursive midpoint traceback on the bivariate
process $(N_a, N_g)$. At each level of the recursion, sample a
midpoint state from
\[
P\!\bigl((a_m, g_m) \,\big|\, (a_0, g_0), (a^*, g^*)\bigr)
\;\propto\;
P\!\bigl((a_0, g_0) \to (a_m, g_m);\, t_b/2\bigr) \cdot
P\!\bigl((a_m, g_m) \to (a^*, g^*);\, t_b/2\bigr),
\]
both factors closed form via the Riccati. The midpoint state space
is finite since $a_m, g_m \leq g^* + a^*$. Recurse on each half.
The base case is reached when the residual subproblem has at most a
small number of events (births and deaths), at which point the
closed-form list of valid event interleavings is enumerable
directly and event times are sampled from the corresponding
hypoexponential simplex density. Total cost is
$O((b + d) \log(b + d))$ per branch where $b, d$ are the branch
totals of births and deaths.

After event times are placed, parentage is sampled forward in time:
at each birth, a uniform alive fragment is chosen as the parent; at
each death, a uniform alive fragment is killed. Parentage uniformity
is the standard exchangeability of the linear Galton--Watson tree.

\subsubsection*{Inserted fragment and residue distributions}

When a child fragment is added at time $u$ to a DP class $C$ that
already contains alive fragments, the inserted residues are drawn
from the conditional Boltzmann
\[
P\!\bigl(x_s(u^+) \,\big|\, x_{C \setminus s}(u^-)\bigr)
\;\propto\;
\pi^{(c_s)}\!\bigl(x_s(u^+)\bigr)\, \exp\!\Big(-\!\!\sum_{s' \in C \setminus s}\!\! H_{c_s c_{s'}}\!\bigl(x_s(u^+),\, x_{s'}(u^-)\bigr)\Big),
\]
which is exactly normalizable on the alphabet $\mathcal{A}$ even
when the joint stationary on $C$ is intractable. This is the proper
coupled generalization of the $\mathrm{Cat}(\pi)$ inserted-residue
rule of standard TKF92.

\subsubsection*{Composition with substitution and with the rest of
the tree}

The traceback gives a piecewise-constant DP-class composition on
each branch, with segment boundaries at every (alive or gravestone)
indel event. The variational ELBO of \S\ref{sec:class-elbo} applies
to each segment as written, with the segment's class composition
and duration. Bridge expectations accumulate over
segments within each branch and across branches.

Internal-node augmented states $(N_a, N_g)$-valued are sampled
jointly across the tree. The standard upwards-downwards
(Felsenstein-style) recursion applies because the closed-form
per-branch transition kernel from the Riccati composes across
branches: the marginal at each internal node, conditional on the
rest of the tree, is a closed-form pointwise distribution
computable in $O(1)$ given the per-branch kernel evaluations. The
augmentation thus extends to a full tree-level closed-form sampler.

\subsubsection*{Marginalizing the augmentation}

Gravestones are an augmentation of the latent state, not of the
data. The marginal likelihood of the observed alignment under the
model is recovered by summing over gravestone configurations at
every internal node and along every branch. The recursive traceback
samples the augmented state, and the variational ELBO of
\S\ref{sec:class-elbo} applied segment-wise on the augmented
branches gives the conditional substitution likelihood. No other
approximation enters at the indel level: the SCFG is the exact
representation of TKF92's fragment-level dynamics, and the
gravestone augmentation is what makes its inside-outside computable.

\begin{remark}[Coalescent priors; ARG augmentation as a parallel construction]
The per-branch SCFG depends on branch length only through $t_b$, so
any prior on the tree --- including a coalescent prior on branch
times --- composes with the gravestone-augmented sampler without
modification. This is a free generalization that supports joint
indel-aware coalescent inference on protein-coding regions, which
existing coalescent frameworks typically handle by treating indels
as missing data. Separately, the structural mechanism --- augmenting
an unbounded latent process with non-observable entities so that an
observed count provides a state-space bound --- carries over to the
ancestral recombination graph: non-ancestral lineages in the ARG
play the role of gravestones, bounding the per-genome-position bridge
state space and admitting (we conjecture) a parallel closed-form
recursive sampler.
\end{remark}

\subsection{Cluster class-composition prior: a mixture of Dirichlet-multinomials}
\label{sec:dm-mixture}

The generative model of \S\ref{sec:tkfdp-model} draws the site classes $c_s$
\emph{iid} from a single global $D_c$, so it says nothing about which classes
prefer to co-occur \emph{within a cluster}. Yet with the partition $z$ and the
labels $c$ jointly sampled, each cluster $C$ furnishes a direct multinomial
observation of its own class composition. We exploit this to give clusters a
learned, \emph{exchangeable}, cap-agnostic preference over class
\emph{combinations} --- the coevolutionary analogue of ``which types like to
sit together'' --- with no pairwise ($O(K_c^2)$, cap-2-bound) Potts term.

\paragraph{Generative model.} With $H$ truncated mixture components (e.g.\
$H=20$), draw stick-breaking weights $\pi\sim\mathrm{TSB}(\alpha_\pi)$ and, per
component, a Dirichlet hyperparameter $\alpha_h\in\mathbb{R}_{>0}^{K_c}$. For
each cluster $C$ of any size $m=|C|$,
\begin{align}
  h_C &\sim \mathrm{Cat}(\pi), &
  \theta_C \mid h_C &\sim \mathrm{Dir}(\alpha_{h_C}), &
  c_s \mid \theta_C &\stackrel{\text{iid}}{\sim} \mathrm{Cat}(\theta_C)\ \ (s\in C).
\end{align}
Only the class \emph{labels} are drawn this way; the residues still evolve
under the same per-cluster field likelihood $\mathcal{L}(C)$ of
\S\ref{sec:dynamic-suppl} (the prior multiplies, it does not replace, the
emission model). With class counts $n_C=(n_{C,1},\dots,n_{C,K_c})$,
$\sum_k n_{C,k}=m$, and $A_h=\sum_k\alpha_{h,k}$, integrating $\theta_C$ out
(Dirichlet--multinomial conjugacy) and marginalising the component gives the
cluster class-composition prior
\begin{equation}
  \Pr(n_C\mid\alpha,\pi) \;=\; \sum_{h=1}^{H}\pi_h\,
    \mathrm{DM}(n_C\mid\alpha_h),\qquad
  \mathrm{DM}(n_C\mid\alpha_h) = \frac{\Gamma(A_h)}{\Gamma(A_h+m)}
    \prod_{k=1}^{K_c}\frac{\Gamma(\alpha_{h,k}+n_{C,k})}{\Gamma(\alpha_{h,k})}.
  \label{eq:dm-mixture}
\end{equation}
Since \eqref{eq:dm-mixture} depends only on the count vector $n_C$ it is
exchangeable in the columns and defined for \emph{any} cluster size (the
cap-$>2$-safe property); the full cluster score is
$\mathcal{L}(C)\cdot\Pr(n_C\mid\alpha,\pi)$ and the joint prior is
$\prod_C\Pr(n_C\mid\alpha,\pi)$.

\paragraph{Co-occurrence, not arrangement.} A component $\alpha_h$ encodes
which classes \emph{co-occur} in a type-$h$ cluster; being exchangeable it does
not by itself prefer any \emph{arrangement} across columns. For the
enumerated-class field model (a class is an ordered pair $(k_0,k_1)$ of
archetypes at the two field states) a compensatory flip is a cluster carrying
classes $(A,B)$ and $(B,A)$: the mixture makes such class \emph{sets} likely,
while the field likelihood $\mathcal{L}(C)$ --- through the shared field
trajectory --- selects the compensatory arrangement over the co-flip
$(A,B),(A,B)$. The Dirichlet--multinomial's positive same-class correlation in
fact \emph{mildly disfavours} the distinct-class (compensatory) pairing, so the
arrangement is carried entirely by $\mathcal{L}$; the concentration $A_h$
should be large enough that this clumping does not fight the likelihood. We add
\emph{no} flip affinity to the prior (a per-cluster transposition kernel
$c_s\mapsto\mathrm{swap}(c_s)$ w.p.\ $p$ preserves exchangeability and is
available, but is omitted) so that flip structure must \emph{emerge}: a fitted
component whose mass concentrates on a charge-complementary class set is then
an unsupervised, directly interpretable ``salt-bridge type.''

\paragraph{Sampling.} All four updates are standard.
\begin{enumerate}
  \item \emph{Component} $h_C$: Gibbs, $\Pr(h_C=h\mid n_C)\propto
    \pi_h\,\mathrm{DM}(n_C\mid\alpha_h)$ (length-$H$ categorical; or marginalise
    $h$ via \eqref{eq:dm-mixture}).
  \item \emph{Site class} $c_s$: the existing class-Gibbs gains the
    Dirichlet--multinomial predictive log-prior. Given $h_C=h$, moving $s$ to
    class $k'$ from $k$ contributes
    \begin{equation}
      \Delta\log\Pr \;=\; \log\frac{\alpha_{h,k'}+n^{-s}_{C,k'}}
                                   {\alpha_{h,k}+n^{-s}_{C,k}},
      \label{eq:dm-cn}
    \end{equation}
    with $n^{-s}_C$ the counts excluding $s$ ($\Gamma(A_h+m)$ cancels) --- the
    ``rich-get-richer within the cluster's type'' term, additive to the
    field-likelihood evidence at $O(1)$ (no forward). With $h$ marginalised it
    is the mixture ratio
    $\sum_h\pi_h\mathrm{DM}(n^{k'}_C\mid\alpha_h)\big/
    \sum_h\pi_h\mathrm{DM}(n_C\mid\alpha_h)$.
  \item \emph{Partition} $z_s$: a column moving $C\!\to\!C'$ updates both
    $n_C,n_{C'}$; the move already rescores the affected clusters' likelihood,
    and \eqref{eq:dm-mixture} for the two clusters is added.
  \item \emph{Hyperparameters} $\alpha_h$: a Dirichlet--multinomial MLE (Minka
    fixed-point) or Metropolis step from the \emph{pooled} counts of clusters
    currently in component $h$ --- each cluster gives only $m$ draws, so the
    components are identified by pooling, not from any single cluster.
  \item \emph{Weights} $\pi$: truncated-stick-breaking posterior (Beta
    conjugacy / Escobar--West auxiliary Gibbs on $\alpha_\pi$), as for the other
    DP concentrations.
\end{enumerate}

\begin{remark}[Interpretability]
The fitted $\{\alpha_h\}$ are a taxonomy of coevolutionary cluster types read
straight off the model: rank each active component's classes by $\alpha_{h,k}$
and inspect their archetypes' charge composition. A component concentrating on
charge-complementary flip-classes is a salt-bridge type; one on diagonal
(static $k_0=k_1$) or hydrophobic classes is a specialiser/packing type --- with
no supervision or built-in flip bias used to obtain them.
\end{remark}

\subsection{Posterior sampling and parameter learning}
\label{sec:tkfdp-svi}

We turn to the inverse problem: given a corpus of MSAs
$\{X_n\}_{n=1}^N$ with associated trees $\{T_n\}$ and fixed
alignment paths, learn the corpus-level substitution hyperparameters
\[
\theta = \bigl(S,\, \alpha_z,\, \alpha_c,\, \alpha_H,\, a_\eta, b_\eta,\,
\kappa_\pi, \bar\pi,\, \{\mu_{kl}, \tau_{kl}\}_{k \leq l}\bigr),
\]
namely the shared exchangeability $S$, the three DP concentrations
$(\alpha_z, \alpha_c, \alpha_H)$, the per-site rate prior
$(a_\eta, b_\eta)$, and the base-measure parameters for $G_0$ on
$\pi^{(c)}$ and for $G_0^H$ on the entries of each Potts atom
$H_h$. The per-site rate multipliers $\{\eta_s\}$, the per-class
profiles $\{\pi^{(c)}\}$, the materialized Potts atoms $\{H_h\}$,
and the class-pair Potts assignments $\{h_{cc'}\}$ are no longer
parameters but \emph{latent} draws from the base measures; their
posteriors are inferred per site, per class, per atom, and per
class-pair respectively, with the corpus-level parameters governing
the priors. Indel parameters $(\lambda, \mu, r)$ decouple cleanly at
the parameter level via the alignment--substitution factorization
of \S\ref{sec:tkfdp-model} and are learned from the alignment
marginal alone by standard TKF92 EM; we focus on the substitution
side. The latent variables on each MSA are the key-DP labels
$\{z_s\}$, the site classes $\{c_s\}$, the per-site rates
$\{\eta_s\}$, the within-class amino-acid trajectories
$\{\sigma_C\}$, the unobserved internal-node states, the Potts-atom
assignments $\{h_{cc'}\}$, and (per \S\ref{sec:gravestone-scfg}) the
gravestone-augmented indel histories on every branch --- comprising
the alive and gravestone counts $(N_a, N_g)$ at every internal
node, plus the parentage tree and event times within each branch.

\subsubsection*{Variational EM with bridge-expectation sufficient
statistics}

For singleton key-DP classes, $Q^s$ reduces to F81 with
class-specific equilibrium, shared exchangeability, and per-site
rate. The bridge expectations~\citep{HolmesRubin2002} --- expected transition counts $N^{(b,s)}_{xy}$ and
dwell times $T^{(b,s)}_x$ along each branch $b$, conditioned on the
endpoint amino acids --- are computable in closed form by
integrating $\int_0^{t_b} e^{Q u} A\, e^{Q(t_b - u)}\,du$ in the
eigenspace of the rate matrix. Aggregated appropriately they give
closed-form posterior updates on every continuous latent on the
substitution side:
\begin{itemize}
    \item $\pi^{(c)}$ has a Dirichlet posterior under the
      $\mathrm{Dirichlet}(\kappa_\pi \bar\pi)$ prior, via the
      secret-destination augmentation of
      \S\ref{sec:base-measure-conj} below.
    \item $\eta_s$ has a Gamma posterior under the
      $\mathrm{Gamma}(a_\eta, b_\eta)$ prior; both shape and rate
      updates are immediate from the bridge expectations.
      The marginal site likelihood
      $\int \mathrm{Gamma}(a_\eta, b_\eta) \cdot \mathrm{Poisson}(N_s^{\mathrm{acc}} \mid \eta_s \tilde T_s)\, d\eta_s$
      is closed-form Negative-Binomial.
    \item Each materialized Potts atom $H_h$ has a Gaussian
      posterior under the $\mathcal{N}(\mu_{kl}, \tau_{kl}^{-1})$
      prior, via the Gaussian (Hessian-based) Laplace approximation
      of the path-DCA likelihood at the current iterate.
\end{itemize}
The shared exchangeability $S$ aggregates over all classes as in
standard GTR EM. Site-class assignments $c_s$ and Potts-atom
assignments $h_{cc'}$ are latent in the CRP sense, integrated over
via the Gibbs sweeps below.

For multi-site DP classes the joint generator on
$\mathcal{A}^{|C|}$ has no usable eigendecomposition, and
the bridge expectations cannot be applied to the coupled bridge directly. The
variational construction of \S\ref{sec:class-elbo} is the
substitute: at the variational fixed point each $s \in C$ has an
inhomogeneous single-site bridge with rate $\tilde Q^{s,*}$ on each
constant-composition segment of the branch
(\S\ref{sec:gravestone-scfg}), and the bridge expectations apply to this
bridge in its eigenspace. The integral identity goes through with
$Q$ replaced by $\tilde Q^{s,*}(u)$ (piecewise-constant on a time
grid within each segment), yielding per-site sufficient statistics
$\tilde N^{(b,s)}_{xy}$, $\tilde T^{(b,s)}_x$ as well as per-pair
co-occupancy times
\[
\tilde T^{(b, s, s')}_{xy} \;=\; \int_0^{t_b} q_s(x; u)\, q_{s'}(y; u)\, du,
\]
the latter immediate from the variational independence of single-site
bridges given endpoints.

The envelope theorem ensures that, evaluated at the variational
fixed point, the gradient of the ELBO with respect to $\theta$ has
no contribution from the implicit $\theta$-dependence of
$\tilde Q^{s,*}$:
\[
\nabla_\theta \mathcal{L}
\;=\; \mathbb{E}_\mathbb{Q}\!\bigl[\nabla_\theta \log \tfrac{d\mathbb{P}}{d\mathbb{Q}}\bigr]
\;=\; \mathbb{E}_\mathbb{Q}\!\bigl[\nabla_\theta \log d\mathbb{P}\bigr].
\]
For the GTR-Potts generator this gradient has three structurally
distinct pieces. The exchangeability gradient $\nabla_S \mathcal{L}$
is a linear combination of the single-site $\tilde N^{(b,s)}_{xy}$
counts as in standard GTR EM. The frequency gradient
$\nabla_{\pi^{(c)}} \mathcal{L}$ aggregates over sites with
$c_s = c$ and is again single-site. The coupling-tensor gradient is
a path-space DCA M-step,
\[
\nabla_{H_{cc'}(x, x')} \mathcal{L}
\;\propto\; \!\!\sum_{n, b}\!\!\sum_{\substack{s, s' \in C(n) \\ c_s = c,\, c_{s'} = c'}}\!\!\!\Bigl[\,\bigl(\text{expected } x_s = x,\, x_{s'} = x' \text{ co-occupancy under } \mathbb{P}\bigr) - \bigl(\text{same under } \mathbb{Q}\bigr)\,\Bigr],
\]
expressible in the same bridge-expectation format using the products
$q_s(x; u) q_{s'}(x'; u)$ for the variational dwell-time terms and
contributions from each variational jump intensity for the
transition terms.

\subsubsection*{MCMC over the discrete latents}

The combinatorial latents $\{z_s\}$, $\{c_s\}$, and $\{h_{cc'}\}$
admit straightforward Gibbs sampling. All three partitions are
integrated over their respective stick-breaking weights via Chinese
restaurant predictives. For the key DP,
\[
P(z_s = k \mid z_{-s}, \alpha_z)
\;\propto\;
\begin{cases}
|C_k^{(-s)}| \cdot \mathcal{L}_{\text{class}}(C_k \cup \{s\}) / \mathcal{L}_{\text{class}}(C_k)
& \text{(existing key class)}, \\
\alpha_z \cdot \mathcal{L}_{\text{class}}(\{s\})
& \text{(new singleton)},
\end{cases}
\]
where $\mathcal{L}_{\text{class}}$ is the per-branch-aggregated
class-level variational ELBO from \S\ref{sec:class-elbo} plus the
contribution of the rest of the tree via Felsenstein pruning.
Site-class updates use the class DP predictive,
\[
P(c_s = c \mid c_{-s}, \alpha_c, \theta)
\;\propto\;
\begin{cases}
n_c^{(-s)} \cdot \mathcal{L}_{\text{site}}(s; \pi^{(c)})
& \text{(existing class, profile $\pi^{(c)}$)}, \\
\alpha_c \cdot \int \mathcal{L}_{\text{site}}(s; \pi')\, G_0(d\pi')
& \text{(new class)},
\end{cases}
\]
with the new-class integral closed form by the
Dirichlet--multinomial conjugacy of
\S\ref{sec:base-measure-conj}. The Potts-atom assignment
$h_{cc'}$ for an unordered class-pair $\{c, c'\}$ is itself drawn
by CRP-Gibbs,
\[
P(h_{cc'} = h \mid \text{rest}, \alpha_H)
\;\propto\;
\begin{cases}
m_h^{(-cc')} \cdot \mathcal{L}_{\text{pair}}(\{c, c'\};\, H_h)
& \text{(existing atom)}, \\
\alpha_H \cdot \int \mathcal{L}_{\text{pair}}(\{c, c'\};\, H')\, G_0^H(dH')
& \text{(new atom)},
\end{cases}
\]
where $\mathcal{L}_{\text{pair}}$ is the path-DCA contribution from
co-occurrences of classes $c$ and $c'$ across the corpus, and the
new-atom integral is approximated by the Laplace scheme of
\S\ref{sec:base-measure-conj}.

Split--merge proposals of Jain--Neal type~\citep{jainneal2004} are
essential for mixing on all three partitions, with acceptance ratios
in terms of class-level (for $z$), site-aggregated (for $c$), or
class-pair-aggregated (for $h$) ELBO ratios. All three concentrations
$\alpha_z, \alpha_c, \alpha_H$ admit closed-form Gamma-prior
posteriors given the respective partitions via the auxiliary-variable
scheme of Escobar \& West~\citep{escobarwest1995}.

\subsubsection*{Truncated stick-breaking on Potts atoms}

The Potts-atom CRP admits a truncated stick-breaking (TSB)
alternative that is cheaper per sweep, requires no new-atom Laplace
branch inside the Gibbs step, and removes the empty-atom bookkeeping
of the CRP. The construction mirrors Ishwaran \&
James~\citep{ishwaranjames2001} and Blei \&
Jordan~\citep{bleijordan2006} as applied to the site-class DP,
transposed to the Potts-atom DP.

Fix a truncation level $K_H^{\max}$. The natural choice is
$K_H^{\max} = K_c(K_c + 1)/2$, the total number of unordered
class-pairs --- with this bound the truncation is exactly tight
(every class-pair could in principle have its own atom). Replace
the CRP-Gibbs sweep over $\{h_{cc'}\}$ with the following two-step
alternation:

\begin{enumerate}
\item \textbf{Beta CAVI / Gibbs on stick proportions.} Maintain
  stick variables $\beta_h \sim \mathrm{Beta}(1, \alpha_H)$ for
  $h < K_H^{\max}$, with $\beta_{K_H^{\max}} = 1$ and
  $\rho_h = \beta_h \prod_{h' < h} (1 - \beta_{h'})$. Given the
  current per-atom counts $m_h = \#\{(c, c') : h_{cc'} = h\}$, the
  conditional posterior on $\beta_h$ is the closed-form Beta
\[
    \beta_h \mid m, \alpha_H \;\sim\;
    \mathrm{Beta}\!\Bigl(1 + m_h,\; \alpha_H + \!\!\sum_{h' > h}\!\! m_{h'}\Bigr).
\]
Use either the posterior mean (CAVI) or a Beta sample (Gibbs).

\item \textbf{Categorical resample on $\{h_{cc'}\}$.} For each
  unordered class-pair $\{c, c'\}$, draw the atom assignment from
  the closed-form Categorical
\[
    P(h_{cc'} = h \mid \text{rest}) \;\propto\;
    \rho_h \cdot \mathcal{L}_{\text{pair}}(\{c, c'\};\, H_h),
\]
where $\mathcal{L}_{\text{pair}}$ is the path-DCA contribution from
co-occurrences of classes $c$ and $c'$ across the corpus, evaluated
at the current materialized atom $H_h$.
\end{enumerate}

\begin{remark}[Bijective initialization]
Initialization at $K_H^{\max} = K_c(K_c + 1)/2$ should set each
unordered class-pair to a \emph{distinct} atom drawn from $G_0^H$,
not collapse all class-pairs onto a single shared atom. The latter
is a chicken-and-egg failure mode: only the populated atom receives
data through the Laplace M-step, so the others stay at the prior;
the next TSB resample therefore prefers the populated atom for all
pairs (its likelihood dominates the others' prior-only score), and
the system collapses. The bijective init guarantees every atom
carries at least one class-pair's data through the first Laplace
M-step.
\end{remark}

\subsubsection*{Locality of DP updates under the SCFG augmentation}

When the augmented indel history is resampled
(\S\ref{sec:gravestone-scfg}), new gravestone fragments appear with
new residues that need DP keys, and old gravestones disappear.
Naively this would suggest a full DP-key resample after every
augmentation step, which would be prohibitive. Exchangeability of
the CRP and the segment-wise structure of the variational ELBO
together remove almost all of this cost.

The decomposition is two-step within a single sweep. The
\emph{indel-skeleton step} resamples the gravestone-augmented history
on every branch by the SCFG traceback of \S\ref{sec:gravestone-scfg}
without touching DP keys; cost $O((b + d) \log(b + d))$ per branch,
independent of DP state. The \emph{gravestone-key step} then
assigns DP keys to each new residue $s$ on a gravestone fragment
with lifetime $[u, v]$ on branch $b$, using the local CRP
predictive over existing classes plus a singleton, evaluated against
$\mathcal{L}_{\text{seg}}$ on segments of $b$ within the gravestone's
lifetime. Adding $s$ to class $C_k$ changes that class's variational
fixed point only on those segments, costing $O(|C_k| \cdot A^3)$
per fixed-point iteration per touched segment. Per gravestone, the
work is one Categorical draw with support equal to the existing
classes plus a singleton, plus one local fixed-point re-solve on
the affected segments.

\subsubsection*{Closed-form base-measure integrals via complete-data augmentation}
\label{sec:base-measure-conj}

The base-measure integrals
\[
    \int \mathcal{L}_{\text{site}}(s; \pi')\, G_0(d\pi'), \qquad
    \int \mathcal{L}_{\text{pair}}(\{c, c'\};\, H')\, G_0^H(dH'),
\]
that appear in the new-class and new-atom branches of the CRP
predictives, and the analogous corpus-level marginals over
$\eta_s$, all admit closed-form approximations under complete-data
augmentation.

\paragraph{Secret-destination augmentation for $\pi^{(c)}$.} The
F81-form rate factors as
$Q^s_{xy} = \eta_s \cdot S_{xy} \cdot \pi^{(c_s)}(y) \cdot \exp(-\tfrac{1}{2}\Delta H_s)$,
which decomposes the dynamics into an $S$-driven Poisson proposal
clock at rate $\eta_s S_{xy}$ from state $x$ to destination $y$,
modulated by a Bernoulli filter that accepts the proposal with
probability $\pi^{(c_s)}(y) \exp(-\tfrac{1}{2}\Delta H_s)$. To
recover full multinomial structure on $\pi^{(c)}$ we augment each
silent failure of a proposal of $y$ with a categorical \emph{secret
destination} $j \neq y$, drawn from $\pi^{(c)}(j)/(1 - \pi^{(c)}(y))$
on the simplex. Each proposal then casts exactly one categorical
vote: the destination is $y$ with probability $\pi^{(c)}(y)$ on
acceptance, and a secret $j$ with probability $\pi^{(c)}(j)$ on
rejection. Marginally over the proposal, the vote is
$\mathrm{Categorical}(\pi^{(c)})$, so the augmented count vector
$N^{(c)} \in \mathbb{Z}_{\geq 0}^A$ satisfies
\[
    N^{(c)} \mid \pi^{(c)},\, \text{augmentation} \;\sim\;
    \mathrm{Multinomial}\bigl(N^{(c)}_{\text{total}},\, \pi^{(c)}\bigr),
\]
and the posterior is the strict closed-form Dirichlet
\[
    \pi^{(c)} \mid N^{(c)} \;\sim\;
    \mathrm{Dirichlet}\bigl(\kappa_\pi \bar\pi + N^{(c)}\bigr).
\]
In the EM/VBEM setting the augmentation is replaced by its
conditional expectation: each proposal contributes $\pi^{(c)}(y)$ to
$E[N^{(c)}_y]$ regardless of which destination was proposed. The
Dirichlet update is then a fixed-point iteration whose limit is the
conjugate posterior under the conditional augmentation. The
new-class integral $\int \mathcal{L}_{\text{site}}(s; \pi')\, G_0(d\pi')$
is the analytic ratio of multivariate Beta functions
$B(\kappa_\pi \bar\pi + N^{(c)}) / B(\kappa_\pi \bar\pi)$.

\paragraph{Gamma--Poisson conjugacy for $\eta_s$.} The F81
substitution-count likelihood at site $s$ is itself Gamma-conjugate
in $\eta_s$ without any need for the secret-destination
augmentation. The expected substitution count along the tree is
Poisson with rate $\eta_s \cdot \tilde T_s$, where
\[
    \tilde T_s \;=\; \sum_b \sum_x T^{(b,s)}_x \cdot
    \sum_{y \neq x} S_{xy} \, \pi^{(c_s)}(y)
\]
is the $\pi$-weighted dwell-time integral. Under the
$\mathrm{Gamma}(a_\eta, b_\eta)$ prior the conjugate posterior is
closed form
\[
    \eta_s \mid \mathrm{path},\, \pi^{(c_s)} \;\sim\;
    \mathrm{Gamma}\!\bigl(a_\eta + N_s^{\mathrm{acc}},\; b_\eta + \tilde T_s\bigr),
\]
with $N_s^{\mathrm{acc}} = \sum_b \sum_{x \neq y} N^{(b,s)}_{xy}$
the total observed substitution count, and the per-site marginal
likelihood is closed-form Negative-Binomial.

\paragraph{Laplace for $H_h$.} Each Potts atom is real-valued and
the path-DCA likelihood is locally Gaussian in $H_h$ to second
order. We approximate
\[
    \log p_{\mathrm{path-DCA}}(\mathrm{data} \mid H_h)
    \;\approx\; \log p_{\mathrm{path-DCA}}(\mathrm{data} \mid \hat H_h)
    \;-\; \tfrac{1}{2} (H_h - \hat H_h)^\top \Lambda_h (H_h - \hat H_h),
\]
where $\hat H_h$ is the MAP under the current iterate, found by a
few Newton or natural-gradient steps starting from a chosen seed,
and $\Lambda_h = -\nabla^2_{H_h} \log p_{\mathrm{path-DCA}}$ is the
Hessian at $\hat H_h$. Combined with the
$\mathcal{N}(\mu_{kl}, \tau_{kl}^{-1})$ prior on each entry, the
posterior is the closed-form Gaussian
\[
    H_h \mid \mathrm{data} \;\sim\;
    \mathcal{N}\!\bigl(\hat H_h^{\mathrm{post}},\, (\mathrm{diag}(\tau_{kl}) + \Lambda_h)^{-1}\bigr).
\]
The Laplace estimate of the new-atom integral is the standard
$\log p_{\mathrm{path-DCA}}(\hat H_h) + \log G_0^H(\hat H_h) + \tfrac{d}{2}\log(2\pi) - \tfrac{1}{2}\log\det(\mathrm{diag}(\tau_{kl}) + \Lambda_h)$
with $d = A(A+1)/2$ the symmetric-slice dimension.

\paragraph{Multi-seed Laplace mixture for multimodal posteriors.}
The Laplace approximation collapses around a single mode and
underestimates posterior mass elsewhere; for atoms that admit
multiple plausible coupling patterns (e.g., charge-complementary
vs.\ hydrophobic-hydrophobic minima) this is a genuine approximation
error. We mitigate by running the few-step optimizer from a small
number $K_{\mathrm{seed}}$ of distinct seeds, obtaining
$K_{\mathrm{seed}}$ Laplace components, and combining them as a
weighted Gaussian mixture.

\subsubsection*{The recommended hybrid}

The natural hybrid is:
\begin{enumerate}
    \item \textbf{Inner variational E-step.} For each MSA, branch
      and DP class, segment the branch by the current sampled
      augmented indel history and solve the geometric-mean fixed
      point of \S\ref{sec:class-elbo} on each segment to obtain
      $\{\tilde Q^{s,*}\}$ and the per-class ELBO. No parameter or
      structure updates here.
    \item \textbf{MCMC over discrete and augmentation latents.}
      Resample augmented indel histories on every branch by recursive
      traceback on the time-indexed pair SCFG of
      \S\ref{sec:gravestone-scfg}, with closed-form midpoint
      marginals from the bivariate Riccati and Felsenstein-style
      upwards--downwards composition for internal-node augmentation
      counts. Gibbs sweep over $\{z_s\}$, $\{c_s\}$, and
      $\{h_{cc'}\}$ using the class-level ELBOs from step~1 as
      evidence; periodic Jain--Neal split--merge proposals on each
      partition.
    \item \textbf{Stochastic M-step.} Compute bridge-expectation
      sufficient statistics $\tilde N$, $\tilde T$,
      $\tilde T_{\mathrm{pair}}$ from the variational single-site
      bridges, summed over augmented constant-composition segments
      within each branch, over a minibatch of MSAs. Take a stochastic
      gradient step on the corpus-level hyperparameters
      $\theta = (S, \alpha_z, \alpha_c, \alpha_H, a_\eta, b_\eta, \kappa_\pi, \bar\pi, \{\mu_{kl}, \tau_{kl}\})$.
      Per-class $\pi^{(c)}$ updates by its closed-form Dirichlet
      posterior under the secret-destination augmentation; per-site
      $\eta_s$ marginalizes in closed form via Gamma--Poisson
      conjugacy; per-atom $H_h$ updates by Laplace with multi-seed
      mixture against the prior $G_0^H$. The three DP concentrations
      update by the Escobar--West auxiliary-variable scheme.
    \item \textbf{Iterate.} The inner variational fixed point in
      step~1 is re-solved cheaply after each step because the
      previous solution is a warm start.
\end{enumerate}

This factors the overall problem into the three regimes where each
tool is strongest: variational for the continuous-time substitution
dynamics where it admits closed-form rates and bridge-expectation
sufficient statistics; MCMC for the discrete combinatorial latents
where it mixes well and avoids mean-field pathologies; stochastic
gradient for the continuous parameters $\theta$ where it scales to
corpus-level training. In SVI form the M-step is over a minibatch
of MSAs at each iteration and amortizes naturally over the corpus.

\subsubsection*{Counts-tensor representation of the training corpus}

Per family we precompute a counts tensor
$(\text{cherry} \times \text{column} \times \text{AA-pair} \times \text{branch length})$,
packed as int8, that decouples inference from raw MSA parsing and
lets one outer SVI iteration scale to $10^3$--$10^4$ \Pfam\ families.
Branch lengths are quantized to $n_t \approx 50$ geometrically spaced
bins between $\tau_{\min}$ and $\tau_{\max}$; geometric spacing
exploits that $\exp(Q\tau)$ varies fastest at small $\tau$, giving a
$\sim 5\times$ cache reduction over linear quantization at no
measurable accuracy cost. Per-class single-site log-transition
matrices $(K_c, n_t, A^2)$ are precomputed once per outer iteration
and shared across all families.

\subsection{Pairwise alignment postprocessing}
\label{sec:postprocessing}

Multiple-sequence alignment tools of the FSA
family~\citep{bradley2009fsa} (and other consistency-based schemes)
operate on pairwise residue posterior matrices
$Q^{(X, Y)}_{ij} = P(X_i \sim Y_j \mid X, Y)$ produced by per-pair
forward--backward passes through a Pair HMM. The TKF-DP postprocessing
target is the joint-marginal alignment-and-partition likelihood
described next, in which an alignment $A$ together with a
size-$\{1, 2\}$ partition $E$ of its Match cells contributes a
product of per-block substitution factors with no class or
rate-multiplier latent left explicit.

\begin{remark}[Historical note: deprecated first-order correction]
\label{rem:deprecated-mean-field}
An early formulation of TKF-DP postprocessing expressed the
coupling as a first-order mean-field correction
$\log Q'_{ij} - \log Q_{ij} \approx \varepsilon \sum_{(i', j')}
  Q_{i' j'} [M(i, j; i', j'; t) - 1]$
to the baseline single-site posterior, with $\varepsilon$ a
per-pair partner prior under the size-$\{1, 2\}$ Ewens
distribution and $M$ a column-class-marginalized Potts boost. That
correction is a first-order expansion of the joint marginal in
the partner-prior strength and is now considered a heuristic
stepping-stone; the modern infinite Pair HMM
(\S\ref{sec:infinite-hmm} of Appendix~\ref{app:phmm}) treats the
same coupling exactly via the joint-marginal target below, with
class assignments and rate multipliers integrated analytically
rather than re-marginalized at every column.
\end{remark}

\subsubsection*{The unified joint-marginal target}

Let $\theta_{\mathrm{indel}} = (\lambda, \mu, r)$ be the TKF92
indel parameters and $\alpha_z$ the Ewens partition concentration.
Write $\mathrm{Match}(A)$ for the Match cells of an alignment $A$,
and let $E$ be a set of unordered Match-cell edges (a partition of
$\mathrm{Match}(A)$ into singletons of size~1 and doublets of
size~2; we let $V(E) \subset \mathrm{Match}(A)$ denote the doublet
endpoints). The TKF-DP joint marginal at branch length $t$ is
\begin{equation}
\label{eq:tkfdp-joint-marginal}
P(X, Y; t) \;=\; \sum_A \pi_{\TKF92}(A \mid t, \theta_{\mathrm{indel}})
  \sum_E \pi_{\mathrm{Ewens}}(E \mid \alpha_z, |\mathrm{Match}(A)|)
  \!\!\!\prod_{(i, j) \in \mathrm{Match}(A) \setminus V(E)}\!\!\!\!\!\!
    P_{\mathrm{singlet}}(x_i, y_j; t)
  \prod_{e = ((i, j), (k, l)) \in E}\!\!\!\!\!
    P_{\mathrm{doublet}}\bigl((x_i, x_k), (y_j, y_l); t\bigr),
\end{equation}
with the per-block primitives
\begin{align}
P_{\mathrm{singlet}}(a, b; t)
  \;&=\; \sum_c \pi_c \,\, \pi^{(c)}(a) \,\, P_c(a \to b;\, t \cdot \eta),
  \label{eq:p-singlet} \\
P_{\mathrm{doublet}}\bigl((a, c), (b, d); t\bigr)
  \;&=\; \sum_{c_1, c_2} \pi_c(c_1) \, \pi_c(c_2)
  \,\, \pi_{\mathrm{joint}}^{c_1, c_2}(a, c)
  \,\, P_{\mathrm{joint}}^{c_1, c_2}\!\bigl((a, c) \to (b, d);\, t, \eta_1, \eta_2\bigr).
  \label{eq:p-doublet}
\end{align}
The doublet endpoint convention (matching the boost API of the
released code) is that $(a, b)$ are the amino acids at the left
endpoint $(i, j)$ and $(c, d)$ at the right endpoint $(k, l)$:
$a = x_i$, $b = y_j$, $c = x_k$, $d = y_l$. Indices $c_1, c_2$ run
over site classes and should not be confused with the
endpoint-amino-acid index $c$.

The empirical class prior $\pi_c$ is the per-class column count
from training (\emph{not} uniform $1 / K_c$); the per-class profile
$\pi^{(c)}(\cdot)$ is the trained \texttt{pi\_class}; the
class-conditional transition is
\begin{equation}
\label{eq:p-class-cond}
P_c(a \to b; t) \;=\; [\exp(Q_c \cdot t)]_{a, b},
\qquad
Q_c \;=\; (S_{\mathrm{LG08}} - \mathrm{diag}\, S_{\mathrm{LG08}})
          \cdot \pi^{(c)}\![\,\cdot\,]^\top,
\end{equation}
the standard F81-form GTR generator with exchangeability
$S_{\mathrm{LG08}}$~\citep{LG08} and class-specific stationary
$\pi^{(c)}$. The joint stationary
$\pi_{\mathrm{joint}}^{c_1, c_2}(a, c)$ at a coupled cell-pair is
the joint stationary of a Potts atom $H = \mathtt{atoms}[\mathtt{assignments}[c_1, c_2]]$
at the trained checkpoint (the trained
\texttt{potts\_dp.atoms} indexed by \texttt{potts\_dp.assignments}),
and $P_{\mathrm{joint}}^{c_1, c_2}$ is its time-evolved transition
$\exp(Q_{\mathrm{joint}}^{c_1, c_2}(H) \cdot t)$ built by
\texttt{build\_joint\_Q\_pair}$(H, \pi_a, \pi_b, S, \eta_{\mathrm{pair}} = (\eta_1, \eta_2))$
as described in \S\ref{sec:gtr-potts}.

\paragraph{Pair-background convention for the joint generator.}
Two natural choices for the pair background $(\pi_a, \pi_b)$
arise. The first is to use the LG08 single-site equilibrium
$\pi_a = \pi_b = \pi^{\mathrm{LG08}}$ for every class-pair
$(c_1, c_2)$. This is the \emph{canonical} convention for the
released $K{=}4$ EM-warmup checkpoint: the trained Potts atoms
were fit as deviations of the joint stationary from the
LG08-pair-background reference, so inference under that checkpoint
must use the same convention for the joint stationary and joint
generator to match training. The second choice,
$\pi_a = \pi^{(c_1)}$, $\pi_b = \pi^{(c_2)}$, uses the per-class
profiles as the pair background; this is the structurally
``cleaner'' choice for future model variants in which the Potts
atoms are retrained against the per-class background, but it is
\emph{not} consistent with the released checkpoint.

\paragraph{Discrete-gamma rate-multiplier integration.}
The per-site rate multiplier $\eta_s \sim \mathrm{Gamma}(a_\eta, b_\eta)$
of \S\ref{sec:gtr-potts} is integrated analytically by the discrete-gamma
quadrature of Yang~\citep{Yang1994a}: $K_r$ equiprobable bin
medians $\eta_1, \ldots, \eta_{K_r}$ are taken from the
$\mathrm{Gamma}(a_\eta, b_\eta)$ prior with uniform weight
$1 / K_r$, and the per-block primitive of
equation~\ref{eq:p-singlet} (resp.~\ref{eq:p-doublet}) is averaged
over rate-multiplier draws via
$P_{\mathrm{singlet}}(a, b; t) = K_r^{-1} \sum_{r = 1}^{K_r}
  P_{\mathrm{singlet}}(a, b; t \cdot \eta_r)$
(resp.\ a double sum over $(\eta_1, \eta_2)$ for the doublet).
The Gamma hyperparameters $(a_\eta, b_\eta)$ at inference time
must match those used at training time; the released
$K{=}4$ EM-warmup checkpoint uses
$(a_\eta, b_\eta) = (2, 2)$. The reduction $K_r = 1$ with
$\eta = 1$ is the canonical evaluation of that checkpoint, since
the stored per-site $\eta$ posteriors all collapsed to $1.0$
during EM warmup; for retrained checkpoints with a non-trivial
$\eta$ posterior, $K_r > 1$ is required.

\paragraph{$t = 0$ reduction identities.}
At $t = 0$ both per-block primitives reduce to the appropriate
marginal stationary:
\begin{align}
\label{eq:p-singlet-t0}
P_{\mathrm{singlet}}(a, a;\, 0)
  \;&=\; \sum_c \pi_c \, \pi^{(c)}(a) \;\equiv\; \pi_{\mathrm{marg}}(a), \\
\label{eq:p-doublet-t0}
P_{\mathrm{doublet}}\bigl((a, c), (a, c);\, 0\bigr)
  \;&=\; \sum_{c_1, c_2} \pi_c(c_1) \, \pi_c(c_2) \,
        \pi_{\mathrm{joint}}^{c_1, c_2}(a, c)
        \;\equiv\; \pi_{\mathrm{marg, pair}}(a, c).
\end{align}
These follow from $P_c(a \to b;\, 0) = \delta_{ab}$ and from the
joint stationary $\pi_{\mathrm{joint}}^{c_1, c_2}$ being a
fixed point of $P_{\mathrm{joint}}^{c_1, c_2}(\cdot;\, 0)$.

\paragraph{The boost tensor.}
The four-residue Potts \emph{boost} entering equation~\ref{eq:tkfdp-joint-marginal}
through the doublet factor is the multiplicative deviation of the
doublet emission from the product of two singlet emissions:
\begin{equation}
\label{eq:M-boost}
M\bigl((a, c), (b, d); t\bigr)
  \;\equiv\;
  \frac{P_{\mathrm{doublet}}\bigl((a, c), (b, d); t\bigr)}
       {P_{\mathrm{singlet}}(a, b; t) \cdot P_{\mathrm{singlet}}(c, d; t)}.
\end{equation}
Indexing convention as above: $(a, b)$ is the left endpoint, $(c, d)$
the right endpoint. By construction $M \equiv 1$ when every Potts
atom is zero (so $\pi_{\mathrm{joint}}^{c_1, c_2}$ factorizes into
$\pi^{(c_1)} \otimes \pi^{(c_2)}$ and $P_{\mathrm{joint}}^{c_1, c_2}$
factorizes into the two single-site $P_{c_i}$ factors). The boost
tensor $M$ is the object that the bounded-edge augmented Pair HMM
methods of Appendix~\ref{app:phmm} fold multiplicatively into the
Match emissions to score doublets relative to singletons. The MCMC
sampler operates on the joint $(A, E)$ directly via
equation~\ref{eq:tkfdp-joint-marginal}, with class assignments and
rate multipliers integrated analytically as above (\emph{not}
sampled).

\subsubsection*{Sequence annealing with coevolutionary scoring}
\label{sec:coupled-annealing}

The joint-marginal of equation~\ref{eq:tkfdp-joint-marginal}
implies a per-Match-cell posterior $Q'_{ij}$ that the augmented
Pair HMM machinery of Appendix~\ref{app:phmm} computes either
exactly under bounded-edge truncation or stochastically under the
infinite-Pair-HMM MCMC. An alternative routing of the same
coupling signal applies the boost \emph{during} assembly: the
greedy column-merging step inside the sequence annealer is extended
to score pairs of merges jointly. With pair-HMM input, the
canonical edge score is
\begin{equation}
\label{eq:tgf-edge}
w(e) \;=\; \frac{2 \, Q_{ij}^{(X, Y)}}{g_i^{(X, Y)} + g_j^{(Y, X)}},
\qquad e = \bigl(\text{column of } X_i,\ \text{column of } Y_j\bigr),
\end{equation}
where $g$ are the per-residue gap posteriors. The coupled-pair
extension admits joint candidates consisting of pairs of merges
scored as
\begin{equation}
\label{eq:coupled-score}
W(e_1, e_2) \;=\; \log Q_{ij}^{(X, Y)} + \log Q_{i'j'}^{(X, Y)}
  + \log M\bigl((X_i, X_{i'}), (Y_j, Y_{j'}); t\bigr),
\end{equation}
with $M$ the boost tensor of equation~\ref{eq:M-boost} evaluated
at the observed amino acids. The first two terms are the baseline
log-posteriors, and the third is the four-residue Potts
contribution; the latter is identically zero when every Potts
atom is zero.

\paragraph{Greedy schedule and conflict resolution.} When a coupled
candidate $(e_1, e_2)$ is popped from the queue, conflicts are
resolved as follows. If both component merges are still admissible
(no same-sequence conflict, no DAG cycle), both are committed in a
single step and the boost-driven joint commit is recorded. If $e_1$
is admissible but $e_2$ has been invalidated by an intervening
merge, the candidate degrades to a single-edge merge with score
$\log Q_{ij}$ and is re-inserted into the queue at that priority.
If neither component is admissible, the candidate is dropped.

\paragraph{Cost analysis.} Under two practical prunings (threshold
$q_{\min} = 0.1$ on the single-edge posterior and threshold
$\mu_{\min} = 0.1$ nat on $|\log M|$), the per-step cost is
$O(L^2)$ amortized and the total cost is $O(L^3)$ in the worst
case, comparable to the $O(L^2 \log L)$ baseline assembler.

\paragraph{Augmented-PHMM versus greedy-pairwise tradeoff.}
The augmented-Pair-HMM and infinite-Pair-HMM methods of
Appendix~\ref{app:phmm} produce a single corrected posterior
matrix $Q'$ that any downstream assembler can consume. The
coupled-pair extension above scores strongly-coupled pairs jointly
on their native four-residue domain inside the assembler, but
commits greedily and is therefore vulnerable to local optima where
an early commit on a coupled pair locks out a globally better
arrangement. The two routes consume the same boost tensor $M$ of
equation~\ref{eq:M-boost} but apply it at different points in the
pipeline; the augmented-PHMM route is the principled default for
TKF-DP postprocessing in \S\ref{app:phmm}.
